\documentclass[a4paper,11pt]{ltjsarticle}
\usepackage[haranoaji]{luatexja-preset}
\usepackage[top=23mm,bottom=23mm,left=25mm,right=25mm]{geometry}
\usepackage{lmodern}
\usepackage{mathtools,amssymb,amsthm}
\usepackage{booktabs}
\usepackage{xcolor}
\usepackage[unicode,hidelinks]{hyperref}

\newcommand{\upstream}{https://github.com/lsernau/grh-lean/blob/72fead102f663d5c7a407964a13dc0cf9d977207}
\newcommand{\leansrc}[2]{\href{\upstream/grh/src/#1.lean}{#2}}
\definecolor{supplementblue}{HTML}{174A7A}
\newenvironment{editorialnote}
  {\begin{quote}\color{supplementblue}\small\textbf{補足}\quad\ignorespaces}
  {\end{quote}}
\newcommand{\RePart}{\operatorname{Re}}
\newcommand{\ImPart}{\operatorname{Im}}
\newcommand{\res}{\operatorname*{Res}}

\theoremstyle{plain}
\newtheorem{theorem}{定理}[section]
\newtheorem{proposition}[theorem]{命題}
\newtheorem{lemma}[theorem]{補題}
\newtheorem{corollary}[theorem]{系}
\theoremstyle{definition}
\newtheorem{definition}[theorem]{定義}
\renewcommand{\proofname}{証明}
\renewcommand{\refname}{参考文献}
\numberwithin{equation}{section}
\allowdisplaybreaks[1]
\setlength{\emergencystretch}{1.5em}

\hypersetup{
  pdftitle={ディリクレ L 関数の一般化リーマン予想},
  pdfsubject={Alpöge の Lean 形式証明に基づく日本語再構成稿}
}
\title{ディリクレ \(L\) 関数の一般化リーマン予想\\[0.45em]
  \large テータ持ち上げを用いる形式証明の再構成}
\author{}
\date{\small 再構成稿 \quad 2026年11月26日}

\begin{document}
\maketitle

\begin{abstract}
Levent Alpöge と Claude による一般化リーマン予想の形式証明を、
\emph{grh-lean} の可読版に沿って記述する。
証明は、二次指標で重みづけた算術級数内の素数冪和の剰余に対する片側評価を与える。
そのため、ユニタリ群の間のテータ持ち上げを反復し、持ち上げた形式の二乗余体積と、
有限素点で整性を得るための指数を比較する。
上界にはテータ微分の整除性とポアンカレ級数の二次モーメントを、
下界には素点ごとに選んだ捻り、グラム行列式の比、および積公式を用いる。
片側評価から半平面での非消滅を得る際には、二次捻りの零点密度評価により
対数微分の差で留数が相殺される可能性を除き、Landau の補題を適用する。
\end{abstract}

\section{主定理}

\begin{theorem}[Alpöge の形式証明]\label{thm:grh}
\(q\ge 1\) を整数とし、\(\chi\) を法 \(q\) のディリクレ指標とする。
\begin{enumerate}
\item \(\RePart s>1/2\) ならば \(L(s,\chi)\ne0\) である。ただし \(\chi\) が主指標で
      \(s=1\) の場合を除く。
\item \(0<\RePart s<1\) において \(L(s,\chi)=0\) ならば、\(\RePart s=1/2\) である。
\item リーマンゼータ関数の零点は、自明零点 \(-2,-4,\ldots\) を除いて
      \(\RePart s=1/2\) 上にある。
\end{enumerate}
\end{theorem}

\noindent
ここで \(L(s,\chi)=\sum_{n\ge1}\chi(n)n^{-s}\) は
\(\RePart s>1\) からの有理型接続であり、\(\chi\) の原始性は仮定しない。
法 \(1\) の指標の \(L\) 関数は \(\zeta(s)\) である。
底本の五つの定理文のうち、ゼータ関数に関する二つは (3) を、
ディリクレ \(L\) 関数に関する三つは (1)、(2) および (2) の原始指標の場合を述べる。
\leansrc{Statements}{定理文}と
\leansrc{Main}{主定理の宣言}を参照されたい。

まず、素数冪和の片側上界から (1) を導く
（第\ref{sec:analytic}節）。次に、補助的な数体の無限素点に重さを割り当てることで、
この剰余をテータ微分の整除性を測る量として表す（第\ref{sec:data}節）。
第\ref{sec:lattices}節では、同じグラム行列式に対して
無限素点と有限素点の情報を合わせる方法を述べる。
第\ref{sec:upper-bound}節から第\ref{sec:determinant-estimates}節までで
比較に必要な上界、次元評価、捻りによる行列式比の評価をそれぞれ与え、
第\ref{sec:conclusion}節で定理を導く。

本稿は Alpöge の公表文 \cite{alpoege} に記された証明のうち、
Sernau らが整理した Lean 開発の可読版 \cite{grh-source} に拠る。
本文の L で始まる参照番号から、巻末にある底本の宣言を参照できる。
とくに局所テータ核の構成や補助捻りを得るための合同の評価には、底本の定理を引用する。
青字の「補足」は、本稿で読解のために付した注記である。
証明の着想と関連資料は、直後の概説にまとめた。

\section{素数冪和と零点}\label{sec:analytic}

\subsection{剰余のラプラス変換}

\(d\) を奇素数、\(q\) を \(d\) と互いに素な正整数とする。
\(\chi_d(n)=(n/d)\) をルジャンドル記号、\(\Lambda\) を von Mangoldt 関数、
\(\varphi\) を Euler のトーシェント関数とし、
\[
  \lambda_q(n)=\mathbf 1_{n\equiv1\bmod q}\Lambda(n)(1-\chi_d(n))
\]
とおく。当面は \(d\) を固定し、記号 \(\lambda_q,R_q,G_q\) 等の添字には \(d\) を書かない。
\(\psi\) が主指標のとき \(L^\circ(s,\psi)=(s-1)L(s,\psi)\) とし、
\(s=1\) では正則に延長した値を取る。\(\psi\) が非主指標のときは
\(L^\circ(s,\psi)=L(s,\psi)\) とする。いずれも整関数である。
\(\psi\chi_d\) は法 \(qd\) の指標とみなす。
\(d\) が奇素数で \((d,q)=1\) なので、\(\psi\chi_d\) は非主指標である。
\begin{align}
  G_q(s)
   &=\frac1{\varphi(q)}
     \left[
       \sum_{\psi\bmod q}
       \left(
        \frac{(L^\circ)'(s,\psi\chi_d)}{L^\circ(s,\psi\chi_d)}
        -\frac{(L^\circ)'(s,\psi)}{L^\circ(s,\psi)}
       \right)-1
     \right], \label{eq:log-derivative}\\
  R_q(y)
   &= \sum_{n\le y}\lambda_q(n)\left(\frac1n-\frac1y\right)
      -\frac{\log y}{\varphi(q)}-c_q,\qquad c_q=\RePart G_q(1).
  \label{eq:prime-remainder}
\end{align}
対数微分は変数 \(s\) に関するものとする。主指標について極を除いた
\(L^\circ\) は \(s=1\) で非零であり、非主指標についても \(L(1,\psi)\ne0\)
であるから、\(G_q\) は \(1\) で正則である。

指標の直交関係により、\(\lambda_q(n)\) は
\(\Lambda(n)(\psi(n)-(\psi\chi_d)(n))\) の法 \(q\) の全指標にわたる平均である。
\(\RePart s>1\) で Euler 積の対数微分を取り、主指標を整関数化したとき
加わる \(1/(s-1)\) を取り出すと
\begin{equation}\label{eq:lambda-dirichlet}
  F_q(s):=\sum_{n\ge1}\frac{\lambda_q(n)}{n^s}
        =\frac{s}{\varphi(q)(s-1)}+G_q(s).
\end{equation}
実軸上 \(s>1\) で右辺は実数なので、\(G_q(1)=c_q\) でもある。

\begin{lemma}\label{lem:prime-transform}
\(\RePart s>2\) において
\begin{equation}\label{eq:prime-transform}
  \mathcal T_q(s):=\int_0^\infty R_q(e^v)e^{(1-s)v}\,dv
     =\frac{G_q(s)/s-G_q(1)}{s-1}.
\end{equation}
右辺による有理型接続は \(s=1\) で正則である。\(\RePart \rho>0\)、\(\rho\ne1\)
での極は高々単純であり、その留数は
\[
  \res_{s=\rho}\mathcal T_q(s)=
    \frac{\nu(\rho)-\nu_d(\rho)}{\varphi(q)\rho(1-\rho)}.
\]
ここで \(\nu,\nu_d\) はそれぞれ \(L^\circ(s,\psi)\)、\(L^\circ(s,\psi\chi_d)\)
の \(\rho\) における零点の位数を法 \(q\) の全指標について加えたものである。
\end{lemma}

\begin{proof}
\(A_r(y)=y^{-r}\sum_{n\le y}\lambda_q(n)n^{r-1}\) とおく。
\(r=0,1\) について各項を \(v\ge\log n\) で積分すれば
\[
 \int_0^\infty A_r(e^v)e^{(1-s)v}\,dv=\frac{F_q(s)}{s+r-1}.
\]
\(\RePart s>2\) では項別積分が絶対収束する。したがって
\[
 \mathcal T_q(s)
 = F_q(s)\left(\frac1{s-1}-\frac1s\right)
    -\frac1{\varphi(q)(s-1)^2}-\frac{G_q(1)}{s-1}.
\]
\eqref{eq:lambda-dirichlet} を代入すると \eqref{eq:prime-transform} を得る。
\(s=1\) では分子が消える。\(s=\rho\ne1\) では、整関数の対数微分の留数が
零点の位数に等しいことを \eqref{eq:log-derivative} に適用する。
対応する宣言は \leansrc{New/Reduction/Zeros/PrimeError}
{\texttt{integral\_primeRemainder}, \texttt{isPolar\_primeRemainder}} である。
\end{proof}

\subsection{二次捻りの選択}

\begin{lemma}\label{lem:select-quadratic}
\(q\) は \(8\) の正の倍数とし、\(d_0\equiv7\pmod8\)、\((d_0,q)=1\) とする。
任意の \(\beta>1/2\)、\(H\ge0\) に対して、
剰余類 \(d\equiv d_0\pmod q\) に属する素数 \(d\) で、
すべての法 \(q\) の指標 \(\psi\) について \(L^\circ(s,\psi\chi_d)\) が
\(\RePart s\ge\beta,\ |\ImPart s|\le H\) で非零となるものが存在する。
\end{lemma}

\begin{proof}
\(\beta\ge1\) なら \(\beta=3/4\) の場合を適用すればよい。
\(1/2<\beta<1\) とし、\(\beta=1/2+9\eta\)、\(\eta>0\) と書く。
底本は、指定した長方形に少なくとも一つの捻りの零点を持つ
素数 \(d\le X\) の個数を \(O_{q,\eta,H}(X^{1-\eta/2})\) と評価する。
以下ではその評価を導く。

\(\mu\) は Möbius 関数とする。\(j\in\{0,1\}\) に対し
\(w_j(u)=u^j e^{-\pi u^2}\) とおき、その Mellin 変換を
\(\Gamma_j(s)=\int_0^\infty w_j(u)u^{s-1}\,du\) と書く。
係数列 \(c=(c(n))_{n\ge1}\) に対し
\[
 \theta_c(u)=u^{1/2}\sum_{n\ge1}c(n)w_j(nu),\qquad
 J(c,T)=\int_{1/T}^\infty
        \bigl(|\theta_c(u)|+|\theta_{\mu c}(u)|\bigr)\,\frac{du}{u}.
\]
\(\chi\) を導手 \(Q\le T^2\)、パリティ \(j\) の非主原始指標とする。
\(\Gamma_j(s)L(s,\chi)\) は \(\theta_\chi\) の \(s-1/2\) における Mellin 変換である。
\(\RePart s>1\) で \(\Gamma_j(s)/L(s,\chi)\) を与える \(\theta_{\mu\chi}\) の積分を
\(u>1/T\) に切ったものを \(M_T(s)\) とする。
\(\RePart s=1/2\) では \(|M_T(s)|\le J(\chi,T)\)。
テータ関数の関数等式は
\(\int_0^\infty |\theta_\chi(u)|\,du/u\) を
\(\int_{Q^{-1/2}}^\infty|\theta_\chi(u)|\,du/u\) の二倍にする。
\(Q^{-1/2}\ge T^{-1}\) なので
\(|\Gamma_j(s)L(s,\chi)|\le2J(\chi,T)\) を得る。

\(J(\chi,T)\le T^{2\eta}\) と仮定し、
\(\Psi(s)=\Gamma_j(s)L(s,\chi)M_T(s)-\Gamma_j(s)^2\) とおく。
左境界 \(\RePart s=1/2\) では \(|\Psi(s)|\ll T^{4\eta}\) である。
右境界 \(\RePart s=2\) では、絶対値が \(1\) 以下の係数列に対する
Mellin 被積分関数の上界 \(3e^{-u}\) により
\[
 |\Gamma_j(s)L(s,\chi)|\le3,\qquad
 \left|\frac{\Gamma_j(s)}{L(s,\chi)}-M_T(s)\right|\le 3/T.
\]
したがって \(|\Psi(s)|\le 9/T\) となる。三線定理に使う帯での
正則性と有界性は、\(\Gamma_j L\) の Mellin 表示と \(M_T\) の打切り積分から従う。
\(\RePart s\in[1/2+9\eta,1]\) では右境界に掛かる補間の重み
\(t=(\RePart s-1/2)/(3/2)\) は \(6\eta\) 以上である。
よって \(|\Psi(s)|\ll T^{4\eta(1-t)-t}\) であり、指数は \(-2\eta\) 以下である。
Euler 積により \(\RePart s>1\) には零点がない。
\(L(s,\chi)=0\) なら \(\Psi(s)=-\Gamma_j(s)^2\) である。
\(\Gamma_j\) の絶対値は \(1/2\le\RePart s\le1,\ |\ImPart s|\le H\) で
正の下界を持つから、十分大きな \(T\) について零点の存在は
\(J(\chi,T)>T^{2\eta}\) を強制する。
\leansrc{Lib/Analytic/QuadraticTwist}{\texttt{eventually\_rpow\_lt\_thetaTailIntegral}}。

次に \(n^\circ\) を \(n\) の奇数部分とし、\(|\alpha_n|\le1\) に対して
\(c_v(n)=\alpha_n(-v/n^\circ)\) とおく。ここでは括弧内は Jacobi 記号である。
\(u\ge T^{-1}\) のとき \(\sum_{1\le v\le T^2}|\theta_{c_v}(u)|^2\) を展開すると、
\(n,n'\) の組に付く指標和は
\[
  \sum_{1\le v\le T^2}\left(\frac{-v}{(nn')^\circ}\right)
\]
となる。奇数部分が平方数なら、その絶対値は \(T^2\) 以下である。
\(|w_j(x)|\ll_A x^{-A}\) に \(A=1/2+\eta\) を使うと、
その組の寄与は \(T^2u^{-2\eta}(nn')^{-1/2-\eta}\) の定数倍以下である。
これらの組では \(k nn'\) が平方数となる \(k\in\{1,2\}\) があり、
\(kn=ab^2,\ n'=ac^2\) と表せる。
したがってこの組だけにわたる \(\sum(nn')^{-1/2-\eta}\) は、
収束する三重和 \(\sum_{a,b,c\ge1}(abc)^{-1-2\eta}\) の定数倍以下である。

奇数部分が平方数でなければ、Pólya--Vinogradov の不等式を
Jacobi 指標に適用し、対数因子を吸収して指標和を
\(O_\eta((nn')^{1/2+\eta/2})\) と評価する。
今度は \(A=3/2+\eta\) を使うと \(u^{-2}\le T^2\) により
各組の寄与は \(T^2u^{-2\eta}(nn')^{-1-\eta/2}\) の定数倍以下となる。
この二重和も収束する。以上を合わせると
\[
 \sum_{v\le T^2}|\theta_{c_v}(u)|^2\ll_\eta T^2u^{-2\eta}.
\]
Cauchy--Schwarz の不等式の後に \(du/u\) で積分すると
\(\sum_{v\le T^2}J(c_v,T)\ll_\eta T^{2+\eta}\) となる。
\(\theta_{\mu c_v}\) の項には \(\alpha_n\mu(n)\) を使う。
\leansrc{Lib/Analytic/QuadraticTwist}{\texttt{sum\_thetaTailIntegral\_le}}。

\(d\equiv7\pmod8\) では二次相互法則と \((2/d)=1\) により
\((n/d)=(-d/n^\circ)\) である。各 \(\psi\bmod q\) を導く原始指標を
\(\psi^\ast\)、その導手を \(Q_0\mid q\) とすると、
\(\psi^\ast\chi_d\) は導手 \(Q_0 d\) の非主原始指標であり、
誘導で加わる有限個の Euler 因子は \(\RePart s>0\) で非零である。
\(T^2=qX\) とし、\(\alpha_n=\psi^\ast(n)\) と置く。二つのパリティをそれぞれ数えると、
零点を持つ素数の個数は \(O(T^{2+\eta}/T^{2\eta})=O(X^{1-\eta/2})\)
である。法 \(q\) の指標は有限個なので、全指標を合わせても同じ指数の評価が成立する。

もし剰余類内のすべての素数がこのような零点を持てば、
この個数評価から \(\sum_{p\equiv d_0\bmod q}(\log p)/p\) が収束する。
これは既約剰余類における同級数の発散に反する。
詳しくは \leansrc{Lib/Analytic/QuadraticTwist}
{\texttt{zero\_density\_quadTwist}, \texttt{exists\_entireL\_quadTwist\_ne\_zero}}。
\end{proof}

\subsection{片側評価が与える零点の排除}

\begin{lemma}[Landau の補題の適用形]\label{lem:landau}
可測な実関数 \(U(v)\) が \(v\ge0\) で \(|U(v)|\le C e^v\) を満たすとする。
その変換 \(T(s)=\int_0^\infty U(v)e^{(1-s)v}\,dv\) は、
\(\RePart s>0\) に高々単純極を持つ有理型接続を持つとする。
ある \(\beta>0\) について、十分大きな \(v\) で
\(U(v)\le e^{(\beta-1)v}\) が成り立ち、
すべての実数 \(\sigma>\beta\) で \(T\) の留数の実部が非負ならば、
\(T\) は \(\RePart s>\beta\) で正則である。
\end{lemma}

\begin{proof}
\(K e^{\beta v}-U(v)e^v\ge0\) が \(v\ge0\) で成り立つように
定数 \(K\) を取る。この非負の密度のラプラス変換を \(A(s)\) とすると、
右方では \(A(s)=K/(s-\beta)-T(s)\) である。
収束横座標を \(\sigma_c\) とし、\(\sigma_c>\beta\) と仮定する。
\(A\) が \(\sigma_c\) で正則なら、
\(\sigma_1>\sigma_c\) を十分近く取ったとき \(\sigma_1\) を中心とする
Taylor 級数はある実数 \(\sigma_2<\sigma_c\) で収束する。
その各項は、収束半平面での積分の微分により、密度に
\(e^{-\sigma_1 v}((\sigma_1-\sigma_2)v)^j/j!\) を掛けた積分に等しい。
各項が非負なので単調収束定理を使うと、
級数の和は \(\sigma_2\) におけるラプラス積分に等しく、有限である。
これは収束横座標の定義に反する。

ゆえに \(A\) は \(\sigma_c\) で極を持つ。ところが同点での留数は
\(T\) の留数の負であり、その実部は非正である。
単純極なら留数の実部は
\(\lim_{\sigma\downarrow\sigma_c}(\sigma-\sigma_c)A(\sigma)\)
で正となる。実際、この極限は非負な実数であり、単純極なので零ではない。
矛盾から \(\sigma_c\le\beta\) である。積分が正則な半平面
\(\RePart s>\beta\) では \(K/(s-\beta)-A(s)\) も正則である。
\leansrc{Lib/Analytic/Landau}{\texttt{IsPolar.landau}}。
\end{proof}

\begin{proposition}\label{prop:one-sided-criterion}
\(q,d_0\) を補題 \ref{lem:select-quadratic} のとおりとし、\(\beta>1/2\) とする。
各素数 \(d\equiv d_0\pmod q\) について、十分大きな実数 \(y\) で
\(R_q(y)\le y^{\beta-1}\) が成り立つと仮定する。
このとき、すべての \(\psi\bmod q\) について
\(\RePart\rho>\beta,\ \rho\ne1\) ならば \(L(\rho,\psi)\ne0\) である。
\end{proposition}

\begin{proof}
\(\rho\) を固定し、補題 \ref{lem:select-quadratic} を \(H=|\ImPart\rho|\) に
適用して \(d\) を選ぶ。補題 \ref{lem:prime-transform} の留数は、
\(\RePart s>\beta,\ |\ImPart s|\le H,\ s\ne1\) で
\(\nu(s)/\bigl(\varphi(q)s(1-s)\bigr)\) である。
とくに実軸の \(\beta<s<1\) では非負である。
\(s\ge1\) では整関数化した \(L^\circ(s,\psi)\) は非零であり、留数も零である。
\(U(v)=R_q(e^v)\) は \(O(e^v)\) であり、
仮定により十分大きな \(v\) で \(U(v)\le e^{(\beta-1)v}\) が成り立つ。
補題 \ref{lem:landau} から留数は \(\rho\) で零である。
分母は零でなく、\(\nu(\rho)\) は非負の零点位数の和なので各項が零である。
\(\rho\ne1\) では \(L^\circ\) の非零性は \(L\) の非零性を含意する。
\leansrc{New/Reduction/Zeros/Criterion}{\texttt{LFunction\_ne\_zero\_of\_primeRemainder\_le}}。
\end{proof}

\subsection{固定法に対する素数定理}

\begin{lemma}\label{lem:fixed-modulus-pnt}
固定した奇素数 \(d\) と \((d,q)=1\) なる正整数 \(q\) について、
任意の整数 \(B\ge0\) に対し十分大きな \(y\) で
\[
 |R_q(y)|\le(\log y)^{-B}.
\]
また \(z\ge y\ge1\) ならば \(|R_q(z)-R_q(y)|\le13(z-y)/y\) である。
\end{lemma}

\begin{proof}
\(\mathcal E_q(t)=\sum_{n\le t}\lambda_q(n)-t/\varphi(q)\) とおく。
\(\lambda_q\le2\Lambda\) と \(\sum_{n\le t}\Lambda(n)\le6t\) によって
\(|\mathcal E_q(t)|\le13t\) である。部分和公式より
\[
 R_q(y)=\int_1^y\mathcal E_q(t)\frac{dt}{t^2}-c_q.
\]
これから差の上界が従う。さらに \(0\le u\le v\le u+1\) で
\(|R_q(e^v)-R_q(e^u)|\le13(e^{v-u}-1)\le26(v-u)\) であり、
区間を長さ \(1\) 以下に分割すれば \(U(v)=R_q(e^v)\) は
\(v\ge0\) で Lipschitz 定数 \(26\) を持つ。

固定法の零点のない領域については \cite[第11章]{mv} も参照されたい。
ここでは底本の零点のない領域と対数微分の評価を
\(\psi\) と \(\psi\chi_d\) の有限族に同時に適用する。
\(F(s)=G_q(s)/s\) は、十分大きい \(\mathcal L\) について
\(\RePart s\ge1-\mathcal L^{-26},\ |s|\le e^{\mathcal L}\)
で正則で \(|F(s)|\le\mathcal L^{26}\) である。
これは各対数微分の評価に現れる指数 \(25\) を一つ増やして、
指標の個数と \(1/|s|\le2\) を吸収したものである。
\leansrc{Lib/Analytic/PrimeNumberTheorem/ZeroFreeRegion}{対数微分の評価}、
\leansrc{New/Reduction/Zeros/PrimeNumberTheorem}
{\texttt{eventually\_norm\_logDerivMean\_div\_le}}。

\(u\) を大きく取り、
\(k=\lfloor u^{2B+4}\rfloor,\ H=(k+1)/u,\ \mathcal L=(2B+5)\log u\)、
\(\eta=\mathcal L^{-26}\) とする。
\(\mathcal T_q(s)=(F(s)-F(1))/(s-1)\) の \(1+H\) での正規化した
Taylor 係数を
\[
 A=\frac{(-H)^k H}{k!}\mathcal T_q^{(k)}(1+H)
\]
と書く。積分を微分すると \(A\) は
\(H^{k+1}v^k e^{-Hv}/k!\) を密度とする \(U(v)\) の平均であり、
この密度の平均は \(u\)、分散は \(u^2/(k+1)\) である。
任意の \(\delta>0\) について
\(|v-u|\le\delta+(v-u)^2/\delta\) と Lipschitz 性を合わせ、
\(\delta=u^{-B-1}\) とすれば \(|A-U(u)|\le52u^{-B-1}\) となる。

中心 \(1+H\)、半径 \(H+\eta\) の円板は上の正則領域内にある。
円周では \(|s-1|\ge\eta\) なので
\(|\mathcal T_q(s)|\le 2\mathcal L^{26}/\eta\) である。
Cauchy の評価に \(H/(H+\eta)\le e^{-\eta/(2H)}\) を使うと
\[
 |A|\le4\mathcal L^{26}H\eta^{-1}
       \exp\!\left(-\frac{(k+1)\eta}{2H}\right)
      =4\mathcal L^{52}H e^{-u\eta/2}.
\]
十分大きな \(u\) では \(\mathcal L^{52}\le u\)、\(4H\le u^{2B+5}\)、
\(u^{3B+7}\le e^{\sqrt{u}/2}\) であり、したがって \(|A|\le u^{-B-1}\)。
\(|U(u)|\le53u^{-B-1}\le u^{-B}\) とおき \(u=\log y\) とすればよい。
\leansrc{Lib/Analytic/PrimeNumberTheorem/Tauberian}
{\texttt{eventually\_\allowbreak abs\_\allowbreak le\_\allowbreak inv\_\allowbreak pow\_\allowbreak of\_\allowbreak transform}}。
\end{proof}

\section{補助データと素数和の恒等式}\label{sec:data}

\subsection{数体、重さ、レベル}

\(q_0\) は \(18480\) の正の倍数、\(d\) は \(q_0/2-1\) に合同な
法 \(q_0\) の素数とする。\(K=\mathbb Q(\sqrt{-d})\) とおく。
このとき \(d\equiv7\pmod8\) である。前節の \(\chi_d\) は \(K\) の二次指標であり、以下 \(\chi_K=\chi_d\) とも書く。

底本の許容データ \(D\) には、正整数 \(x,x_1,b\) と階数を定める整数 \(m,r\)
が含まれる。\(a=xb\)、\(k=r+1-m\)、\(t=5r\) とする。
\(m\) は \(2\) の正の冪、\(r\) は \(3\) の正の冪であり、次の関係を満たす。
\begin{gather*}
 (\log b)^{100}\le x,\qquad
 x^{10^6}\le x_1,\qquad x_1^{10^6}\le b,\qquad
 (\log x_1)^6\le\log b,\\
 \sqrt{x}\le m\le 2\sqrt{x},\qquad x^5\le r\le 3x^5.
\end{gather*}
\(\log\log b\) は \(10^{12}C_{q_0,d}\) 以上を要求する。
ここで \(C_{q_0,d}\) は底本が \(q_0,d\) だけから定める定数である。\leanref{Parameters.constantBound}また \(2^s=m^{22}\) とすると \(s\le32\log x\) であり、
\[
  l_0=3^{\lfloor\log_3 s\rfloor+4},\qquad
  q_1=3l_0\,2^s q_0,\qquad q_2=(q_0d)^{s+6},\qquad
  E=K(\zeta_{q_1})
\]
とおく。\(b\) は \(q_1x_1\) の倍数である。
\(E\) の無限素点の集合 \(\Phi\) は \((\mathbb Z/q_1\mathbb Z)^\times\)
と同一視する。\(f=|\Phi|\) とし、法 \(q_0\) で \(1\) に合同な素点の集合を
\(\Phi_a\) とする。法 \(q_0\) で \(-1\) に合同な素点から
部分集合 \(\Phi_b\) を選び、
\[
  f=|\Phi_a|\varphi(q_0),\qquad
  |\Phi_a|\le|\Phi_b|\log b,\qquad
  \kappa(\tau)=
  \begin{cases}
    b=a/x,&\tau\in\Phi_b,\\
    a/x_1,&\tau\notin\Phi_b
  \end{cases}
\]
とする。

\(\Phi_b\) は次の疎性条件を満たす。任意の素数 \(\ell\le kt\) と
\(0<i<j<l_0\) に対し、\(\sigma,\tau,\rho\in\Phi_b\) が
\(\tau/\sigma\equiv\pm\ell^i\) および
\(\rho/\sigma\equiv\pm\ell^j\) を同時に満たすことはない。
合同は \(q_1\) の \(\ell\) と互いに素な部分を法とし、二つの符号は独立である。
この条件が、底本で \(\Phi_b\) に課す疎性である。\leanref{Parameters.IsSparse}

レベルを定める素数 \(p_0,p_1,N\) は \(q_2^2\) を法として \(1\) に合同である。
相異なる \(p_0,p_1\) は \([b^{10},b^{11}]\) に属し、\(M=p_0p_1\) とする。
\(K\) の整数 \(\pi_M\) と \(E\) の整数 \(\pi_N\) が
\(\pi_M\overline{\pi_M}=M\)、\(\pi_N\overline{\pi_N}=N\) を満たすように取る。
\(\bar\kappa=\exp(f^{-1}\sum_{\tau\in\Phi}\log\kappa(\tau))\) を重さの幾何平均とし、
\(\mathrm{rd}_E=|\operatorname{disc} E|^{1/[E:\mathbb Q]}\) を根判別式とする。
\(N\) の範囲は
\begin{equation}\label{eq:level-window}
 \log N\in[\log N_-,\log N_-+\sqrt{\log b}],\qquad
 \log N_-=
 m\log\bar\kappa+(m-1)\log\frac{M\mathrm{rd}_E}{2\pi}
 -\frac{15}{2}\log b.
\end{equation}
この位置でレベルを選ぶことで、後述の無限素点側の評価における
\(m\) 回分の増加を \(a\log N\) の項で相殺できる。

\subsection{許容データの選択}

\begin{proposition}\label{prop:admissible-data}
\(q_0,d\) を上のとおり固定し、\(0<\varepsilon<1/100\) とする。
十分大きな任意の整数 \(x\) に対して上の許容データが存在する。
\(x_1=\lfloor e^{x^{\varepsilon/6}}\rfloor\)、\(m\) は \(\sqrt{x}\) 以上の最小の
\(2\) の冪、\(r\) は \(x^5\) 以上の最小の \(3\) の冪に取れる。
\(b\) は \(e^{x^\varepsilon}\) より大きい最小の \(q_1x_1\) の倍数とでき、
\[
 x^{\varepsilon/6}-1\le\log x_1,\qquad
 x^\varepsilon\le\log b\le x^\varepsilon+1
\]
を満たす。
\end{proposition}

\begin{proof}
\(W=x^{\varepsilon/6}\) とおく。階数 \(m,r\) は主張どおりに選ぶ。
\(x_1=\lfloor e^W\rfloor\) について \(W-1\le\log x_1\le W\) である。
また \(q_1x_1\le e^{W^6}\) が十分大きな \(x\) で成り立つので、
指定の倍数 \(b\) は \(e^{W^6}<b\le2e^{W^6}\) を満たす。
これから \(\log b\le W^6+1\) を得る。
\(s=O(\log x)\)、\(W\) は \(x\) の正の冪であることにより、
許容データの対数スケールの条件が従う。

疎集合には Behrend の構成を使う。\leanref{exists_three_term_free}
座標を \(0,\ldots,H-1\) とする \(n\) 次元の箱の中で、二乗長の同じ
点を少なくとも \(H^n/(nH^2)\) 個選ぶ。
各点の座標を底 \(2l_0H\) の数字と読むことで指数の集合が得られる。
このとき、正整数 \(a_1,a_2\) が \(a_1+a_2\le2l_0\) を満たせば、
\[
 (a_1+a_2)e_2\equiv a_1 e_1+a_2 e_3 \pmod {2^s}
\]
は \(e_1=e_2=e_3\) を含意する。
実際、\((2l_0H)^n\le 2^s\) の下では両辺の整数代表は \(2^s\) 未満であり、
桁ごとの等式に繰り上がりはない。中間の点が両端の点の真の重み付き平均になるので、
同じ球面上にあるなら三点は一致する。

\(g=1+3l_0q_0\) とし、素点を \(-g^e\) の形に取る。
\(2^v\) を \(q_0\) を割る最大の \(2\) の冪とすれば、
\(g\) は法 \(2^{s+v}\) で位数 \(2^s\) を持つ。よって選んだ指数に対応する素点は
相異なり、すべて法 \(q_0\) で \(-1\) に合同である。
\(\ell=2\) の場合は比が法 \(3l_0\) で \(1\) になるが、
\(4\) の同じ法での位数は \(l_0\) なので禁じられた比は生じない。
奇素数 \(\ell\) については二つの比をそれぞれ \(2j,2i\) 乗して
法 \(2^{s+v}\) で比較すると
\(2j e_2\equiv2(j-i)e_1+2i e_3\pmod{2^s}\) となる。
指数の性質から \(e_1=e_2=e_3\) である。しかしこの場合
\(\ell^{2i}\equiv1\pmod{2^{s+1}}\) となり、底本の位数評価
\(2^s\le(\ell+1)2i\) に反する。実際、ここでは
\(2l_0(kt+1)<2^s\) を満たす。
整数 \(n\) を \(\varepsilon\) に応じて固定し、\(H\) は
\(H^n\le 2^s/W<(2H)^n\) を満たす \(2\) の冪に選ぶ。
底本の数値条件
\[
 W\le 2^s\le W^n,\qquad (2l_0)^n\le W,\qquad
 3l_0 n2^n\le W^3
\]
の下で \(H\le W\) であり、上の箱は収まる。選んだ素点の個数は
\(3l_0 2^s/W^6\) 以上である。
\(3l_0 2^s=|\Phi_a|\) および \(W^6\le\log b\) より必要な個数を確保できる。\leanref{Parameters.exists_isSparse}

レベルの素数は、Siegel--Walfisz 型の評価から得る
\(\log p\) の幅 \(W\) の窓に取る。
\(q_0,d\) は固定され、\(s=O(\log x)\) であるから、\(q_2^2\le W^{n+7}\) となるように
\(n\) を固定できる。\(W\le U\le W^{n+7}\) に対し
\([U,U+W]\) に \(\log p\) が入る素数を法 \(q_2^2\) の剰余類 \(1\) に取る。
使用する範囲は \([\log b,24m\log b]\) に含まれる。
\(h_K\) を \(K\) の類数とすると、
\([10\log b,11\log b]\) に \(h_K+1\) 個の互いに交わらない窓を並べられる。
そこから取った各素数 \(p\) は \(K\) で分裂し、上の二つの素イデアルを
\(I_p,J_p\) とすれば \([I_p][J_p]=1\) である。
同じ類にある \(I_{p_0},I_{p_1}\) を取ると \(I_{p_0}J_{p_1}\) が主イデアルになり、
生成元を \(\overline{\pi_M}\) と取る。
\(N\) については始点を \(\log N_-\) に取る。
\(W\le\sqrt{\log b}\) なのでこの窓は \eqref{eq:level-window} の範囲に収まる。
\(N\equiv1\pmod4\) から \(N=u^2+v^2\) と書け、
\(E\) には \(\sqrt{-1}\) が入るので \(\pi_N=u+v\sqrt{-1}\) とできる。
窓の存在とスケールの不等式の確認は、底本の許容データの構成を用いる。\leanref{eventually_exists_data}
\end{proof}

\subsection{持ち上げの始域と目標空間}

底本は \(E\) 上のエルミート空間 \(X_j,Y_j\) を用いる。
\(\dim_E X_j=k+j(t+1)\)、\(\dim_E Y_j=r+jt\) であり、
\(Y_m\subset X_m\) は余次元 \(1\) の定符号空間である。
\(X_0\) 上の代数的保型形式のうち、無限素点での重さが
\(\Phi_a\) では \((a,0,\ldots,0)\)、その他では \(0\) で、
指定のレベルを持つ空間を \(S_0\) とする。
\(M\) の上の選んだ素点で主合同レベルを課し、
\(\pi_N\) では最後の行を \((0,\ldots,0,1)\) に合同とする。
\(S\subseteq S_0\) は \(\pi_N\) を割る各素点 \(w\) について、
\(w\) とその共役で支持される有限アデールの中心スカラー作用の平均が
零になる形式の空間である。\leanref{ramifiedSourceForms}

一方、\(Y_0\) 上の主合同レベルの形式には、\(p_0,p_1\) の上で
\(\pi_M\) により選ぶ有限素点 \(w\) に対応した
\(\bar G=\prod_w\mathrm{GL}_r(\mathbb F_w)\) が作用する。
中心冪等元 \(e_0\) による射影の像を \(V\) とし、\(h=\dim_{\mathbb C}V\) と書く。
\(e_0\) は次のラベルを持つ既約指標の射影を足したものである。
各 \(w\) のラベルは次数 \(r\) のモニックな平方因子のない多項式で定数項は非零、
その剰余体にある根は \(q_2\) 乗根を含まない。
また \(p_0,p_1\) それぞれの上に、次数 \(r\) の拡大体の乗法群の生成元を根に持つ
原始多項式のラベルが少なくとも一つある。
\leansrc{New/UpperBound/Objects/LabelsAndIndexTypes}{有限群と中心冪等元}、
\leansrc{New/UpperBound/Objects/Rows}{目標空間とその無限型}。

捻り \(\xi\) ごとに \(Y_0\to\cdots\to Y_m\) の \(m\) 段の
テータ持ち上げがあり、その合成を \(V\) に制限して \(\Theta_\xi\) を得る。
第 \(j\) 段の指標は底本の固定指標 \(\mu_j^0\) に \(\xi^j\) を掛けたものである。
自明な捻りを \(1\) と書く。\(X_0\) から \(Y_0\) へのポアンカレ写像に対し、
第\ref{sec:upper-bound}節の正規化と射影を施した写像を
\(\mathfrak p_g\colon S\to V\) と書く。
\leansrc{New/UpperBound/Objects/RowChain}{持ち上げの連鎖}、
\leansrc{New/Reduction/IntegralTargetForms}{\(\Theta_\xi\) と整構造}。

\subsection{局所整除性に現れる素数和}

整数の \(\Phi\) における像を、\(q_1\) に素ならその剰余類、
そうでなければ単位元とする。この規約の下で \(-\nu\tau\) と書く。
\(u_+=\max(u,0)\) とし、\(0\le y\le a\) で
\begin{align*}
 s_{\mathrm{fin}}(y,\tau)
   &=\sum_{1\le\nu\le a}\lambda_1(\nu)
      \left(\frac{y}{\nu}-\kappa(-\nu\tau)\right)_+,\\
 s_{\infty}(j)&=\log(j!)+j\log(M/\Omega),\qquad
 \Omega=
    \exp\!\left(\gamma-\RePart\frac{L'}{L}(1,\chi_d)\right)
        \frac{2\pi}{\mathrm{rd}_E},\\
 s_{\mathrm{main}}&=a\bigl(\log(a/\bar\kappa)+c_1\bigr),\qquad
 \mathcal R=\sum_{\tau\in\Phi_a}\bigl(s_{\mathrm{fin}}(a,\tau)-s_{\mathrm{main}}\bigr),\\
 A_0&=\sum_{\tau\in\Phi}
       \bigl(s_\infty(\kappa(\tau))-s_{\mathrm{fin}}(\kappa(\tau),\tau)\bigr),
 \qquad H_0=mkA_0,\qquad E_j=fa(\log b)^j .
\end{align*}
\(\gamma\) は Euler 定数、\(c_1=G_1(1)\) は \eqref{eq:prime-remainder} の定数である。
使う量の非負性 \(H_0\ge0\) は、
各 \(\tau\) について
\(s_\infty(\kappa(\tau))-s_{\mathrm{fin}}(\kappa(\tau),\tau)\ge\kappa(\tau)\log b\)
であることから従う。後者は底本のスケールと
\(s_{\mathrm{fin}}(y,\tau)\le4y\log x_1\) を用いる。\leanref{primeSum_le}

\begin{proposition}\label{prop:total-remainder}
許容データについて
\begin{equation}\label{eq:total-remainder}
 \mathcal R =
 a\left[
 |\Phi_b|\bigl(R_{q_0}(x)-R_{q_0}(x_1)\bigr)
 +|\Phi_a|R_1(x_1)
 \right].
\end{equation}
\end{proposition}

\begin{proof}
\(-\nu\tau\in\Phi_b\) となる \(\tau\in\Phi_a\) の個数は、
\(\nu\equiv1\pmod{q_0}\) ならば \(|\Phi_b|\)、それ以外では零である。
実際、\(\nu\) が \(q_1\) に素なら乗法で素点の剰余類が移り、
\(\Phi_b\) は法 \(q_0\) で \(-1\) の剰余類にある。
\(\nu\) が \(q_1\) に素でない場合は \(-\nu\) の像は単位元である。
また \(q_1\) と \(q_0\) は素因数を同じくするので、この場合に
\(\nu\equiv1\pmod{q_0}\) となることはない。

したがって \(\sum_{\tau\in\Phi_a}s_{\mathrm{fin}}(a,\tau)\) を計算するときは、
まず \(|\Phi_a|\) 箇所すべてに重さ \(a/x_1\) を入れ、次いで
\(\nu\equiv1\pmod{q_0}\) の各項の \(|\Phi_b|\) 箇所を \(a/x\) の重さで補正すればよい。
\(\bigl(a/\nu-a/y\bigr)_+\) を \(\nu\le y\) の和に直せば、
\[
 \frac1a\sum_{\tau\in\Phi_a}s_{\mathrm{fin}}(a,\tau)
  =
 |\Phi_a|B_1(x_1)+|\Phi_b|\bigl(B_{q_0}(x)-B_{q_0}(x_1)\bigr),
 \qquad B_q(y)=R_q(y)+\frac{\log y}{\varphi(q)}+c_q .
\]
ここで差に現れる \(c_{q_0}\) は消える。一方、重さの定義から
\[
 |\Phi_a|\bigl(\log(a/\bar\kappa)+c_1\bigr)
 =|\Phi_a|(\log x_1+c_1)
   +\frac{|\Phi_b|}{\varphi(q_0)}(\log x-\log x_1).
\]
これが上の和の主項なので、差を取れば \eqref{eq:total-remainder} を得る。\leanref{totalRemainder_eq}
\end{proof}

\section{捻りつき格子と積公式}\label{sec:lattices}

\subsection{有限素点の整性と無限素点の大きさ}

\(\iota\) は、素数 \(\ell\) と体同型
\(\mathbb C\simeq\overline{\mathbb Q_\ell}\) の組を表す。
\(|z|_\iota=|\iota(z)|_\ell\) と書き、\(|\ell|_\ell=\ell^{-1}\) と正規化する。
複素共役を先に施した同型を \(\bar\iota\) と書く。
捻り \(\xi\) は無限素点で重さ \(0\) の \(U(1)\) の指標で、
分岐集合 \(Q_\xi\) は \(MN\) の上の素点から離れ、その共役と交わらない。
各 \(1\le j\le m\) で \(\xi^j\) は \(Q_\xi\) の各素点で分岐する。
法のノルム \(N_\xi\) が \(\ell\) と互いに素な場合、
捻り \(\xi\) は \(\ell\) で許容されるという。
\leansrc{New/Reduction/Data/LevelAndTwists}{捻りと許容条件}。

\begin{editorialnote}
添字 \(\iota\) は素数だけでなく体同型も指定する。
同じ \(\ell\) の上で \(\iota\) と \(\bar\iota\) を区別するのは、
エルミート内積の一方の変数に複素共役が入るからである。
\end{editorialnote}

\(W_\xi\) を \(\Theta_\xi\) の値域側の形式の空間とし、
その中の \(\iota\) で整な形式の加群を \(W^{\mathrm{int}}_{\xi,\iota}\) と書く。
この整構造は形式の各値の \(\ell\)-進 Weyl ゲージが \(1\) 以下であることによる。
\(\xi\) が \(\ell\) で許容されるなら、
\[
  \phi\in W^{\mathrm{int}}_{\xi,\bar\iota},\quad
  \psi\in W^{\mathrm{int}}_{\xi,\iota}
  \quad\Longrightarrow\quad
  |\langle\phi,\psi\rangle|_\iota\le1 .
\]
実際、内積は類を代表する点での Weyl 内積の和である。
このレベルでは類の安定化群が自明であり、
その位数で割ることによる分母は生じない。
各値の位置で局所行列を \(A\)、Weyl 表現を \(\rho_{\mathrm{W}}\)、整数 Gram 行列を \(G\) とする。
第一の変数に複素共役を取り、二つの変数にそれぞれ \(\rho_{\mathrm{W}}(A^{\mathsf T})\)、
\(\rho_{\mathrm{W}}(A^{-1})\) を作用させると、ゲージの仮定により両方の座標が整になる。
\(\rho_{\mathrm{W}}(A^{\mathsf T})^{\mathsf T}G\rho_{\mathrm{W}}(A^{-1})=G\) なので、
この整な座標で計算する内積はもとの内積の \(\iota\) での像に等しい。
\leansrc{Lib/Adelic/AutomorphicForms/PAdicGauge}{\texttt{gramForm\_mem\_embValRing}}。
自明な捻りでは類の代表における値を標準 Weyl 基底で読んだ座標により、
内積は有理数成分の正定値行列で表せる。
\leansrc{New/Reduction/IntegralTargetForms}{\texttt{Data.targetLattice}}。

\begin{definition}
\(y\in V\) が \(\iota\) で整であるとは、
\(\ell\) で許容されるすべての \(\xi\) について
\(\Theta_\xi y\in W^{\mathrm{int}}_{\xi,\iota}\) となることをいう。
有限族 \(y=(y_1,\ldots,y_n)\) の対数指数が \(\delta\) 以下であるとは、
次の条件を満たす非負の実数 \(\delta_\ell\) が存在することをいう。
\begin{enumerate}
\item 任意の有限素数集合 \(P\) に対し
\(\sum_{\ell\in P}\delta_\ell\log\ell\le\delta\)。
\item \(\ell\) の上の各 \(\iota\) について、\(\iota\) で整な成分を持つ
\(n\times n\) 行列 \(U_\iota\) があり、
\(yU_\iota\) の各成分は \(\iota\) で整、かつ
\(|\det U_\iota|_\iota\ge\ell^{-\delta_\ell}\)。
\end{enumerate}
同じ \(U_\iota\) が、その素数で許容されるすべての捻りに対して使える。
\end{definition}

自明な捻りの座標に \(\sigma\in\operatorname{Aut}(\mathbb C)\) を成分ごとに施し、
有理な内積行列で測った二乗ノルムを \(\|\Theta_1 y_i\|_\sigma^2\) と書く。
この座標と内積で計算した \(y\) の持ち上げのグラム行列式が、
すべての \(\sigma\) で \(e^c\) 以下なら、対数二乗余体積は \(c\) 以下という。
とくに \(\|\Theta_1 y_i\|_\sigma^2\le e^B\) がすべての \(i,\sigma\) で成り立てば、
Hadamard の不等式から \(c=nB\) が使える。
これらは \leansrc{New/Reduction/Lattices/TwistedLattice}
{\texttt{IndexLe}, \texttt{CovolumeLe}, \texttt{covolumeLe\_of\_normSqAt\_le}}
の条件である。

\subsection{高さの積公式}

\begin{definition}
族 \(u=(u_\iota)_\iota\) の高さが \(C\) 以下である
（\(\operatorname{ht}(u)\le C\) と書く）とは、
共通の乗数 \(H\in\mathbb C^\times\)、実数 \(\alpha\)、
非負の実数 \(e_\ell\) があり、すべての \(\sigma,\iota\) および
有限素数集合 \(P\) に対し
\[
 |\sigma H|\le e^\alpha,\qquad
 |H u_\iota|_\iota\le\ell^{e_\ell},\qquad
 \alpha+\sum_{\ell\in P} e_\ell\log\ell\le C
\]
となることをいう。
\end{definition}

二つの族を掛けるときは乗数を掛け、\(\alpha,e_\ell\) を足せばよいので、
高さの上界を加えられる。定数族 \(1\) の高さの上界は非負である。
このことを示すため、定数族 \(1\) の乗数 \(H\) を取る。
\(\ell=2\) で \(|H|_\iota\) が全 \(\iota\) にわたり有界なので \(H\) は代数的である。
実際、超越数なら同型を選んで任意の超越的な像に移せる。
\(\overline{\mathbb Q_2}\) では任意の上界より大きい絶対値を持つ開集合に
超越数がある。代数的数が可算で、その補集合が稠密だからである。
同型の延長に使う条件は
\leansrc{Lib/General/NumberTheory/PadicLimitsValuation}
{\texttt{isAlgebraic\_of\_norm\_complexToCp\_le}}
に示されている。

\(Q=\operatorname{Nm}_{\mathbb Q(H)/\mathbb Q}H\)、\(n_H=[\mathbb Q(H):\mathbb Q]\)
とおくと、\(Q\ne0\) であり、
\(\log|Q|\le n_H\alpha\)、\(\log|Q|_\ell\le n_H e_\ell\log\ell\)
である。各絶対値で \(Q\) を \(H\) の共役の積に展開した評価である。
分子と分母を割る素数をすべて含む \(P\) を取り、
\(\mathbb Q\) の積公式を使うと
\[
 0=\log|Q|+\sum_{\ell\in P}\log|Q|_\ell
   \le n_H\left(\alpha+\sum_{\ell\in P} e_\ell\log\ell\right)
   \le n_H C .
\]
したがって \(C\ge0\)。特に、\(z\ne0\) の定数族の高さが \(C\) 以下で
\(|\sigma z|\le e^{C'}\) が全 \(\sigma\) で成り立てば、\(C+C'\ge0\) である。
実際、\(z^{-1}\) の定数族は乗数 \(H=z\) を取れば高さ \(C'\) 以下となり、
二つの族の積が \(1\) になる。
\leansrc{New/Reduction/Lattices/Heights}{\texttt{HeightLe.nonneg\_of\_one}, \texttt{HeightLe.neg\_le\_of\_norm\_map\_le}}。

\subsection{基底に対するグラム行列式}

\(\Theta_\xi\) は単射であるとし、
\(Z(\xi)=\det(\Theta_\xi^*\Theta_\xi)\) とおく。
各 \(\iota\) で許容される捻りを一つ選んで \(\xi_\iota\) とする。
\(y=(y_1,\ldots,y_n)\) を一次独立とし、\(z=(z_1,\ldots,z_{n'})\) を補って
\(V\) の基底とする。\(w=(y,z)\) と書き、
\[
 G_\xi(w)=\bigl(\langle\Theta_\xi w_i,\Theta_\xi w_j\rangle\bigr)_{ij},
 \qquad \Delta_\xi(w)=\det G_\xi(w)
\]
とおく。

\begin{proposition}\label{prop:gram-product}
\(y\) の対数二乗余体積は \(c\) 以下、対数指数は \(\delta\) 以下とする。
各捻りについて \(\gamma(\xi)\ne0\) があり、
\(\gamma(\xi)\Theta_\xi z_j\) はすべての許容素点で整であるとする。
また、すべての \(\sigma,j\) に対して \(\|\Theta_1 z_j\|_\sigma^2\le e^{B_{\mathrm{Ar}}}\)
とする。族 \((|\gamma(\xi_\iota)|^{-2})_\iota\) と
\((Z(1)/Z(\xi_\iota))_\iota\) の高さの上界をそれぞれ \(B_{\mathrm{fin}},B_Z\) とすれば、
\begin{equation}\label{eq:gram-product}
  c+2\delta+n'(B_{\mathrm{Ar}}+B_{\mathrm{fin}})+B_Z\ge0.
\end{equation}
\end{proposition}

\begin{proof}
\(\bar\iota\) は \(\iota\) と同じ \(\ell\) の上にある。
\(y\) の指数を与える \(U_\iota,U_{\bar\iota}\) はともに、
その \(\ell\) で許容される \(\xi=\xi_\iota\) に対して使える。
\[
  P_\iota=\operatorname{diag}(U_\iota,\gamma(\xi)I_{n'}),\qquad
  P_{\bar\iota}=\operatorname{diag}(U_{\bar\iota},\gamma(\xi)I_{n'})
\]
とおく。
\(\overline{P_{\bar\iota}}^{\,t}G_\xi(w)P_\iota\) の各成分は
\(\bar\iota\) で整な形式と \(\iota\) で整な形式の内積なので、\(\iota\) で整である。
行列式を取ると
\begin{align*}
 1&\ge|\det P_{\bar\iota}|_{\bar\iota}
             |\Delta_\xi(w)|_\iota\,|\det P_\iota|_\iota,\\
 \bigl|\,|\gamma(\xi)|^{2n'}\Delta_\xi(w)\,\bigr|_\iota
     &\le\ell^{2\delta_\ell}.
\end{align*}
第一式から第二式に移る際、
\(\overline{\gamma(\xi)}\,\gamma(\xi)=|\gamma(\xi)|^2\) と
二つの指数行列の行列式の下界を使った。
従って族 \((|\gamma(\xi_\iota)|^{2n'}\Delta_{\xi_\iota}(w))_\iota\) の高さは
乗数 \(H=1\) を取って \(2\delta\) 以下となる。

基底 \(w\) は \(V\) を張るから、グラム行列式の変換則により
\(\Delta_\xi(w)=Z(\xi)\det(\langle w_i,w_j\rangle)_{ij}\) である。
各 \(\iota\) で
\[
 \Delta_1(w)
 =\bigl(|\gamma(\xi_\iota)|^{2n'}\Delta_{\xi_\iota}(w)\bigr)
     \bigl(|\gamma(\xi_\iota)|^{-2}\bigr)^{n'}
     \frac{Z(1)}{Z(\xi_\iota)} .
\]
ゆえに \(\Delta_1(w)\) の定数族の高さは
\(2\delta+n'B_{\mathrm{fin}}+B_Z\) 以下である。

無限素点では、座標を \(\sigma\) で移した全基底のグラム行列式に
Fischer の不等式を適用する。行列式は \(y\) と \(z\) の二つの対角ブロックの
行列式の積以下である。補助形式 \(z\) のブロックに Hadamard の不等式を用いると、
全基底のグラム行列式は \(e^{c+n'B_{\mathrm{Ar}}}\) 以下となる。
\(\sigma\) が複素共役と可換とは限らないので、
\(\sigma(\Delta_1(w))\) 自体の評価には外積を取る。
\(\sigma^\dagger(z)=\overline{\sigma(\bar z)}\) とすると、内積行列は有理なので、
\(\sigma(\Delta_1(w))\) は \(\sigma^\dagger\) と \(\sigma\) を用いた
二つの外積ベクトルの内積である。
それぞれの二乗ノルムが \(e^{c+n'B_{\mathrm{Ar}}}\) 以下であり、
Cauchy--Schwarz の不等式から
\(\left|\sigma(\Delta_1(w))\right|\le e^{c+n'B_{\mathrm{Ar}}}\) となる。

\(\Theta_1\) は単射なので \(\Delta_1(w)\ne0\) である。
その高さの上界と無限素点の上界を積公式で合わせれば
\eqref{eq:gram-product} を得る。
\leansrc{New/Reduction/Lattices/TwistedLattice}
{\texttt{covolume\_add\_index\_nonneg}}。
\end{proof}

\section{テータ持ち上げの上側評価}\label{sec:upper-bound}

分岐条件を課した始域 \(S\) から、テータ持ち上げの像の長さと
目標空間 \(V\) での指数を同時に制御できる形を取る。
指数に対しては各有限素点で許容される捻りをすべて扱い、
長さに対しては自明な捻りだけを考える。

\begin{theorem}[上側評価]\label{thm:upper-bound}
\begin{enumerate}
\item すべての有限素点で整であり、すべての \(g\in\overline G\) と
\(\sigma\in\operatorname{Aut}(\mathbb C)\) に対して
\[
  \|\Theta_1\mathfrak p_g\phi\|_\sigma^2
  \le \exp\!\bigl(H_0-m\mathcal R+E_4\bigr)
\]
を満たす \(\phi\in S\) が \(S\) を張る。
\item \(\mathcal T\subseteq\overline G\) を有限集合とする。
各 \(y_i\) が、ある \(g_i\in\mathcal T\) とすべての有限素点で整な
\(\phi_i\in S\) によって \(y_i=\mathfrak p_{g_i}\phi_i\) と書けるならば、
族 \((y_i)_i\) の対数指数は \(2|\mathcal T|(\dim S)E_3\) 以下である。
\end{enumerate}
\end{theorem}

\subsection{シーソー恒等式と射影}

主レベルの \(Y_0\) 上の形式には、\(M\) の上で選んだ素点の剰余体上の
一般線型群の積 \(\overline G\) が作用する。その群環の中心冪等元 \(e_0\)
が \(V\) への射影を与える。以下、群作用を \(\rho\)、ポアンカレ級数による
写像を \(\mathcal P\) と書く。主レベルからの下側の持ち上げ
\(Y_0\to\cdots\to Y_m\) を \(L_\xi\) と書くと、\(\Theta_\xi\) はその
\(V\) への制限である。

\emph{下側の積分を開く。}
\(F=E^+\) とし、ここでのアデールは \(F\) のアデールとする。
\(G_j=U(Y_j)\) とし、\([G_j]=G_j(F)\backslash G_j(\mathbb A)\) と書く。
\(J^\circ_j\) を有限側の標準格子の固定群、\(J_{\xi,j}\subset J^\circ_j\) を
捻り \(\xi\) におけるレベルとする。
下側のテータ核 \(K_{\xi,j}(g,h)\) は、シュワルツ・データを
\((g,h)\in G_j(\mathbb A)\times G_{j+1}(\mathbb A)\) で Weil 作用させた
テータ値である。下側の持ち上げ \(L_{\xi,j}\) は、
\(K_{\xi,j}(g,h)f(g)\) の \([G_j]\) 上の積分に
\(d_{\xi,j}=[J^\circ_j:J_{\xi,j}]\) を掛けたものとする。
共役型の形式 \(z\) に対して
\[
 L_{\xi,j}^{\mathrm t}z(g)=d_{\xi,j}
          \int_{[G_{j+1}]}K_{\xi,j}(g,h)z(h)\,dh
\]
と置くと、二重積分の順序を交換して
\[
  \int_{[G_{j+1}]} z\,L_{\xi,j} f
      = \int_{[G_j]}(L_{\xi,j}^{\mathrm t}z)f
\]
となる。これを \(j=m-1\) から \(0\) まで適用すると、
試験形式 \(x_m\) とのペアリングは、ポアンカレ級数
\(\mathcal P\phi\) と、\(\overline{x_m}\) を逆順に下ろして得た形式との積分に移る。
ここで積分の絶対収束には各段の急減少性を用いる。
\leansrc{New/UpperBound/Seesaw/SeesawIdentity}{持ち上げの転置と絶対収束}。

各 \(Y_j\) の次数付き分解から得るベッセル部分群を \(R_j\) とする。
\([R_j]\) 上の最初の関数 \(\mathcal W_0\) は、\(X_0\) 上での
\(\phi\) の行列係数を指定された指標でユニポテント部分に延ばしたものである。
先行する外側変数を固定して、\(j\) 段開いた関数を \(\mathcal W_j\) とする。
次の外側変数を \(g_j=(a_j,y_j)\)（無限側、有限側）に固定する。
\(j>0\) では \(W_{j,g_j}(r)=\mathcal W_j(rg_j)\)、
初段では \(W_{0,g_0}(r)=\mathcal W_0(r)\) と置き、
この関数と核の積を \([R_j]\) の \(r\) について積分する。
テータ和を開くと、次数を上げる写像に沿う定数項が残り、
これが次の関数 \(\mathcal W_{j+1}\) を与える。従って
\[
 \begin{aligned}
 \int_{[R_j]}W_{j,g_j}(r)(L_{\xi,j}^{\mathrm t}z)(rg_j)\,dr
   &=d_{\xi,j}\,\Xi_{j+1}(z),\\
 \Xi_{j+1}(z)
   &=\int_{R_{j+1}(F)\backslash G_{j+1}(\mathbb A)}
                   \mathcal W_{j+1}(h)z(h)\,dh .
 \end{aligned}
\]
\leansrc{New/UpperBound/Seesaw/UnfoldingStages}{ベッセル周期の一段の展開}。

初段では構成時の主レベル \(J_M\) の有限側の台を、
\(J_{\xi,0}\) の剰余類 \(\Omega_0=J_M/J_{\xi,0}\) に分解する。
ポアンカレ級数の積分を開くと、この剰余類の平均が得られるため、
初段に残る倍率は
\[
 \frac{d_{\xi,0}}{[J_M:J_{\xi,0}]}=[J^\circ_0:J_M]
\]
である。これは捻りによらない。
中間段 \(1\le j<m\) では、
\(\mathcal W_j\) が零でないところの有限成分の \(R_j(\mathbb A_f)\) 軌道を、
レベル \(J_{\xi,j}\) の有限個の剰余類の集合 \(\Omega_j\) が覆う。
\(c\in\Omega_j\) の重複度 \(m_j(c)\) は、局所格子の固定群を法として、
その軌道に沿って \(c\) をこの集合へ戻す類を数えたものである。
重み \(1/m_j(c)\) を置くと各軌道上で重みの積分が \(1\) になる。
\(J^\circ_j\) の体積を \(1\) としているので、
各 \(J_{\xi,j}\) の剰余類の体積は \(d_{\xi,j}^{-1}\) であり、
先ほどの周期に掛かる \(d_{\xi,j}\) を打ち消す。
無限側ではこの重み付けから得る測度で積分し、
有限側では \(\prod_{j<m}\Omega_j\) 上の和を取る。
初段の重複度は \(1\) で、他の重複度は
\(\mathrm NQ(\xi)\) の冪を割る。
ここで \(\mathrm NQ(\xi)\) は捻りの法のノルムである。
\leansrc{New/UpperBound/Seesaw/IntegralAsCosetSum}{軌道の重みとレベル体積の計算}。
\leansrc{New/UpperBound/Seesaw/Cut}{有限個の剰余類とその重複度}。

\emph{二段ずつ核を分解する。}
固定した外側変数について、連続する二段 \(Y_j,Y_{j+1},Y_{j+2}\)
（\(j\) は偶数）の定数項の積を考える。
第一の核の写像を随伴で反転し、第二の核と合わせると、
\(Y_{j+1}\to Y_j^{-}\oplus Y_{j+2}\) 上のデータを得る。
ここで \(Y_j^{-}\) はエルミート形式の符号を反転した空間である。
二つの次数付き列に沿う写像を \(T,T'\) とすると、
合わせたデータの平行移動は \((T^\ast,T')\) である。
次数を上げる写像に沿う二つの平均は、合わせたデータでの
次数零部分への射影に対応する。
外側の \([R_j]\) のユニポテント部分では \(\mathcal W_j\) と
定数項の指標が打ち消し合い、中間の \([R_{j+1}]\) の
ユニポテント部分では、残った双曲平面の等方座標を平均する。
このとき \(-TT^\ast+(T')^\ast T'=0\) なので、
二つの平行移動を合わせる際のヘイゼンベルクの位相は \(1\) となる。
従って残るのは、定数倍を除いて、
\(X_{j+1}\to X_j^{-}\oplus X_{j+2}\) 上のテータ核である。
これを \([U(X_{j+1})]\) で積分し、続いて \([U(X_j)]\) で
先行する行列係数を掛けて積分すると、上側の二段の合成の行列係数となる。
実際、無限側の Fock 多項式の積に対して核は二つのテータ核の積に分かれ、
二つの定符号群の積分がそのまま合成を与える。
一般の無限側の因子には、Fock 多項式の稠密性と積分の連続性により移る。
\leansrc{New/UpperBound/Seesaw/UnfoldingStages}{二つの定数項の積と定符号群上の二重積分}。

この計算では無限側と有限側が別々に作用する。
捻りの無限側の重みは固定されているので、無限側は行列係数を取る
線型形式だけを変え、その線型形式は捻り、\(\phi\)、有限側の外側変数によらない。
有限側は、外側変数の剰余類の組 \(\boldsymbol c\) における
二段の評価済みデータ \(\mathcal E_i(\xi;\boldsymbol c)\) を与える。
従って二段ずつの帰納法で、終段の \(\mathcal W_m\) は
\[
  \mathcal E(\xi;\boldsymbol c)
  =\bigotimes_{i=0}^{m/2-1}\mathcal E_i(\xi;\boldsymbol c)
\]
をデータとする \(X_0\to\cdots\to X_m\) の持ち上げの行列係数になる。

終段では \(Y_m\) は定符号で \(R_m=U(Y_m)\) であり、
その点は \(X_m\) の \(\operatorname{diag}(h,1)\) に入る。
行列係数を \(\overline{x_m}\) と掛けて積分する際、無限側の直交関係によって
整数係数の分岐写像を施した値の有限類平均を得る。
同じ平均を \(Y_m\) への制限の内積として書くと、
\(Y_m\) の Weyl 加群の次元と
\(\prod_\tau\tau(\overline{\pi_N})^{u_\tau}\) で割ったものになる。
ここで \(u_\tau\) は制限に用いる分岐指数である。
従って終段周期は、上側の持ち上げを制限したものと \(x_m\) の内積に、
外側の無限変数だけで決まる定数を掛けた形になる。
その無限側の積分と \([J^\circ_0:J_M]\) を合わせた定数を \(c_0\) とすると、
\(\langle x_m,L_\xi\mathcal P\phi\rangle\) は、
次のデータによる上側の持ち上げを制限して \(x_m\) と内積を取り、
\(c_0\) を掛けたものになる。
\[
 \mathcal A_\xi
   =\sum_{\boldsymbol c\in\prod_{j<m}\Omega_j}
         \left(\prod_{j<m}m_j(c_j)^{-1}\right)\mathcal E(\xi;\boldsymbol c).
\]
任意の \(x_m\) に対して等しいので、持ち上げ自体が等しい。
\leansrc{New/UpperBound/Seesaw/SeesawIdentity}{終段の分岐写像と上下の持ち上げの一致}。

\emph{有限データの分母をそろえる。}
\(\mathsf B=M^{f n_m^2}\), \(n_j=\dim_E X_j\) と置く。
各評価済みデータを有限フーリエ展開すると、接合した平行移動の和になる。
一組の二段に対する分母は \(\mathsf B^4\) を割り、
真空を整な真空に正規化する倍率は捻りによらない。
従って \(\mathcal A_\xi=\sum_t r_t C\,\mathcal V_t\) と書ける。
\(\mathcal V_t\) は接合した整な真空の平行移動で、
\(C\) は捻りによらず、\(\mathsf B^{2m} r_t\) は
\(\mathrm NQ(\xi)\) に素な場所で整である。
自明な捻りでは各評価済みデータの係数の絶対値和は \(1\) 以下であり、
各切断の和の倍率は高々 \(\mathsf B\) なので
\(\sum_t|\sigma(r_t)|\le\mathsf B^m\) となる。
レベルの作用は剰余類の組を置換し重複度を保つため、\(\mathcal A_\xi\) は
上側の中間レベルの同時作用で不変である。
\leansrc{New/UpperBound/Seesaw/GoodDatum}{接合したデータの有限フーリエ展開}。

\(c_0C/\mathsf B^{2m}\ne0\) ならこれを \(c_{\mathrm{ss}}\) とし、
\(\Psi_\xi(t)=\mathsf B^{2m}r_t\) とする。
積 \(c_0C\) が零なら全体の写像は零なので
\(c_{\mathrm{ss}}=1,\ \Psi_\xi=0\) と取ることもできる。
レベル \(M\mathrm NQ(\xi)\) に素な場所では、許容された平行移動を
零に置き換えても真空への作用は同じである。
こうして、対応する上側の合成持ち上げを \(B_{\xi,\Psi_\xi}\) とすれば
\begin{equation}\label{eq:upper-seesaw}
 L_\xi(c_{\mathrm{ss}}^{-1}\mathcal P\phi)=B_{\xi,\Psi_\xi}\phi
 \qquad(\phi\in S).
\end{equation}
平行移動は \(M\mathrm NQ(\xi)\) の上に支持を持ち、
\(\overline{\pi_M}\) を掛ければ \(\mathrm NQ(\xi)\) に素な場所で整である。
係数もその場所で整であり、自明な捻りでは
\(\sum_t|\sigma(\Psi_1(t))|\le\mathsf B^{3m}=M^{3fmn_m^2}\) である。
接合したテンソルは中間レベルの同時作用で不変である。\leanref{exists_isAdmissibleCombination_seesawDatum}

合成持ち上げの評価では、接合を \(N\) を法とする有限和として開く。
各項は、二段の整な真空の平行移動に
最後の列の \(-y/N\) と次の段の最後の行の \(-\overline y/N\) を加えたものになる。
項を集めると、係数の整性を保った接合前の平行移動データの和を得る。
平行移動の分母には \(N\) が許され、自明な捻りでの係数の絶対値和は
\(\exp(x^{60})\) 以下である。
\leansrc{New/UpperBound/Divisibility/ChainDatumExpansion}{接合テンソルの展開}。

射影を施した後の持ち上げは、\(L_\xi z\) の右平行移動の和で書ける。すなわち、
\(z\) を主レベルの形式、\(u\) を \(U(Y_m)\) の有限アデール点とすると
\begin{equation}\label{eq:upper-project}
 L_\xi\bigl(|\overline G|e_0\rho(g)z\bigr)(u)
 =
 \sum_{v\in\overline G}|\overline G|e_0(v)\,
            (L_\xi z)\bigl(u\,\widetilde{vg}_{\,\xi}\bigr).
\end{equation}
\(\widetilde{vg}_{\,\xi}\) は \(vg\) の第1ブロックへのレベル持ち上げである。
式 \eqref{eq:upper-project} は \(L_\xi\) の群作用との両立性から従う。
\(\overline G\) の既約指標による中心射影の表示から、
\(|\overline G|e_0(v)\) は代数的整数であり
\(|\sigma(e_0(v))|\le 1\) である。
右平行移動は整性の判定に使う格子を保つので、非アルキメデス的三角不等式
により、\(cL_\xi z\) が整なら \(cL_\xi(|\overline G|e_0\rho(g)z)\) も整である。
一方、各値の二乗ノルムの上界が \(\exp(B)\) なら、平方根を取った上界は
\(|\overline G|^2\exp(B/2)\) になる。自明な捻りでの類の個数は
\(\exp(x^{30})\) 以下であり、二乗して類和を取った上界は
\(\exp(B+x^{99})\) である。射影とレベルについて用いた仮定と数値評価は
\leansrc{New/UpperBound/Size/ArchimedeanSize}{射影による整性とノルムの評価}
に対応する。

\subsection{テータ級数の微分と有限素点の寄与}

\(K=\mathbb Q(\sqrt{-d})\) とし、有限素点では \(v_p(p)=1\) と正規化する。
\(p\) が \(K\) で非分解なら
\[
  \pi_p=\begin{cases}p & (p\ne d),\\ \sqrt{-d}&(p=d),\end{cases}
  \qquad
  \rho_p=\begin{cases}(p^2-1)^{-1}&(p\ne d),\\
                        (2p-2)^{-1}&(p=d).\end{cases}
\]
\(I_p^0\subseteq\mathbb Z_{\ge0}\) を、惰性素点では正の奇数の集合、
\(p=d\) ではすべての非負整数の集合とし、\(I_p=I_p^0\setminus\{0\}\)
とする。\(p\) が分解する場合には両集合を空とする。
\(\tau_e=(-p^e)\tau\) は無限素点の添字集合上の作用である。
ここでは、円分法において \(-p^e\) が単元を与えない場合に単位元で置き換える
という底本の規約を含め、この記号を用いる。

\paragraph{テータ級数の標準形}
\(p\) を非分解素数とする。以下では一つの \(p\)-進埋め込みを固定し、
その像を \(\mathbb C_p\) の中で書く。\(\mathcal O\) をその付値環とし、
\(q=p^2\)（惰性）、または \(q=p\)（分岐）とする。このとき
\(\rho_p=v_p(\pi_p)/(q-1)\) である。CM 楕円曲線のテータ級数を
\(\theta\) とし、正規化された局所パラメータを \(y=z+O(z^2)\) とする。
曲線の2等分点 \(t\) を原点に平行移動した座標
\(x_{\rm CM}(z)=\wp(z+t)-\wp(t)\) は \(\mathcal O[[y]]\) に属し、
\(x_{\rm CM}=z^2+O(z^3)\) である。

\emph{形式パラメータの拡大。}
CM 整数による作用は \(\mathcal O[[y]]\) を保つ。この作用を
\(p\)-進閉包に延ばすと、\(\zeta^{q-1}=1\) に対して
\(y(\zeta z)\in\mathcal O[[y]]\) を得る。
\[
 T(z)=\frac{1}{q-1}\sum_{\zeta^{q-1}=1}\zeta^{-1}y(\zeta z)
\]
と置く。この元は \(y\) のうち次数が \(1\bmod(q-1)\) の部分を取り出す。
\(q-1\) は単元で、\(T=y+O(y^2)\) の逆代入の各係数は次数ごとに
\(\mathcal O\) の中で解ける。したがって
\(\mathcal O[[T]]=\mathcal O[[y]]\) である。
\(T(\zeta z)=\zeta T(z)\) より、
\[
 T(\pi_p z)=\pi_p T(z)+\sum_{i\ge1}b_i T(z)^{1+i(q-1)}
 \qquad (b_i\in\mathcal O).
\]
剰余体上でこの級数の消滅次数は \([\pi_p]\) の非分離次数に等しく、
その次数は \([\pi_p]\) の次数 \(q\) 以下である。係数 \(\pi_p\) は剰余体で
零になり、次数 \(2,\ldots,q-1\) の項はない。従って \(b_1\) は単元である。
\leansrc{Lib/CM/FormalGroup/InseparableDegree}{CM 作用の還元と非分離次数}。

\(\varpi^{q-1}=-\pi_p/b_1\) を選び、\(Y(z)=\varpi^{-1}T(\varpi z)\) と置く。
\(v_p(\varpi)=\rho_p\) であり、上の式を代入して
\[
 \frac{Y(\pi_p z)}{\pi_p}=Y(z)-Y(z)^q+\pi_p H_1(Y(z)),
 \qquad H_1\in z^2\mathcal O[[z]]
\]
を得る。実際、\(i\ge2\) の項の係数は
\(b_i\varpi^{i(q-1)}/\pi_p\in\pi_p\mathcal O\) である。
\(Y(0)=0,\ [z]Y=1\) であり、次数 \(n\ge2\) では
\((\pi_p^{n-1}-1)[z^n]Y\) が既に求めた係数の整係数多項式で与えられる。
\(\pi_p^{n-1}-1\) は単元なので、帰納的に \(Y\in\mathcal O[[z]]\) となる。
\(\mathcal O/(\pi_p)\) は標数 \(p\) であり、還元した式は \(Y-Y^q=z\)
である。定数項零の解は次数ごとに一意に決まり、
\begin{equation}\label{eq:upper-formalparam}
 Y(z)\equiv\sum_{\nu\ge0}z^{q^\nu}\pmod{\pi_p}.
\end{equation}
また任意の \(H\in\mathcal O[[y]]=\mathcal O[[T]]\) は
\(T(\varpi z)=\varpi Y(z)\) に代入した後に \(H(\varpi z)\in\mathcal O[[z]]\)
を満たす。拡大前でも \(T=z+O(z^q)\) のため \(H\) の \(q\) 未満の係数は整である。
\leansrc{Lib/CM/FormalGroup/FormalParameter}{形式パラメータの構成と合同式}。

\emph{一変数の指数部。}
\(g=-\log\theta\) とする。\([\pi_p]\) による \(\theta\) の引き戻しは
次数 \(q\) の偶切断でねじれ点での整性を保つ。偶切断の展開を適用すると
\(\theta(\pi_p z)=\theta(z)^qR(x_{\rm CM}(z))\) で
\(2\deg R<q\) となる。ここで次数の狭義不等号は \(q\) が奇数であることによる。
\leansrc{Lib/CM/Integrality/IntegralSections}{偶切断の多項式表示}。
これから
\begin{equation}\label{eq:upper-frobenius}
 q g(z)-g(\pi_p z)=\pi_p H_2,\qquad H_2\in\mathcal O[[y]]
\end{equation}
を得る過程を述べる。\(g_0=g_1=0\) であり、\(g_n=[z^n]g\) とすれば
\(n<q\) について \(\pi_p g_n\in\mathcal O\) である。
\(n=2\) では \(\sqrt{-d}\,g_2\) の整性と
\(v_p(\sqrt{-d})\le v_p(\pi_p)\) を使う。
\leansrc{Lib/CM/Integrality/LogTheta}{テータ対数の二次係数の整性}。
\(n>2\) では \(g''=g''(0)+x_{\rm CM}\) だから
\(n(n-1)g_n=[z^{n-2}]x_{\rm CM}\) である。
右辺は \(n<q\) なら整であり、惰性の場合 \(v_p(n(n-1))\le1\)、
分岐の場合 \(v_p(n(n-1))=0\) となる。
\(q=\pi_p^2\) または \(q=-\pi_p^2\) なので、
\((q-\pi_p^n)g_n\in\pi_p\mathcal O\) が次数 \(q\) 未満で従う。

従って \(\exp(qg-g(\pi_p z))-1\) の \(q\) 未満の各係数は
\(\pi_p\) で割り切れる。実際、
\(v_p(j!)\le(j-1)/(p-1)\le (j-1)v_p(\pi_p)\)（\(j\ge1\)）であり、
指数の定数項は零である。この指数関数は \(R(x_{\rm CM})\) に等しい。
\(x_{\rm CM}^i\) の最低次項は \(z^{2i}\) なので、次数 \(2i\) の係数を
\(i=0\) から順に比較すると \((R-1)\) の全係数が
\(\pi_p\mathcal O\) に属する。よって \(R(x_{\rm CM})=1+\pi_p H\) で
\(H\in\mathcal O[[y]]\), \(H(0)=0\) となる。対数の第 \(j\) 項に現れる
\(\pi_p^j/j\) も \(\pi_p\mathcal O\) に属するので
\(\log(1+\pi_p H)\in\pi_p\mathcal O[[y]]\) である。
\leansrc{New/Derivatives/FormalGroup/FrobeniusCongruence}{テータ対数のフロベニウス合同式}。

\(g^\flat\) を次数 \(p+1\bmod(q-1)\) の部分とする。この次数類の
最小の非負代表 \(i_0\) は、惰性の場合 \(p+1\)、分岐の場合 \(2\) である。
CM 作用に対するテータの整性は、\(\zeta^{q-1}=1\) に対して
\[
 \exp\bigl(\zeta^{-(p+1)}g(\zeta z)-g(z)\bigr)\in\mathcal O[[y]]
\]
を与える。これは CM 整数での変換式をノルムと組にして
\((\zeta^{-1},\zeta^{-(p+1)})\) に延長し、変数を \(\zeta z\) に置き換えた式である。
\leansrc{Lib/CM/FormalGroup/FormalParameter}{単位根によるテータ級数の変換}。
これらの式をすべて掛けると、次数類を取る平均の公式により
\(\exp(g^\flat-g)^{q-1}\in\mathcal O[[y]]\) となる。定数項 \(1\) の
\((q-1)\) 乗根は、各次数で新しい係数に掛かる \(q-1\) を割って求められる。
従って \(\exp(g^\flat-g)\) 自身が \(\mathcal O[[y]]\) に属する。
\(c_\varpi=\varpi^{p+1}/p\) と置くと
\begin{equation}\label{eq:upper-one-variable}
 \begin{aligned}
 \theta(\varpi z)&=W_1(z)\exp(c_\varpi\Xi(z)),\\
 W_1(z)&=\exp\bigl((g^\flat-g)(\varpi z)\bigr)\in\mathcal O[[z]],\\
 \Xi(z)&=-g^\flat(\varpi z)/c_\varpi .
 \end{aligned}
\end{equation}

式 \eqref{eq:upper-frobenius} の次数 \(i_0\bmod(q-1)\) の部分を取る。
\(T(\zeta z)=\zeta T(z)\) なので、\(H_2\) を \(\mathcal O[[T]]\) で表して
同じ次数類を取ってよい。この部分は
\(\lambda' T^{i_0}+T^{i_0+q-1}H_3\)（\(\lambda'\in\mathcal O,\ H_3\in\mathcal O[[T]]\)）
である。変数を \(\varpi z\) に替えて \(-q c_\varpi\) で割ると、
\(-\pi_p\varpi^{i_0}/(q c_\varpi)\) は単元であり、
第二項には \(\varpi^{q-1}\in\pi_p\mathcal O\) が残る。従ってある
\(\lambda\in\mathcal O\) に対して
\[
 (1-\pi_p^n/q)[z^n]\Xi\equiv\lambda[z^n]Y^{i_0}\pmod{\pi_p\mathcal O}.
\]
\(n<i_0\) の係数は零である。\(n>2\) なら \(1-\pi_p^n/q\) は
\(\pi_p\) を法として \(1\) に合同な単元であり、分岐時の \(n=2\) では
この因子は \(2\) となる。従って \(\Xi\) は整であり、還元すれば
\[
 \Xi\equiv
 \begin{cases}
    \lambda Y^{p+1}& (p\text{ が惰性}),\\
    \lambda(Y^2-z^2/2)& (p=d)
 \end{cases}
 \pmod{\pi_p}.
\]
惰性の場合、式 \eqref{eq:upper-formalparam} を \(Y^pY\) に代入すると
指数 \(p^{2i+1}+p^{2j}\) が現れる。二つの (p) の冪指数の差の絶対値は正の奇数であり、
小さい方を \(\nu\)、差を \(e\) とすれば \((e,\nu)\) ごとに一組ずつ生じる。
分岐の場合、\(Y^2\) の対角項には係数 \(1\)、二つの指数の異なる項には
係数 \(2\) が付き、\(z^2\) の係数だけが \(1/2\) に変わる。こうして
\[
  \Xi(z)\equiv
  \sum_{e\in I_p^0}\sum_{\nu\ge0}\lambda_{e,\nu}(z^{p^e}z)^{p^\nu}
  \pmod{\pi_p}
\]
を得る。惰性の場合の係数はすべて \(\lambda\) である。分岐の場合は
\(\lambda_{0,0}=\lambda/2\)、\(\lambda_{0,\nu}=\lambda\)（\(\nu>0\)）、
\(\lambda_{e,\nu}=2\lambda\)（\(e>0\)）である。
\leansrc{New/Derivatives/FormalGroup/Exponent}{一変数の指数部の構成と係数合同式}。

\emph{多変数のテータ積。}
\(u_i\) を \(p\) の上のすべての素点で整なベクトルとし、
\(\mu_i\in\mathbb Z_{\ge0}\) は \(\sum_i\mu_i u_i u_i^\ast=1\) を満たすとする。
線型形式全体を \(L_i(z)=u_i^\ast z=\sum_{\tau,c}\overline{\tau(u_{i,c})}z_{\tau,c}\)
と書く。\(F\) を \(\prod_i\theta(L_i(z))^{\mu_i}\) と
\(y(L_i(z))\) の整係数級数の積とする。
式 \eqref{eq:upper-one-variable} を各 \(L_i\) に代入する。指数部の整理には
\(\overline{\tau(u_{i,c})}^{\,p^e}\equiv\tau_e(u_{i,c})\pmod{\pi_p}\)
を使う。整元については、円分体の生成元である単位根と \(\sqrt{-d}\)
でフロベニウスと \(\tau_e\) の作用を比べればよい。惰性の場合は
\(e\) が奇数なので \(\sqrt{-d}\) に対する共役と \(p^e\) 乗の符号が一致し、
分岐の場合は \(\sqrt{-d}\equiv0\pmod{\pi_p}\) となる。
\(p\) の上で整な元は \(p\) と素な整数分母を除いて整元にでき、
分母の \(p^e\) 乗は元の分母に合同なので、同じ結論が成り立つ。
したがって
\[
 \begin{aligned}
 \sum_i\mu_i L_i(z)^{p^e}L_i(z)
 &\equiv \sum_{\tau,c}z_{\tau,c}^{p^e}
              \sum_i\mu_i\tau_e(u_{i,c}) L_i(z)\\
 &=\sum_{\tau,c} z_{\tau,c}^{p^e}z_{\tau_e,c}\pmod{\pi_p}.
 \end{aligned}
\]
最後の等号は \(\sum_i\mu_i u_i u_i^\ast=1\) の \((\tau_e,c)\) 成分である。
さらに \(p^\nu\) 乗と和は還元後に交換し、\(\mu_i^{p^\nu}\equiv\mu_i\)
である。従って \(\sum_i\mu_i\Xi(L_i(z))\) と次の指数中の和との差は
\(\pi_p\mathcal O[[z]]\) に属し、定数項は零である。
\begin{equation}\label{eq:upper-normal}
 F(\varpi z)=W(z)\exp\!\left(
       \frac{\varpi^{p+1}}{p}
       \sum_{e\in I_p^0}\sum_{\nu\ge0}\lambda_{e,\nu}\Gamma_e(z)^{p^\nu}
                              \right),
 \qquad
 \Gamma_e(z)=\sum_{\tau,c}z_{\tau,c}^{p^e}z_{\tau_e,c}
\end{equation}
ここで \(W\in\mathcal O[[z]]\) である。実際、置換に伴う誤差を
\(c_\varpi\) 倍して指数関数を取ると、
\(v_p(c_\varpi\pi_p)=1/(p-1)\) と \(v_p(j!)\le j/(p-1)\) により整な因子となる。
これを \(\prod_i W_1(L_i(z))^{\mu_i}\) および \(\varpi z\) への拡大で
整となった \(y(L_i(z))\) の級数と掛ければ \(W\) を得る。
指数中の和は各係数について有限である。\leanref{exists_hasNormalForm}

\paragraph{小行列式の積に沿う微分}
\(Q\) を多項式とし \(J_Q(F)=(Q(\partial)F)(0)\) と置く。
\(Q\) が整係数の \(B\) 次斉次多項式、\(G\) の \(B\) 次の係数がすべて整で、
変数の個数が \(w\) ならば
\begin{equation}\label{eq:upper-legendre}
  v_p(J_Q(G))\ge
  \frac{B}{p-1}-w\bigl(1+\lfloor\log_p\max(B,1)\rfloor\bigr).
\end{equation}
実際、\(|D|=B\) の単項式に対する微分は係数に \(D!=\prod_c D_c!\)
を掛ける。Legendre の公式によれば
\(v_p(D_c!)=(D_c-s_p^{\mathrm{dig}}(D_c))/(p-1)\) で、
\(s_p^{\mathrm{dig}}\) は \(p\) 進表示の桁の和である。
各桁の大きさは \(p-1\) 以下で、桁数は
\(1+\lfloor\log_p\max(B,1)\rfloor\) 以下である。
各項の評価を足し合わせる際は非アルキメデス的三角不等式を使う。

\emph{シュール行列成分から小行列式の積へ。}
一段の始域と終域の階数を \(n=n_j,\ n'=n_{j+1}\) とする。
各場所 \(\tau\) の図形に高さが \(n\) と \(n'\) の間に狭義に入る列はない。
終域の標準充填 \(I'\) を固定し、\(Z_\tau=(z_{\tau,b,a})\) を
\(n'\) 行 \(n\) 列の変数行列、\(Y_\tau=(y_{\tau,a,k})\) を整な方向の行列とする。
\(Q_Y(Z)\) は \(Z_\tau Y_\tau\) の小行列式の積である。場所 \(\tau\) の高さ
\(h<n'\) の各列について、充填 \(I'\) が指定する \(h\) 行と最初の \(h\) 列を取る。
シュール行列の \((I',I)\) 成分を \(P_{I'I}(Z)\) と書くと、
シュール行列を定める標準充填の展開は
\[
 Q_Y(Z)=\sum_I B_I(Y)P_{I'I}(Z)
\]
を与える。\(B_I(Y)\) は始域の充填 \(I\) に対応する小行列式の積と、
始域の階数に等しい高さの列（\(n<n'\) のとき）から生じる共通の行列式因子との積である。
この共通因子は剰余体上でも零多項式でなく、標準充填の基底は剰余体上でも一次独立である。
従って \(B_I\) の剰余類は一次独立である。
\leansrc{New/Derivatives/Divisibility/Count}{シュール行列の展開と標準充填の基底}。

\(\alpha_I=J_{P_{I'I}}(F)\) とする。すべて零なら上界は自明である。
そうでなければ \(|\alpha_{I_0}|_p\) を最大に取ると
\(\sum_I(\alpha_I/\alpha_{I_0})B_I\) の還元は零多項式でない。
剰余体は無限なので、この還元が消えない点を選んで座標を \(\mathcal O\) に持ち上げる。
そこで得た \(Y\) について
\(J_{Q_Y}(F)=\sum_I\alpha_I B_I(Y)\) の絶対値は
\(|\alpha_{I_0}|_p\) となる。従ってすべての整な \(Y\) に対する
\(J_{Q_Y}(F)\) の付値下界を示せば、シュール微分の各成分にも同じ下界が従う。\leanref{exists_norm_schurJet_le_norm_jet_minorProd}

\emph{方向に課す合同条件。}
高さ \(h<n'\) の列数を \(N_{\tau h}\) とする。
高さ \(1\) の列数を \(t_\tau\)、残る列数を \(r_\tau\)、その箱数を \(B_\tau\) と書く。
全箱数は \(B=\sum_\tau(t_\tau+B_\tau)\) である。
印を付ける場所の集合を \(A\) とし、合同条件を数えるときだけ
\(\tau\in A\) において \(t_\tau(A)=0\) とする。それ以外では \(t_\tau(A)=t_\tau\) とする。
各 \(\tau,h\) で \(h\) 次の正方行列 \(g_{\tau h}\) を用いて方向を混ぜ、
\[
 \beta_{\tau h k,a}=\sum_{\ell<h}(g_{\tau h})_{k\ell}\,y_{\tau,a,\ell}
\]
とする。\(e\in I_p^0\), \(e\le\lfloor\log_p\max(B,1)\rfloor\) と
\(\tau'=\tau_e\) に対し、抑えるべき対は
\begin{equation}\label{eq:upper-direction-pairing}
 \sum_a\beta_{\tau h k,a}^{\,p^e}\beta_{\tau' h' k',a}\equiv 0
   \pmod{\pi_p}.
\end{equation}
例外は、\(h=1\) で、さらに \(\tau\in A\) または \(h'=1,\ \tau'\in A\) が成り立つ場合である。

課す合同条件は、\(g_{\tau h}\) の行に関して線型である。
\(\alpha_e^{p^e}=\pi_p\) とし、左側では法 \(\alpha_e\)、右側では法 \(\pi_p\) を使う。
ある行の全成分を左側で \(\alpha_e\) の倍数にすると、
新方向の各成分の \(p^e\) 乗が \(\pi_p\) で割れる。
右側では全成分を \(\pi_p\) の倍数にすればよい。
また左の高さが \(1\) なら、その新方向は元の \(y_{\tau,\cdot,0}\) の倍数なので、
右側の各行について
\(\sum_a y_{\tau,a,0}^{p^e}\beta_{\tau' h' k',a}\equiv0\pmod{\pi_p}\)
という一つの合同条件で足りる。右の高さが \(1\) のときは、各
\(y_{\tau',a,0}\) の \(p^e\) 乗根 \(\eta_a\in\mathcal O\) を選び、
\(\sum_a\eta_a\beta_{\tau h k,a}\equiv0\pmod{\alpha_e}\) を課せばよい。
この式を \(p^e\) 乗すると、\(\pi_p\mid p\) のため還元後に和と累乗が交換し、
\(\sum_a y_{\tau',a,0}\beta_{\tau h k,a}^{p^e}\equiv0\pmod{\pi_p}\) となる。
このようにペアリングの合同条件は、反対側の元の方向だけで定まる。

条件は次のように割り当てる。

\begin{enumerate}
\item 左の高さが \(1\) の対には、左側の高さ \(1\) の全成分に条件を課すか、
右側に元の左の方向とのペアリング条件を課すかを選ぶ。
費用はそれぞれ \(t_\tau(A)/p^e\) と \(r_{\tau'}+t_{\tau'}(A)\) であり、
小さい方を取る。選んだ側で印を付けた高さ \(1\) の列は除く。
\item 左の高さが \(1\) より大きく右の高さが \(1\) の対には、
左側に元の右の方向とのペアリング条件を課す。費用は \(r_\tau/p^e\) である。
\item 両側の高さが \(1\) より大きい対には、一方の側の該当する行に
全成分の条件を課す。費用は \(B_\tau/p^e\) と \(B_{\tau'}\) の小さい方である。
\end{enumerate}

費用の単位は \(v_p(\pi_p)\) である。各高さの正方行列について、
ペアリング条件は一つの線型形式、全成分の条件は高さの数だけの線型形式に対応する。
左側の法は \(v_p(\alpha_e)=v_p(\pi_p)/p^e\) を持つので \(p^{-e}\) を掛ける。
ここで全成分の条件を数える際に掛かる高さは線型形式の数である。
各形式を行列の全行で整除させる際に必要な行列式の損失は、その法一つ分で済む。

実際、一つの線型形式に対して行ごとの値を \(a_k\) とする。
零でない値があれば付値が最小の \(a_{k_0}\) を選び、
他の各行からその行の \(a_k/a_{k_0}\) 倍を引く。
この商は整で、他の行の値は零になる。残る行を法で掛けると、
行列式はその法だけを余分に持つ。値がすべて零なら操作は不要である。
単位行列から始めて線型形式を一つずつ処理する。
新しい行は前の行の整係数結合なので、すでに課した整除は保たれる。
従って \(\det g_{\tau h}\) は、その場所と高さに集めた法の積を割る。
係数行列は同じ高さの各列に一度ずつ現れるので、
\[
 \begin{aligned}
  \sum_{\tau}\sum_{h<n'}N_{\tau h}v_p(\det g_{\tau h})
   &\le v_p(\pi_p)\mathcal D_p(A),\\
  \mathcal D_p(A)
   &=\sum_{\substack{e\in I_p^0\\e\le\lfloor\log_p\max(B,1)\rfloor}}
       \sum_{\tau\in\Phi}
       \biggl\{\min\!\left(\frac{t_\tau(A)}{p^e},\,r_{\tau_e}+t_{\tau_e}(A)\right)\\
   &\hphantom{=\sum_e\sum_\tau\biggl\{}{}
          +\frac{r_\tau}{p^e}
          +\min\!\left(\frac{B_\tau}{p^e},\,B_{\tau_e}\right)\biggr\}.
 \end{aligned}
\]
以上で例外を除く式 \eqref{eq:upper-direction-pairing} がすべて同時に成り立つ。\leanref{exists_newDirection_dvd}

\emph{連鎖律による付値評価。}
新変数 \(s_{\tau h k b}\) により
\[
 z_{\tau,b,a}(s)=\sum_{h<n'}\sum_{k<h}s_{\tau h k b}\,p^{w_{\tau h b}}\beta_{\tau h k,a}
\]
と置く。印がなければ \(w_{\tau h b}=0\) とする。
印があるときは、各例外の対と各行 \(b\) について
\(w_{\tau 1 b},w_{\tau' h' b}\in\mathbb Z_{\ge0}\) の少なくとも一方が正であるようにする。
すると \(\Gamma_e(z(s))\) は \(\pi_p\) で割り切れる。
実際、還元後に \(p^e\) 乗を和に分配すると、各項の係数は
\eqref{eq:upper-direction-pairing} のペアリングに
\(p^{p^e w_{\tau h b}+w_{\tau' h' b}}\) を掛けたものである。
例外の対もこの因子により消える。

\(Q_*(s)\) を各列について、充填 \(I'\) が与える行 \(b\) と
\(k<h\) を添字とする小行列式 \(\det(s_{\tau h k b})\) を掛けたものとする。
\(W_{\rm ch}\) を各箱の行 \(b\) について \(w_{\tau h b}\) を足した和とする。
連鎖律によって
\begin{equation}\label{eq:upper-minor-chain-rule}
 J_{Q_*}\bigl(F(\varpi z(s))\bigr)
 = \varpi^B p^{W_{\rm ch}}\!
    \left(\prod_{\tau,h<n'}(\det g_{\tau h})^{N_{\tau h}}\right)J_{Q_Y}(F).
\end{equation}
一列だけを見れば、新方向の微分行列は \(Y_\tau\) の方向に沿う微分行列に
右から \(g_{\tau h}^{\,t}\) を掛け、左から行因子
\(\varpi p^{w_{\tau h b}}\) の対角行列を掛けたものである。
行列式の積を取れば \eqref{eq:upper-minor-chain-rule} となる。

\(Q_*\) は整係数で次数 \(B\)、変数は高々 \(f(n')^3\) 個ある。
式 \eqref{eq:upper-normal} に代入した際、条件を課していない
\(e>\lfloor\log_p\max(B,1)\rfloor\) の項は次数が \(B\) を越える。
残る結合項には \(\pi_p\) が掛かるので、先ほどの指数関数の付値計算により
\(F(\varpi z(s))\) の次数 \(B\) までの係数は整である。
式 \eqref{eq:upper-legendre} と \eqref{eq:upper-minor-chain-rule} から
\begin{equation}\label{eq:upper-jet}
 \begin{aligned}
 v_p(J_{Q_Y}(F))+W_{\rm ch}
   &\ge \gamma_p B-v_p(\pi_p)\mathcal D_p(A)\\
   &\quad{}-f n_{j+1}^3\bigl(1+\lfloor\log_p\max(B,1)\rfloor\bigr),\\
 \gamma_p&=\frac1{p-1}-\rho_p .
 \end{aligned}
\end{equation}
を得る。\(B\) には印を付けた列の箱も含まれている。\leanref{norm_jet_minorProd_le_of_charges}
\(e>0\) について
\(\gamma_p=v_p(\pi_p)\sum_{e\in I_p}p^{-e}\) である。
\(\gamma_p B\) から \(\mathcal D_p(A)\) の \(e>0\) の控除を引くと、
\(v_p(\pi_p)\) を単位として、次の正部が利得に現れる。
\[
 \left(\frac{t_\tau(A)}{p^e}-t_{\tau_e}(A)-r_{\tau_e}\right)_+
 +\left(\frac{B_\tau}{p^e}-B_{\tau_e}\right)_+
 -\frac{r_\tau}{p^e}.
\]
実際、\(x-\min(x,y)=(x-y)_+\) を二つの最小値に適用している。
控除を \(I_p\) 全体に延ばしても非負の項を引くだけなので下界が得られる。
さらに印を付けた各箱の \(\gamma_p\) を加える。\(p=d\) では \(e=0\) の控除も引く。

\paragraph{持ち上げ一段の利得}
局所素数和と全段階の利得を
\begin{align}
 s_p(y,\tau)
  &=v_p(\pi_p)\sum_{e\in I_p}
       \left(\frac{y}{p^e}-\kappa((-p^e)\tau)\right)_+,
       \label{eq:upper-sp}\\
 G_p
  &=m\left\{\sum_{\tau\in\Phi_a}s_p(a,\tau)
           +k\sum_{\tau\in\Phi}s_p(\kappa(\tau),\tau)\right\}
       \label{eq:upper-gain}
\end{align}
とする。分解する \(p\) には \(s_p=G_p=0\) とする。
\(0\le y\le a\) ならば、\(p\ge x_1\) でもすべての正部が零になる。
\(p\) が惰性のとき \(1-\chi_K(p^e)\) は奇数 \(e\) で \(2\)、
偶数 \(e\) で零である。\(p=d\) ではこの係数が \(1\) で
\(v_p(\pi_p)=1/2\) である。従って素数冪をまとめると
\begin{equation}\label{eq:upper-gain-sum}
  2\sum_p G_p\log p =
    m\sum_{\tau\in\Phi_a}s_{\mathrm{fin}}(a,\tau)
       +mk\sum_{\tau\in\Phi}s_{\mathrm{fin}}(\kappa(\tau),\tau).
\end{equation}

高さ 1 の列数は \(\Phi_a\) で \(a\)、その他で零である。
残る列数は \(\kappa(\tau)+x^{60}\) 以下、その箱数は
\(k\kappa(\tau)\) と \(k\kappa(\tau)+x^{60}\) の間にある。
これを上の正部に代入する。一段での付値評価を \(D_{p,j}(A)\) と書く。
ここで \(A\subseteq\Phi\) は高さ 1 の列を印付けする場所であり、
\(D_{p,j}(A)\) は印を除いた図形の予測付値に除いた箱の
\(\gamma_p\) を加え、微分の対数誤差を差し引いた量である。
\(p<x_1\) が非分解のとき
\begin{equation}\label{eq:upper-one-step}
 D_{p,j}(A)\ \ge\ \frac{G_p}{m}
  - C_p(A)
  +v_p(\pi_p)|\Phi_a\cap A|\,\frac{a}{x_1}
  -\frac{fa(\log b)^2}{mp},
 \qquad
 C_p(A)=\sum_{\substack{\tau\in\Phi,\ e\in I_p\\
                 \tau,\,(-p^e)\tau\in\Phi_a\setminus A}}\frac{a}{p^e}.
\end{equation}
この比較では、対応先にも高さ 1 の列が残っていれば、その重なりに
高々 \(a/p^e\) を差し引く。高さ 1 以外の列の誤差は上の \(x^{60}\)
の範囲であり、\(p=d,e=0\) の控除は \(\tau,-\tau\) がともに
\(\Phi_a\) に属することも、ともに \(\Phi_b\) に属することもないことから抑える。
一場所の誤差は \(a/(mp)\) を単位として、
正部の比較に高々 \(2\)、分岐素点の控除に高々 \(2\)、微分の対数誤差に
高々 \(1\) である。これらは \(5\le(\log b)^2\) で吸収できる。\leanref{le_divisibilityBound}

\paragraph{局所核の級数と格子}
一段 \(X_j\to X_{j+1}\) の局所核には、\(n'\times n\) 行列を集めた
標準空間のシュワルツ・データが入る。レベル \(1\) の標準データとは、
整な真空ベクトルを \(p\) の上で整な点で平行移動し、
それらを \(p\)-進整な係数で有限和にしたものである。
非分解素数 \(p\) は円分法 \(q_1\) を割らない。
以下では任意のレベル \(1\) の標準データのテータ級数に
式 \eqref{eq:upper-normal} が成り立つことを示す。

\emph{真空のテータ積展開。}
自己双対格子 \(\mathfrak a_0\) の元 \(v_i\) と \(p\) に素な整数 \(b_0\) を取り、
\(\sum_i \tau(v_i)\overline{\sigma(v_i)}=b_0^2\delta_{\tau,\sigma}\)
を満たすようにする。底本のタイトフレームの構成では、この条件を使って
各座標に \(v_i/b_0\) を置いた線型形式上の \(\theta\) の積を作る。
これに対応する Weil ベクトル \(F_0\) のテータ値は \(1\) で、
そのテータ関数は \(b_0\mathfrak a_0\) の各座標方向に指定された
符号付き周期を持つ。
\leansrc{New/Derivatives/Vacuum/TightFrames}{タイトフレームからテータ積を作る構成}。

整な真空を \(F_0\) の平行移動で表すには、この部分格子に対する
有限フーリエ反転を用いる。各剰余類の代表を \(u\) とすると、
フーリエ係数に真空を掛けたものは、\(F_0\) を \(-u\) で平行移動して
さらに周期の代表で平行移動した和で表される。
一方、反転公式にテータ汎関数を適用すると、これらのフーリエ係数を
対応する真空平行移動のテータ値で重み付けした総和は
\((b_0^2)^{f n n'}\) となる。
これによって整な真空は \(F_0\) の平行移動の有限和になる。
係数の分母は \((b_0^2)^{f n n'}\) で、分子は整な真空の値と
ヘイゼンベルク作用の位相因子の積であるため、\(p\) で整である。
現れる平行移動は \(b_0\) 倍で標準格子に入る。
\leansrc{New/Derivatives/Vacuum/StandardDataProducts}{真空の有限フーリエ反転と分母}。

標準データに現れる各平行移動をこの展開に適用する。
ヘイゼンベルク作用の積は平行移動の和に位相因子を掛けたものであり、
\(p\) での整性が保たれる。各平行移動は \(F_0\) への同じ作用を与える
有理代表で置き換えられる。その代表 \(\gamma\) は
ある \(p\) に素な整数 \(a_0\) を掛けると \(\mathfrak a_0\) の各座標に入る。
\leansrc{New/Derivatives/Vacuum/StandardDataProducts}{整な分母を持つ有理代表の選択}。
\(\gamma\) のテータ積に現れる各平行移動点 \(t_i\) は
\(a_0 b_0 t_i\in\Omega_0\mathcal O_K\) を満たす。
実際、各 \(t_i\) は \(-\Omega_0/b_0\) と、\(v_i\) と
\(\gamma\) の該当座標とのエルミート対の積である。
\(a_0\gamma\) と \(v_i\) は自己双対格子にあるので、その対は
\(\mathcal O_K\) に入る。有限個の平行移動について、すべての
\(a_0 b_0\) の倍数で \(p\) に素な整数 \(k_0\) を一つ取る。

各座標に \(v_i/(b_0 k_0)\) を置いたベクトルを \(u_i\) とし、
重複度を \(k_0^2\) とする。上のフレーム恒等式を \(b_0^2\) で割ると
\(\sum_i k_0^2u_i u_i^\ast=1\) を得る。
\(k_0 t_i\) が周期であるため、平行移動した一変数テータを
\(\theta(z)^{k_0^2}\) と \(\mathcal O[[y]]\) の元の積に分けられる。
ここでは \([k_0]\) による引き戻しを用い、級数の引数を \(k_0\) 倍した
平行移動テータの次数は \(k_0^2\) となる。
\leansrc{Lib/CM/Integrality/IsogenyIntegrality}{ねじれ点での平行移動テータの引き戻し}。
\(k_0\) をすべての項で共通に選んだので、有限和でも
\(\prod_i\theta(u_i^\ast z)^{k_0^2}\) をくくり出せる。
残る因子は \(y(u_i^\ast z)\) の整な級数環に属する。

この構成で整と分かるのは、まず \(u_i\) のアデール像
\(\alpha^{-1}u_i\) である。標準格子を自己双対にするイデール \(\alpha\) は
\(q_1\alpha\) が整であり、\(p\nmid q_1\) より \(p\) の上で \(\alpha\) は整である。
従って \(u_i=\alpha(\alpha^{-1}u_i)\) は \(p\) の上のすべての素点で整になる。
以上のテータ積展開は、テータ関数の変数を
\(\sqrt{2\pi/\sqrt d}/\Omega_0\) 倍した規格化で得られる。
実際の核ではこれにさらに周期因子と \(\pi_M\) を掛ける。
ここで \(\pi_M\overline{\pi_M}=M\)、周期因子は \(p\nmid q_1\) で単元である。
この追加のスカラーを \(c\) とすると \(v_p(c)\ge0\) である。
式 \eqref{eq:upper-normal} で変数を \(c\) 倍すると
\(W(cz)\) は整であり、指数係数には \(c^{(p^e+1)p^\nu}\) が掛かる。
従って実際の核の規格化でも標準形が保たれる。
\leansrc{New/Derivatives/Vacuum/Divisibility}{局所核のテータ級数の規格化と整性}。

\emph{格子座標での評価。}
標準充填による Weyl 加群の基底を使う。有限アデール点 \(g\) におけるベクトル
\(\eta\) の局所ノルム \(\|\eta\|_g\) は、
\(g^{-1}\) を重みの表現で作用させたベクトルの \(p\)-進座標の最大値である。
各 \(\tau\) の重み作用は埋め込み \(x\mapsto\iota_p(\tau(\overline x))\) を使う。
このノルムは格子を保つ点の右乗法で変わらず、有理点を
有限アデール点とベクトルの両方に作用させても変わらない。

\(\psi\) を標準データとし、\(K_j(\psi;g,h)\) を
捻った Weil 作用 \(\omega_j(g,h)\psi\) のシュール微分行列とする。
ここで \(g\) は \(X_j\)、\(h\) は \(X_{j+1}\) の有限アデール点である。
標準データを \(\omega_j(1,x d_\ast y)\psi\) に取り替えた場合も考える。
\(x,y\) は \(p\) の上で整な点とし、\(d_\ast\) は
各重み埋め込みで成分 \(u_{\tau,b}\) の対角行列になるとする。
弱近似により有理点 \(\gamma,\epsilon\) を選び、
\(g_0=\gamma g,\ h_0=\epsilon h x d_\ast\) と
\(d_\ast h_0 d_\ast^{-1}\) を \(p\) の上で整にする。
\(g_0,h_0 y\) は標準格子を保つ。したがってそれらによる Weil 作用は
真空の整な平行移動の張る加群を保ち、捻りの指標の値も単元である。
よって \(\psi'=\omega_j(g_0,h_0 y)\psi\) は標準データである。

核の共変性と局所ノルムの不変性により、
入力の点 \(g\) と出力の点 \(h\) での比較は、
標準データ \(\psi'\) を点 \(1\) で評価した微分行列を用い、
入力の点を \(1\)、出力の点を \(d_\ast^{-1}\) とした比較に移る。
実際、入力側の \(g_0\) は整であり、出力側では
\(\epsilon h=d_\ast^{-1}(d_\ast h_0d_\ast^{-1})x^{-1}\) の右の二因子が整だからである。
このとき充填 \(I'\) の行には、各箱の添字 \(b\) に対する \(|u_{\tau,b}|_p\) と、
終域の重みの末項を \(e'_\tau\) とした行列式因子
\(\prod_\tau(\prod_b|u_{\tau,b}|_p)^{e'_\tau}\) が掛かる。
行列成分の最大絶対値が作用素ノルムを抑えるので、
これらの積とシュール微分成分の積を評価すればよい。
\leansrc{New/UpperBound/Kernels/ThetaLiftGauge}{捻った Weil 作用と格子座標での核の評価}。

印なしの標準データでは \(d_\ast=1\) とすれば行因子はすべて \(1\) である。
式 \eqref{eq:upper-jet} で \(W_{\rm ch}=0\) とし、控除をすべての \(e\) に延ばし
対数の床を取らない形にゆるめると、各成分の付値は \(D_{p,j}(\varnothing)\) 以上である。
よって
\begin{equation}\label{eq:upper-kernel-gauge}
 \|K_j(\psi;g,h)\eta\|_h
 \le p^{-D_{p,j}(\varnothing)}\|\eta\|_g
\end{equation}
がすべての \(g,h,\eta\) に対して成り立つ。
\leansrc{New/UpperBound/Divisibility/ValuationAtNonSplitPrimes}{一段の核の付値評価}。

\emph{合成における付値の加算。}
各段階のデータを \(\psi_{t,j}\) とし、核のノルムの倍率を \(p^{c_j}\)
で抑える。これは \(\Psi(t)\ne0\) のすべての \(t\) に共通の上界とする。
中間レベルの同時不変性により、合成した持ち上げの値は、
各段階の類の代表点 \(g_j\) を渡る
\[
  \Psi(t)\,K_{m-1}(\psi_{t,m-1};g_{m-1},g_m)
        \cdots K_0(\psi_{t,0};g_0,g_1)\phi(g_0)
\]
の有限和である。この表示には各類の安定化群の位数による分母が入らない。
点 \(g_{j+1}\) における出力の格子が次の核の入力の格子になるので、
積に対して各倍率を掛けることができる。
\(\phi\) の局所ノルムが \(1\) 以下であり、\(\Psi(t)\) は \(p\) で整だから、
和に非アルキメデス的三角不等式を使うと上界は \(p^{\sum_{j<m}c_j}\) になる。
レベル作用と可換な線型作用素を各データに施しても同時不変性は保たれる。
最後の \(Y_m\subset X_m\) への制限は分岐行列が整数係数で、
\(p\ne N\) では \(\overline{\pi_N}\) が \(p\) 進単元であるから、
同じ上界を保つ。
\leansrc{New/UpperBound/Divisibility/ChainValuation}{接合した鎖の有限類和と制限の付値}。

\subsection{有限素点での余長}

\(\iota\) を \(p\) の上の付値を定める同型、\(\mathcal O_\iota\) を
その付値環とする。\(M_\iota\) は \(S_0\) の \(\iota\)-整な形式の加群である。
部分加群の余長が \(p^j\) を割るとは、\(M_\iota\) の任意の有限族が
\(\mathcal O_\iota\) 上の正方行列によりその部分加群に移り、その行列式が
\(p^j\) を割ることをいう。とくに行列式の付値は \(j\) 以下である。
この意味での余長を使うので、加群全体の基底を選ぶ必要はない。

\(\mathcal U\) を \(x_1\) 未満の非分解素数の集合とし、
\[
 \ell_p=\frac{fa(\log b)^2}{p}
      \begin{cases} m & (p\equiv\pm1\bmod m),\\1 & (\text{その他}) \end{cases},
 \qquad
 \varepsilon_p=
      \begin{cases}
       2\ell_p &(p\in\mathcal U),\\
       2mfb &(p=N),\\
       0 &(\text{その他})
      \end{cases}
\]
と置く。\(N\ge x_1\) である。各 \(p,\iota\) に対して、
\(M_\iota\cap S\) に余長が \(p^{j_\iota}\) を割る部分加群 \(M'_\iota\) があり
\begin{equation}\label{eq:upper-local}
 j_\iota\le\varepsilon_p\dim S_0,\qquad
 v_p(B_{\xi,\Psi}\phi)\ge G_p-\varepsilon_p
  \quad(\phi\in M'_\iota,\ (p,\mathrm NQ(\xi))=1)
\end{equation}
となることを示す。形式に対する付値の下界 \(v\) は、対応する格子座標
で測るすべての値の \(p\) 進ノルムが \(p^{-v}\) 以下であることをいう。
\(\Psi\) は式 \eqref{eq:upper-seesaw} に課した条件を満たす任意のデータで
あり、\(M'_\iota\) は \(\xi,\Psi\) に先立って選ばれる。

\(p\notin\mathcal U\cup\{N\}\) では \(G_p=0\) である。\(d<x_1\) から
\(p\ne d\) なので、標準格子に
対する各段階の核と最終段階の制限は整性を保つ。
\(p\in\mathcal U\) のとき、まず印なしで
\eqref{eq:upper-one-step} を各段階に適用する。
核の変換則により各段階の付値の下界は加算でき、
合成持ち上げの係数は整なので、その有限和で下界を保つ。
従って \(M_\iota\) 上の下界は \(G_p-\ell_p-mC_p(\varnothing)\) である。
\(\tau,(-p^e)\tau\in\Phi_a\) を満たすすべての \(e\) が \(e\ge e_0\) なら、
等比級数から \(C_p(\varnothing)\le 2fa/p^{e_0}\) である。
\(p\equiv\pm1\bmod m\) の場合は \(e_0=1\) でよい。
それ以外では、重なりがあれば \(p\) は惰性であり、
\(e\) は \(o_0\) の奇数倍である。ここで \(o_0\) は
\(-p\) の場所表示を \(q_0\) を法として読んだ元の位数である。
\(o_0>\lceil\log_p m\rceil\) なら \(p^{e_0}\ge pm\) と取れる。
いずれの場合も \(mC_p(\varnothing)\le\ell_p\) となり
\(M'_\iota=M_\iota\cap S\) でよい。

残る例外は、\(p\) が惰性、\(p\not\equiv\pm1\bmod m\) で
\(o_0\) が奇数かつ \(o_0\le\lceil\log_p m\rceil\) の場合である。
例外素数での印付けの補題により、\(\Phi_a\) が誘導する \(E\) の
有限素点とそれらの共役の各対から片方を取り、集合 \(\Lambda\) を作る。\leanref{exists_isFullMark}
\(\Lambda\) を誘導する無限素点の集合を \(A\) とすれば \(C_p(A)=0\) となる。

\(r_p=\lceil\log_p m\rceil\)、\(\nu_p=r_pf(n_m+r_p)\) と置く。
\(w\in\Lambda\) に対し \(a_w=a\,|\{\tau\in\Phi_a:\tau\text{ が }w
\text{ を誘導する}\}|\) とする。
\(Y\) は、各 \(w\in\Lambda\) で指数 \(r_p\) のヘッケ作用素を
\(p^{a_w}\) 倍して得た作用素 \(T_w\) の積である。
\emph{隣接格子における印付きの評価。}
各 \(w\in\Lambda\) で、指数が \(r_p\) 以下のヘッケ二重剰余類の元を
右から作用させた核を考える。両側の整な行列を除けば、
その \(w\) での成分は行の集合 \(J_w\)（\(|J_w|\le r_p\)）上で \(p\)、
それ以外で \(1\) の対角行列である。共役素点での成分は逆行列になる。
重み埋め込みは共役素点の成分を読むので、行因子の絶対値は
\(p^{\eta_{\tau,b}}\) となる。ここで \(\eta_{\tau,b}\) は、
\(\tau\) が \(w\in\Lambda\) を誘導し \(b\in J_w\) なら \(1\)、
\(\overline w\) を誘導し \(b\in J_w\) なら \(-1\)、
それ以外では \(0\) である。
行列式因子に現れる
\(\sum_\tau e'_\tau\sum_b\eta_{\tau,b}\) は、\(w,\overline w\) 上の
末項重みの和が等しいことにより零である。

式 \eqref{eq:upper-jet} で使う \(w_{\tau h b}\) は、\(\tau\) が \(w\) を誘導し
\(h=1,\ b\notin J_w\) の場合と、\(\tau\) が \(\overline w\) を誘導し
\(b\in J_w\) の場合に \(1\)、他では \(0\) とする。
\(p\) は惰性であり、関係する \(e\) は奇数なので、
\(\tau_e\) が誘導する素点は \(\tau\) が誘導する素点の共役である。
例外の対で、\(b\in J_w\) のときは共役側に、\(b\notin J_w\) のときは
印のある高さ \(1\) の列に \(p\) の因子が付く。従って
式 \eqref{eq:upper-direction-pairing} から除いた項にも必要な因子が現れる。

一箱あたりの指数 \(\eta_{\tau,b}+w_{\tau h b}\) は、
\(\tau\) が \(\overline w\) を誘導するときは \(0\) である。
\(w\) を誘導する場所では、高さ \(1\) の列の箱は \(1\) を寄与する。
他の列では標準充填の行添字が列内で相異なるので、
\(J_w\) に属する箱は一列につき高々 \(r_p\) 個である。
従って全箱について足した指数の上界は
\(\sum_{w\in\Lambda}a_w+r_p\sum_{\tau\in A} r_\tau\) となり、
この量から \(D_{p,j}(A)\) を引いたものが、隣接格子での核の損失を抑える。
\leansrc{New/UpperBound/Divisibility/ValuationAtNonSplitPrimes}{隣接格子の重みと方向の冪指数の計数}。

\emph{ヘッケ作用素を各段のデータに移す。}
\(C_{\boldsymbol x,j}\) を、平行移動の組 \(\boldsymbol x\) に対する
最初の \(j\) 段の合成とする。\(w\in\Lambda\) での剰余体の位数を
\(Q_w=z_w^2\), \(z_w=p^{h_w}\), \(2h_w\le f\) と書き、
\(X_j\) での指数 \(i\) のヘッケ作用素を \(T_{w,i}^{(j)}\) とする。
階数差 \(n_{j+1}-n_j\) を \(t_{\rm rk}+1\) と書く。
この素点では捻りの指標が不分岐で、局所データは真空である。
ヘッケ多項式
\[
 {\cal H}_{w,j}(U)=\sum_{i=0}^{n_j}(-U)^i Q_w^{\binom{i}{2}}T_{w,i}^{(j)}
\]
に対する不分岐テータの恒等式は、一段の持ち上げを \(L_j\) とすれば
\begin{equation}\label{eq:upper-hecke-polynomial}
 {\cal H}_{w,j+1}(U)L_j
    =R_j(U)L_j{\cal H}_{w,j}(z_w^{t_{\rm rk}+1}U)
\end{equation}
である。ここで \(\chi_{w,j+1}\) は持ち上げ指標の基底変換を \(w\) の
一様化元で評価した値で、
\(R_j(U)=\prod_{0\le k\le t_{\rm rk}}(1-Q_w^k z_w^{n_j}\chi_{w,j+1}U)\) である。
\(\chi_{w,j+1}\) は \(w,\overline w\) 上で指標の重みの和が等しいため
\(p\) 進単元であり、\(R_j\in\mathcal O[U]\) で \(R_j(0)=1\) である。
\leansrc{Lib/Weil/Unramified/HeckeTransfer}{不分岐テータにおけるヘッケ多項式の恒等式}。

式 \eqref{eq:upper-hecke-polynomial} に \(R_j^{-1}\in\mathcal O[[U]]\)
を掛けて \(U^i\) の係数を比較する。
\(z_w^{(t_{\rm rk}+1)i}Q_w^{\binom{i}{2}} L_j T_{w,i}^{(j)}\) は、
\(i'\le i\) にわたる \(T_{w,i'}^{(j+1)}L_j\) の整係数線型結合となる。
この前置因子は \(p^{h_w i(t_{\rm rk}+i)}\) であるから、
\(i\le r_p\le t_{\rm rk}\) と \(2h_w\le f\) より
\(p^{f r_p t_{\rm rk}}\) を割る。
よって各段でこの冪を払うと、指数の上限 \(r_p\) を保ったまま
整係数で移送できる。\(j+1\) 段について繰り返し、
\((j+1)f r_p t_{\rm rk}\le \nu_p\) を使って不足する \(p\) の冪を掛けると、
ある \(p\) 進整な係数 \(c_{w,i}\) によって
\begin{equation}\label{eq:upper-hecke-transfer}
 \sum_{i\le r_p}c_{w,i}\,T_{w,i}^{(j+1)}C_{\boldsymbol x,j+1}
       = p^{\nu_p} C_{\boldsymbol x,j+1} T_{w,r_p}^{(0)}
\end{equation}
となる。この式は \(w\) で球的な形式上で成り立ち、
係数は \(\boldsymbol x\) によらない。
\leansrc{New/UpperBound/Divisibility/ValuationAtNonSplitPrimes}{一段の移送から合成の移送へ}。

式 \eqref{eq:upper-hecke-transfer} の左の作用素をデータに作用させ、
\(p^{a_w-\nu_p}\) を掛ける。同じ操作を相異なる \(w\) ごとに施す。
終域でのヘッケ作用をデータに移せるのは持ち上げの共変性による。
データが各 \(w\) で球的なので、終域レベルとの可換性が保たれ、
始域の作用とも可換である。得られたデータによる一段の持ち上げを
\(L'_{\boldsymbol x,j}\) とすると
\(L'_{\boldsymbol x,j}C_{\boldsymbol x,j}
 =C_{\boldsymbol x,j+1}Y\) である。
帰納的に \(m\) 段の変更後の合成は \(C_{\boldsymbol x,m}Y^m\) となる。

各作用素 \(T_{w,i}^{(j+1)}\) は二重剰余類の片側剰余類代表による
平行移動の有限和である。これらはすべて、上で評価した
隣接格子の条件を満たしている。整な \(c_{w,i}\) 付きの有限和は
非アルキメデス的三角不等式により損失を増やさない。
\(p^{a_w-\nu_p}\) を掛けると \(w\) ごとに \(\nu_p-a_w\) の損失が加わる。
従って一段の総損失は
\[
  |\Lambda|\nu_p+r_p\sum_{\tau\in A}r_\tau-D_{p,j}(A)
\]
以下である。可換性により接合したデータの同時不変性が保たれるので、
この上界を全段で加算できる。\leanref{gaugeLe_scaledHeckeOperator_pow}

\leansrc{New/UpperBound/Divisibility/MarkedPrimes}{印を付けた列での損失の評価}の
数値評価により
\[
 |\Lambda|\nu_p+r_p\sum_{\tau\in A}r_\tau
 \le \frac{\ell_p}{m}+|\Phi_a\cap A|\frac a{x_1}.
\]
例外素数では \(v_p(\pi_p)=1\) なので、余分の \(a/x_1\) が消え、
各段階での損失は \((2\ell_p-G_p)/m\) 以下である。
一方、\(j_p=\lceil 2mfa/p\rceil\) とすれば、
\(p^{j_p}M_\iota\) の像の付値は \(G_p-\ell_p\) 以上である。
これは印なしの評価と \(mC_p(\varnothing)\le 2mfa/p\le j_p\) から従う。
従って \(Y^m M_\iota+p^{j_p}M_\iota\) 上で
\(G_p-2\ell_p\) の下界を得る。

この加群の余長には、\leansrc{New/UpperBound/Divisibility/MarkOperator}
{正規化ヘッケ作用素の単因子評価}を適用する。
この評価によれば、\(T_w\) は \(M_\iota\) を保ち、
\(T_wM_\iota+p^{j_p}M_\iota\) の余長は、ある \(p^{c_w}\) を割る。ただし
\[
  c_w\le\left(\frac{12a_w}{mp}+\frac{j_p}{m^2f}\right)\dim S_0 .
\]
この評価に至る計算を記す。\(w,\overline w\) を誘導する \(\Phi_a\)
の素点は \(-p^{o_0}\) の作用で一つの軌道をなす。
その長さを \(2u_w\)、\(P_p=p^{o_0}\) とする。
この軌道の付値環への埋め込みを、共役側の埋め込みに複素共役を合成して
並べる。隣接する埋め込みが \(P_p\) 乗と合同になり、偶数番が \(w\) を
誘導する。従って \(a_w=au_w\) である。
\leansrc{New/UpperBound/ElementaryDivisors/ElementaryDivisorTheorem}
{Frobenius サイクルの単因子定理}では、格子類を結ぶヘッケ作用を
剰余類ごとの制限とペアリングによる行列積に分解し、
各行列の像の余長を比較する。作用の座標表示および
\[
 \begin{gathered}
  0<r_p\le k/2,\qquad a^4\le P_p^{2u_w},\qquad \theta_0=(m^2f)^{-1},\\
  (2u_w k\log(P_p^{2u_w})/\theta_0)^5 P_p^{3r_p}\le a
 \end{gathered}
\]
のもとで、使われる一成分あたりの上界は \(3C_{\mathrm{div}}+\theta_0 j_p\) である。
ここで \(C_{\mathrm{div}}=2u_wa(P_p^{-2r_p}+r_p^2/(kP_p))\)。
上の尺度条件とヘッケ作用の座標表示は
\leansrc{New/UpperBound/Divisibility/MarkOperator}{単因子定理の適用}
で確認される。\(m\le p^{r_p}\)、\(p\le P_p\) より
\(P_p^{-2r_p}\le1/(mp)\) である。また \(mr_p^2\le k\) と
\(p\le P_p\) から第2項も \(1/(mp)\) 以下である。
そこで \(3C_{\mathrm{div}}\le12a_w/(mp)\) を得る。
二つの作用素の積について余長指数が加算できることは、次の行列操作で
分かる。有限族をそれぞれ \(T_1M_\iota+p^{j_p}M_\iota\) に移した後、
各像の \(T_1\) による原像に \(T_2\) の余長行列を適用する。
\(T_1\) は \(p^{j_p}M_\iota\) を保つので、合成行列は
\(T_1T_2M_\iota+p^{j_p}M_\iota\) に移す。その行列式は二つの行列式の積
である。この手順を繰り返すと、\(\sum_w a_w\le af\)、
\(|\Lambda|\le f\) と \(j_p\le 2mfa/p+1\) から、
\(Y^mM_\iota+p^{j_p}M_\iota\) の余長指数は
\[
  m\sum_{w\in\Lambda}
      \left(\frac{12a_w}{mp}+\frac{j_p}{m^2f}\right)\dim S_0
    \le \frac{14fa+p/m}{p}\dim S_0
    \le \ell_p\dim S_0
\]
となる。\(S\) との共通部分を取っても、最初の有限族を \(S\) に取れば
行列操作後も \(S\) に属するので、同じ余長上界が得られる。

最後に \(p=N\) では、\(\mathrm{Up}^{m/2}\phi\) が変更後の標準格子で整に
なる部分加群を取る。
\leansrc{New/UpperBound/Divisibility/ColengthAtLevel}{レベル素数での余長評価}
により、その余長指数は \(mfb\dim S_0\) 以下である。
\leansrc{New/UpperBound/Divisibility/ValuationAtLevel}{レベル素数での付値評価}
では、分岐条件を用いて接合した二段の持ち上げを
\(\mathrm{Up}\) と直接の持ち上げの合成で表す。この部分加群上の
合成持ち上げの付値は \(-2mfb\) 以上である。
\(S\) との共通部分を取っても余長上界は保たれる。
\(G_N=0\) なので、これで \eqref{eq:upper-local} のすべての場合を得た。

局所誤差の合計は
\begin{equation}\label{eq:upper-error-sum}
  \sum_p \varepsilon_p\log p\le E_3.
\end{equation}
実際、\(p\in\mathcal U\) について \(\log p\le\log x_1\) とする。
\(p\equiv\pm1\bmod m\) なら \(p=mc\pm1\), \(1\le c\le p\) と書き
\(m/p\le2/c\) を使う。それ以外は \(c=p\) とする。
\(p=N\) の項にも \(c=1\) が使え、
\(2mfb\log N\le4fa(\log b)^2\log x_1\) である。
各 \(c\) に対応するのは \(N,mc-1,c,mc+1\) の高々四項だから、総和は
\(16fa(\log b)^2\log x_1(1+\log x_1)\) 以下である。
許容データのスケールの不等式により、この量は \(E_3\) 以下である。\leanref{sum_allowedLoss_mul_log_le}

\subsection{無限素点でのノルムと正規化}

上側の鎖の一段で、無限素点 \(\tau\) における階数 \(n_j\) までの
各行長を \(l_{j,\tau,i}\) とし、
\(U_j=\sum_{\tau}\sum_{i<n_j}s_\infty(l_{j,\tau,i})\) と置く。
周期を \(\Omega\) とすると \(s_\infty(q)=\log(q!)+q\log(M/\Omega)\)
である。これが各核のノルムに現れることを示す。
Weyl 加群の Gram ノルムを用い、シュール行列の作用素を双対ノルム
\(1\) 以下の行ベクトルで評価する。一つの多項式行列 \(Z_\tau\) に対し、
その特異値の二乗を大きい順に \(\rho_{\tau,i}\ge0\) とすると、シュール行列の
作用素ノルムの二乗は \(\prod_{\tau,i}\rho_{\tau,i}^{l_{j,\tau,i}}\) 以下である。
これは左右のユニタリ行列が Gram ノルムを保ち、対角行列では
最高重みの行長が指数になることによる。\(\sum_i\rho_{\tau,i}\) は
\(Z_\tau\) の成分の絶対値二乗和である。
\leansrc{Lib/GLn/SchurModule/SchurNorm}{特異値によるシュール行列のノルム評価}。

\(\lambda=a+1\), \(u=\log(M/\Omega)\) と置く。
\(v=u+1/\lambda\) とすれば、指数級数の \(l\) 次の項を使って
\(\rho^l\le l!\exp(lv+\rho e^{-v})\) となる。
\(\sum_{\tau,i}\rho_{\tau,i}\le e^u(2\lambda+2)R\) のとき、
\(e^{1/\lambda}\ge1+1/\lambda\) により
\(e^{-v}\sum_{\tau,i}\rho_{\tau,i}\le2\lambda R\) である。
従ってシュール多項式の成長を抑える倍率は
\[
 \exp\!\left(\frac{U_j}{2}
       +\frac1{2\lambda}\sum_{\tau,i}l_{j,\tau,i}+\lambda R\right)
\]
となる。ここで共役後のガウス積分の変数に掛かるスカラーの二乗絶対値は
\(M/\Omega\) である。

微分の共役を、埋め込みの置換を施したシュール多項式と
共役後の真空の平行移動とのガウス積分で表す。
\(s_j=f n_j n_{j+1}\) 個の複素変数 \(y\) について
\(R=\sum|y|^2/(2\lambda+2)\) を用いると、多項式の
\(\exp(\lambda R)\) と真空の \(\exp(-\sum|y|^2/2)\) の積は
\(\exp(-R)\) となる。各複素座標の測度はルベーグ測度を \(\pi\) で
割ったものなので、このガウスの積分は \((2\lambda+2)^{s_j}\) である。
\(K\) を固定する共役で積分表示を適用すればよい。
他の共役は複素共役を合成して絶対値を保つものに取り替えられる。
\leansrc{New/UpperBound/Kernels/JetConjugatesAndBounds}{ガウス積分と行ごとの階乗評価}。

共役後の整な真空の評価から来る対数定数は
\(\mathrm{const}(P)(q_1 n_j n_{j+1}+1)^2\) である。
これと \(s_j\log(2\lambda+2)\)、および
\((2\lambda)^{-1}\sum_{\tau,i}l_{j,\tau,i}\) の和は、
行長が \(2a\) 以下であることとデータの尺度条件により \(x^{50}\) 以下である。
従って任意の共役について一段の持ち上げでノルムを増やす倍率は
\(\exp(x^{50}+U_j/2)\) 以下である。\leanref{sqrt_weylNormSq_jetMap_le}

行長の主部は、第一行の \(a\,1_{\Phi_a}(\tau)\) と、
続く \(k-1\) 行の \(\kappa(\tau)\) である。
各余分の箱は \(25\log b\) 以下を加える。実際、すべての行長は
\(2a\) 以下であり、一箱を足した際の増分は
\(\log(2a+1)+\log(M/\Omega)\le25\log b\) である。
各段階での余分な箱の総数は
\(f(b+a/x_1+n_m r^{16})\) 以下である。
ここでは \(\Phi_b\) 以外の重さ \(a/x_1\) と、持ち上げの重さの間隔の半分
\(r^8\) が現れる。行の長さの評価を全段について足すと、\leanref{sum_rowSizeSum_le}
\begin{equation}\label{eq:upper-rows}
 \sum_{0\le j<m}U_j
 \le m\left\{|\Phi_a|s_\infty(a)+k\sum_\tau s_\infty(\kappa(\tau))\right\}
         +\frac{E_3}{2}.
\end{equation}

自明な捻りでの接合テンソルの展開において、各段階の類の個数は
\(\exp(x^{30})\) 以下であり、平行移動の係数の共役の絶対値和は
\(\exp(x^{60})\) 以下である。\(\phi\) のすべての値の共役の二乗ノルムが
\(\exp(B)\) 以下なら、\(m\) 個の核の積の評価を足し上げた二乗ノルムは
最終段階の制限前に
\[
  \exp\!\left(B+\sum_j U_j
             +2(mx^{30}+x^{60}+mx^{50})\right)
\]
以下である。\(Y_m\subset X_m\) への制限は、
\leansrc{New/UpperBound/Divisibility/ChainValuation}{最終段の制限の評価}
により評価する。制限写像の分岐行列は整数係数で、二乗ノルムを
高々 \(\exp(x^{70})\) 倍する。さらに \(X_m,Y_m\) の重みの差を
各 \(\tau\) で \(u_\tau\) とすると、制限写像には
\(\prod_\tau \tau(\overline{\pi_N})^{-u_\tau}\) が掛かる。
\(\pi_N\overline{\pi_N}=N\) より、この因子はどの共役でも二乗絶対値
\(N^{-\sum_\tau u_\tau}\) を持つ。
従って指数から少なくとも
\((|\Phi_a|a-fm\,r^8)\log N\) が差し引かれる。

\(s_{\mathrm{main}}\) を \(s_{\mathrm{fin}}(a,\tau)\) の比較に用いる主項、
\(\mathcal R=\sum_{\tau\in\Phi_a}(s_{\mathrm{fin}}(a,\tau)-s_{\mathrm{main}})\) とする。
ここで使うレベル \(N\) の条件は、\(\overline{\log\kappa}
=f^{-1}\sum_\tau\log\kappa(\tau)\), \(D_E=\operatorname{rd}(E)\) と書いて
\[
 \log N\ge m\!\left(\overline{\log\kappa}
                 +\log \frac{M D_E}{2\pi}\right)
          -\log\frac{M D_E}{2\pi}-\tfrac{15}2\log b.
\]
周期の公式
\(\log(M/\Omega)=1+c_1+\log(M D_E/(2\pi))\) と、
\(s_{\mathrm{main}}=a(\log a-\overline{\log\kappa}+c_1)\)、
Stirling の上界
\(\log(a!)\le a\log a-a+1+\frac12\log a\) を併せると、
\[
 s_\infty(a)-s_{\mathrm{main}}
 \le a\!\left(\overline{\log\kappa}+\log \frac{M D_E}{2\pi}\right)
        +1+\tfrac12\log a .
\]
従って \(m(s_\infty(a)-s_{\mathrm{main}})-a\log N\le32a\log b\)
である。ここでは
\(m(1+\frac12\log a)\le a\log b\) と
\(\log(M D_E/(2\pi))\le23\log b\) を用いた。
式 \eqref{eq:upper-gain-sum} を代入すれば
\[
 H_0-m\mathcal R+2\sum_p G_p\log p
 =m|\Phi_a|s_{\mathrm{main}}+mk\sum_\tau s_\infty(\kappa(\tau)).
\]
この等式と \eqref{eq:upper-rows} を使い、
\(32fa\log b+x^{70}+fm\,r^8\log N+
  2(mx^{30}+x^{60}+mx^{50})\) を \(E_3/2\) に納めると、
\(B_{1,\Psi}\phi\) のすべての値の共役の二乗ノルムは
\begin{equation}\label{eq:upper-arch}
  \exp\!\left(B+H_0-m\mathcal R+2\sum_p G_p\log p+E_3\right)
\end{equation}
以下となる。\leanref{weylNormSq_restrictedChain_le}

\paragraph{両方の評価に用いる有理数}
\(g_p=\lfloor G_p-\varepsilon_p\rfloor\) とし
\[
 C_B=\prod_{p\in\mathcal U\cup\{N\}}p^{-g_p},\qquad
 \mathfrak p_g\phi
   = C_B|\overline G|e_0\rho(g)c_{\mathrm{ss}}^{-1}\mathcal P\phi .
\]
\(C_B\) は正の有理数である。\eqref{eq:upper-local} により
\(C_B B_{\xi,\Psi_\xi}\phi\) は \(M'_\iota\) 上で整であり、
\eqref{eq:upper-seesaw} と \eqref{eq:upper-project} から
\(\Theta_\xi\mathfrak p_g\phi\) も整になる。すべての許容捻りに
同じ \(M'_\iota\) が使える。

一方、丸めによる損失は \(\sum_{p\in\mathcal U\cup\{N\}}\log p\le b\)
であり、\eqref{eq:upper-error-sum} によって
\begin{equation}\label{eq:upper-cb}
 \log C_B\le-\sum_p G_p\log p+E_3+b .
\end{equation}
二乗ノルムに倍率 \(C_B^2\) を掛けると、有限素点の利得と合わせた
指数への寄与は
\[
 2\log C_B+2\sum_p G_p\log p
    =2\sum_{p\in\mathcal U\cup\{N\}}(G_p-g_p)\log p
    \le 2(E_3+b)
\]
である。従って式 \eqref{eq:upper-arch} の素点の主項が相殺され、
射影の損失 \(x^{99}\) とこの誤差が残る。

\begin{proof}[定理~\ref{thm:upper-bound}の証明]
\(S_0\) は全有限素点で整で、任意の共役の値の二乗ノルムが
\(\exp(E_3)/4^f\) 以下の形式で張られる。\leanref{span_isShort_eq_top}
この評価では各格子類の安定化群が自明なので、格子類ごとに
その類に支持を持つ形式を取る。格子類を \(c\) として、
\(c\) の格子と重みの行列式部分を表す格子に小さい基底を選ぶ。
Weyl 加群で1 の冪根を代入した評価ベクトルは空間を張り、
それらに選んだ基底の行列を作用させても生成性を保つ。
局所では基底が格子内にあり、1 の冪根の絶対値は \(1\) なので、
得られる形式はすべての有限素点で整である。
無限素点での小さい基底の定数は \(C=6(q_1d)(fk)^2\) である。
高さ \(k\) 未満の列の箱は全場所で \(fa\) 個以下であり、
同補題の二乗ノルムの評価は \(C^{6(fk+fa)}\) を与える。
許容データにより \(4^f C^{6(fk+fa)}\le \exp(E_3)\) である。

\(N\) の上の指定素点 \(w\) でのスカラー作用の生成元を \(z_w\)
とすると、\(\prod_w(1-z_w)\) の像が \(S\) である。
スカラー平行移動は整性を保ち、一つの差 \(1-z_w\) は
二乗ノルムの上界を高々 \(4\) 倍する。指定素点は \(f\) 個なので、
これらの像は \(S\) を張り、\(B=E_3\) が使える。
\eqref{eq:upper-arch} の指数、射影の \(x^{99}\)、および
\eqref{eq:upper-cb} の二倍を合わせると
\[
 \|\Theta_1\mathfrak p_g\phi\|_\sigma^2
 \le \exp\bigl(H_0-m\mathcal R+x^{99}+4E_3+2b\bigr).
\]
許容データの \(x^{99}+4E_3+2b\le E_4\) から (1) が従う。

(2) について、まず \(g_i=g\) を固定する。各 \(\iota\) で
\((\phi_i)_i\) に \eqref{eq:upper-local} の余長行列を施すと、
その各成分が \(M'_\iota\) に入る。同じ行列が
\((\mathfrak p_g\phi_i)_i\) をすべての許容捻りについて整な族に移す。
この行列式の付値は \(\varepsilon_p\dim S_0\) 以下である。
分岐条件で失う次元は
\(\frac{f}{N-1}\dim S_0\) 以下であり、\(\frac{f}{N-1}\le 1/2\) であるから
\(\dim S_0\le2\dim S\) となる。式 \eqref{eq:upper-error-sum} により
固定した \(g\) に対する対数指数は \(2(\dim S)E_3\) 以下である。
一般の族は \(g_i\) ごとの部分族に分け、それぞれの補正行列を
対角ブロックに並べる。その行列式は各ブロックの行列式の積となり、
対数指数を足すと \(2|\mathcal T|(\dim S)E_3\) を得る。
\end{proof}

\section{ポアンカレ写像による生成}\label{sec:spanning}

\begin{theorem}[生成定理 B]\label{thm:spanning}
目標空間 \(V\) の次元 \(h\) は正である。さらに、有限集合 \(T\subset\bar G\) で
\[
  |T|\dim S\leq 5h\log(1+2b^2),\qquad
  \dim\sum_{g\in T}\mathfrak p_g S\geq (1-4/b^2)h
\]
を満たすものがある。
\end{theorem}

ポアンカレ級数が張る空間を調べるには、有限レベル群による平行移動の間の内積を評価する。
以下では、まず全始域 \(S_0\) を使い、最後に分岐条件を課す。
\(Y_0\) 上の主合同レベルの形式全体を \(\mathcal H\) とし、\(\bar G\) のユニタリ作用を
\(\rho\colon\bar G\to U(\mathcal H)\) とする。\(V=e_0\mathcal H\) への直交射影
\(e_0\) はこの作用と可換である。
\(D_X=\dim S_0>0\) と置く。絶対値による作用素の評価には、底本の双剰余類と
Bessel 座標で \(S_0\) に入る内積を用いる。以下、作用素のノルムを Hilbert--Schmidt
ノルム \(\|\cdot\|_{\mathrm{HS}}\) とする。

\subsection{展開式と二次モーメントの局所化}

無限素点で選んだ Bessel 系の形式次数の積を \(\deg_B\)、階数 \(r\) の標準エルミート
格子の種の質量を \(\mathrm{mass}_r\) とする。ここで種の質量は、格子の同型類を
自己同型群の位数の逆数で重み付けした和である。目標空間について
\begin{equation}\label{eq:span-dimension-mass}
 h\leq |\bar G|\deg_B\,\mathrm{mass}_r
\end{equation}
が成り立つ。実際、底本の多重度一致により \(h\) は定符号空間 \(W_m\) 上の
対応する射影像の次元に等しい。\(W_m\) 上では主合同レベルにおける類の
自己同型群が自明なので、類の個数は最大コンパクトレベルでの質量
\(\mathrm{mass}_r\) にレベルの指数 \(|\bar G|\) を掛けたものになる。
各類に付随する代数表現の次元の積は、Weyl の次元公式により \(\deg_B\) に等しい。
したがって、射影前の空間の次元は \eqref{eq:span-dimension-mass} の右辺である。\leanref{targetDim_eq_finrank_goodForms_spaceW}\leanref{finrank_principalForms_spaceW}

\(J\) を主合同レベル、\(\nu\) を類によらない実ユニポテント格子の余体積とする。
ポアンカレ写像を底本の座標で \((\operatorname{vol}(J)\nu)^{-1/2}\) 倍したものを
\(P\colon S_0\to\mathcal H\) と書き、
\[
 P_s=P^*\rho(s)P \qquad(s\in\bar G)
\]
と置く。\(c,c'\) を \(U(X_0)\) 上の始域の類、\(h_c,h_{c'}\) をその代表の \(U(Y_0)\) での像、
\(W\) を Bessel 関数、\(\beta_c,\beta_{c'}\) を類の切断関数とする。
\(P_s\) の \((c,c')\) ブロックは、次の行列の成分ごとの複素共役になる。
\begin{equation}\label{eq:span-unfolded}
 \frac1\nu
 \sum_{\substack{\gamma\in U(Y_0)(F)\\
       h_c^{-1}\gamma h_{c'}\in K,\
       \operatorname{red}(h_c^{-1}\gamma h_{c'})=s}}
 \int
    \beta_c(y)\beta_{c'}(\gamma^{-1}y)\,
    W(y)W(\gamma^{-1}y)^*\,dy .
\end{equation}
ここで \(F\) は \(E\) の最大総実部分体、\(K\) は有限アデールでの格子安定化群であり、
積分は無限素点の積にわたる。この式はポアンカレ級数の内積を展開して得られる。\leanref{hasAmplitudeBlocks_scaledPoincare}

有理点の Bessel 部分群を \(R_0\)、その整点の還元像を \(\bar R_0\subset\bar G\) とする。
\eqref{eq:span-unfolded} を \(\gamma\in R_0\) に限ると、\(\bar R_0\) 上の
ユニタリ表現 \(r_\psi\) を \(\bar G\) に零で延長した作用素が得られる。
\[
  E_s=P_s-r_\psi(s),\qquad
  D_{\mathrm{id}}=|\bar R_0|D_X,\qquad
  D_{\mathrm{main}}=\frac{D_X^2}{\deg_B\,\mathrm{mass}_r}.
\]
\(\bar G=G_0\times G_1\) と書く。\(G_i\) は \(p_i\) の上で選んだ素点での局所因子の積で、
\(\bar G\) の部分群としても扱う。\(R_0\) 上の項の計算では
\begin{equation}\label{eq:span-equivariance}
 P_t=P_1r_\psi(t),\qquad E_{ts}=r_\psi(t)E_s
 \quad(t\in\bar R_0,\ s\in\bar G)
\end{equation}
を得る。これは \(r_\psi(t)\) が始域で作用し、展開式の項を移し替えることによる。\leanref{besselTranslate_mul_left}

広い範囲の二次モーメントを評価する際には、\eqref{eq:span-unfolded} の項を、
Bessel フレームの等方部分 \(I\) に付随する対応ごとにまとめる。
項 \(\gamma\) の対応は、\(I\) と \(\gamma I\) の両方に直交する \(y\) に対し、
\(y\) と \(\gamma^{-1}y\) の \(I^\perp/I\) での像を組にした部分空間である。
以下、\(\mathcal L(\Omega)\) を、二つの類、還元 \(s\in\Omega\)、この対応を固定して
切断した Bessel 関数の振幅を合計し、その二乗ノルムを足した値を \(\nu^2\) で割った
ものとする。この和には、フレーム積の行列式 \(\Delta\) とその一つ手前の首座小行列式
\(\Delta'\) に対し
\(\mathfrak d(\gamma)
   =\Delta\overline{\Delta}-N^2\Delta'\overline{\Delta'}\ne0\)
の項を入れる。\leanref{linkSecondMoment}

以下では \(E\) に原始 \(5040\) 乗根と底本が選ぶ Weil 表現の分裂指標があることを用いる。
許容データから作る Bessel 系とレベルに対し、底本は三つの評価を与える。
底本が構成する Bessel 系は、以下の評価に必要な大きさとレベルの仮定を満たす。\leanref{filling}

\begin{enumerate}
\item
\(\Omega\) が非自明な \(1\) の冪根による \(\bar R_0\) の剰余類と交わらなければ、
\begin{equation}\label{eq:span-cutoff}
 \Bigl(\sum_{s\in\Omega}\|E_s\|_{\mathrm{HS}}^2\Bigr)^{1/2}
 \leq \mathcal L(\Omega)^{1/2}+b^{-15}D_{\mathrm{id}}^{1/2}.
\end{equation}
\ref{sec:span-pair-count}節で、振幅の二乗和の交差項と切断の誤差を評価する。\leanref{sqrt_secondMoment_compl_besselS_le}\leanref{secondMoment_remainingTerms_le}
\item
選んだ素点とその共役のうち、還元したフレームの張る空間が自己の \(s\) による像と
非零の共通ベクトルをもつところのノルムの積を \(\operatorname{Meet}(s)\) とする。
正整数 \(z\) が \(\Omega\) 上で \(\operatorname{Meet}(s)\) 以下なら
\begin{equation}\label{eq:span-intersect}
 \mathcal L(\Omega)\leq
       \left(\frac{b^{2f}}z\right)^k D_{\mathrm{main}} .
\end{equation}
\ref{sec:span-pair-count}節で、共通部分の格子を使って同じ対応をもつ項の組を数える。\leanref{linkSecondMoment_le_of_le_meetNorm}
\item
\(p_0\) の上の選んだ素点の集合 \(A\) で、各 \(w\in A\) について
フレームの張る空間の非零ベクトルを \(w\) またはその共役で指定し、
\(s\) がこれらを固定するとする。そのような \(s\) の部分群を \(\Gamma_A\) と書く。
\(p_0\) の上の選んだ素点のうち少なくとも一つで、フレーム積の正の位数の真の首座
小行列式がすべて非零になる \(s\) を開な還元と呼ぶ。
\(\Omega_A\subset\Gamma_A\) を開な還元の集合とすると
\begin{equation}\label{eq:span-open-input}
 \mathcal L(\Omega_A)\leq (1+b^{-4})b^{-6k|A|}D_{\mathrm{main}}.
\end{equation}
\ref{sec:span-open-proof}節で、Siegel の評価と Poisson 和公式を用い、
局所質量の評価から固定条件一つ当たりの因子 \(b^{-6k}\) を導く。\leanref{linkSecondMoment_isOpenElement_le}
\end{enumerate}

\subsection{共通格子による計数}\label{sec:span-pair-count}

\(n=\varrho+1\)、\(n_X=\#\{c\}\)、\(\tau=b^{kf/4}\)、
\(n_p=|\bar R_0|n_X/N^{kf}\) と置く。切断 Bessel 関数の項を
整ユニポテント群の二重軌道ごとにまとめる。ある無限素点の積の点 \(y\) について
\(y\) と \(\gamma^{-1}y\) がともに、簡約条件と Bessel 長の上界を満たす領域に入るとき、
項 \(\gamma\) を近接項と呼ぶ。各軌道には近接項があれば、その中から代表元を一つ選ぶ。
還元 \(s\) が左右のユニポテント指標を一致させる条件は、積が
\(us=su'\) となる整ユニポテント元に対して両側の指標の値が等しいことである。
条件が破れると、指標の比を掛ける変数変換で振幅が消える。

以下で数える項は、\(\mathfrak d(\gamma)\ne0\) の項と、
\(\mathfrak d(\gamma)=0\) で \(\gamma I\ne I\) となり有理ユニポテント指標が
一致する項である。有理点での一致条件は、\(u\gamma=\gamma u'\) を満たす
有理ユニポテント元の指標値が等しいことである。この項の集合は有理ユニポテント群の
左右作用で安定である。
\leansrc{New/Spanning/Unfolding/TermStructure}{残りの項と有理ユニポテント指標}

近接項 \(\gamma\) の格子 \(L_\gamma\) は、類 \(c\) の格子内の \(X_0\) のベクトルのうち、
\(M\) 倍の整等方ベクトルを足して \(\gamma\) のフレームに直交させられるものからなる。
同じ \((c,s,c')\) 上の代表元 \(\gamma,\gamma'\) に対し、
二つの対応の共通部分で両座標の最後のフレーム方向への成分を零とする。
この部分を第一座標で \(X_0\) に移した空間に \(L_\gamma\cap L_{\gamma'}\) を制限した格子を
\(L_{12}\) とする。その余階数を \(l\)、判別式を \(d\) と書く。判別式は格子基底の
Gram 行列式の一般化指数の平方根である。各対応は次元 \(k+1\) で、その共通部分の
余次元を \(i\) とすると、最後の座標を落とす条件は二つなので \(l\le i+1\) となる。
\(i=0\) で片方の \(\mathfrak d\) が零なら両方が零で、このとき \(l=0\) である。
近接項に対する長さと格子の評価は \(l\le n^3\) と \(1\le d\le b^{n^3 f}\) を与える。
\leansrc{New/Spanning/RemainingTerms/Gain}{共通格子の階数と判別式の上界}

細胞の階数 \(j\) をもつ項の量 \(G_\gamma\) を、\(N^f\) 倍した首座小行列式の
正規化ノルムと \(\mathfrak d(\gamma)\) のノルムの最大値に
\(M^{2\min(\varrho,j)f}\) を掛け、左右の等方行格子の指数の積で割ったものとする。
\(1\le z\le\operatorname{Meet}(s)\) に対し、\(i=0\) で \(\mathfrak d(\gamma)\ne0\) なら
\(\zeta=z\) とし、ほかでは \(\zeta=1\) とする。二つの格子の切り分けから
\begin{equation}\label{eq:span-pair-gain}
 (d\zeta)^2\le (2/M)^{2if}G_\gamma G_{\gamma'}
\end{equation}
が従う。ここで \(\zeta=z\) のときの因子は次の行列式から得られる。
フレームが交わる剰余体で、共通ベクトルを二通りに表示し、係数が同時には零でない
最大の添字を取る。指標の一致条件により、その添字で現れた係数の対は、次の添字でも
二つの旗の共通ベクトルの係数として現れる。これは、次の添字より前の両方の列を零にする
線形形式が、次の列にその二係数を掛けたものに同じ値を取ることから従う。
この操作を繰り返すと、最後の二つの係数が同時には零でない関係式を得る。
\leansrc{New/Spanning/Unfolding/Climbing}{共通ベクトルの係数を次の段へ移す補題}
この関係式を、先行部分に直交する二つの整リフトとの積で写すと、最後の二列に核が生じる。
行列式を \(\delta\) とすれば、交わるすべての素点で \(\delta\) は零である。
従って \(\widehat{\mathcal O}_E/\delta\widehat{\mathcal O}_E\) からそれらの
剰余体の積への全射があり、\(\delta\) の一般化指数は \(\operatorname{Meet}(s)\) 以上となる。
\leansrc{New/Spanning/RemainingTerms/Cut}{交わる素点での切断行列式}

具体的には、対応の共通部分を両方の最後の座標が消える部分とその補空間に分け、
さらに各対応で共通部分の直交補空間を取る。最後の座標が消える部分における
\(L_\gamma\) の判別式を \(d_\gamma^0\)、等方行から生じるねじれ指数を
\(a_\gamma\)、共通部分の外の補格子の判別式の二乗を \(\Delta_\gamma\) とする。
切断行列式の因子と Gram 行列式の積公式により
\((a_\gamma d_\gamma^0\zeta)^2\Delta_\gamma\le G_\gamma^2\) となる。
二つの補空間の間のペアリングは非退化で、還元が同じなので、Hermite 形式の左右成分の差からなるペアリングを
\(M\) で割ると整になる。その行列式を Gram 行列式で評価すると
\(1\le(2/M)^{4if}\Delta_\gamma\Delta_{\gamma'}\) を得る。
また \(L_\gamma\) の共通部分に \(L_{\gamma'}\) の条件を課すと増える指数は
\(a_{\gamma'}\) で抑えられ、逆向きにも同じことが成り立つ。
従って \(d\le d_\gamma^0a_{\gamma'}\) と \(d\le d_{\gamma'}^0a_\gamma\) を掛けた後、
補格子の判別式を消去すると \eqref{eq:span-pair-gain} が得られる。
\leansrc{New/Spanning/RemainingTerms/Gain}{二つの対応の判別式比較}

余階数 \(l\) を固定する。\(\bar R_0\) の一つのラベルを持つ項は、
類 \(c\) のレベル \(M\) のユニポテント作用までのフレームで決まる。
\(L\subset L_\gamma\) のフレームを数えるため、まず細胞の階数 \(j\) と
フレーム積の行列 \(B\) を固定する。\(B\) の最初の \(j\) 行を取り、
残る行を恒等行列の対応する行で補った行列を \(B_j\) とする。
その行列式 \(\delta_j\) は \(B\) の位数 \(j\) の首座小行列式で、非零である。
\([a]\) は整有限イデール \(a\) の指数
\(|\widehat{\mathcal O}_E/a\widehat{\mathcal O}_E|\) を表す。

格子 \(L\) の基底を \(Y=yA_L\) と書く。\(y\) の \(k-l\) 本の列は有理で
一次独立、\(A_L\) は有限アデール上で可逆である。\(L\subset L_\gamma\) の条件は、
\(X_0\) 成分の行列 \(X_\gamma\) と \(B\) の等方行 \(B_{\mathrm{iso}}\) に対して
\[
 X_\gamma^*y+M B_{\mathrm{iso}}^*w=0,\qquad
 wA_L\text{ は整}
\]
となる持ち上げ \(w\) を与える。ここでは \({}^*\) は随伴行列を表す。
\(wA_L\) の各行を \(Y^*Y\) の行格子で分類する。その指数は
\(\operatorname{disc}(L)^2\) なので \(\varrho\) 行の分類は
\(\operatorname{disc}(L)^{2\varrho}\) 通り以下である。同じ類では
ユニポテント作用で上の直交条件の右辺を共通にできる。

従って \(X_\gamma\) の差の \(n\) 本の列は、\(L\) に直交する
階数 \(l\) の格子に入る。これを法 \(M\delta_j\) で分類すると
\([M\delta_j]^{ln}\) 通りであり、同じ類の \(X_0\) 成分は
ユニポテント作用で一致させられる。残る等方成分については、
\(j\) 段以降の共通ベクトルを記録する随伴小行列を用いた整行列で
まず分類する。近接条件はその成分を各無限素点で \(b^{182n^2}\) 以下に
抑えるので、その行列の数は \(b^{366n^4 f}\) 以下である。さらに等方成分の
\(\varrho n\) 個の成分を法 \(M\det(B_j B_j^*)\) で分類する。
同じ類に入るフレームはユニポテント作用で一致する。
\leansrc{New/Spanning/RemainingTerms/FrameParts}{持ち上げの直交条件とユニポテント作用}
\leansrc{New/Spanning/RemainingTerms/FramesOverLattice}{固定格子を含むフレームの合同類}

固定した \(j,B\) の上では、これらの類の個数から
\[
 \operatorname{disc}(L)^{2\varrho}\,[M\delta_j]^{ln}
       b^{366n^4 f}[M\det(B_j B_j^*)]^{n\varrho}
\]
を得る。二つの法の指数と \(\operatorname{disc}(L)^2\) はいずれも
\(b^{2n^3f}\) 以下で、\(j\) の選択は \(b^{2f}\) 通り、\(B\) は
\(b^{270n^3f}\) 通り以下なので、\(L\) を含む近接代表元は
\(b^{2000n^6(l+1)f}\) 個以下である。ラベルは左右の指標が一致する
還元に対して選べ、同じ組の二つの項に共通である。従って二つの
代表元の選び方は \(b^{4000n^6(l+1)f}\le\tau\) 通り以下になる。

格子の \(d^{-k}\) の和は、ランク \(k-l\) の格子について各類を動かして平均する。
極大レベルの部分群に関する剰余類は少なくとも \((N^k/b)^f\) 個ある。
一つのユニモジュラー格子を固定し、\(L\) の飽和格子の基底を
環 \(\widehat{\mathcal O}_E\) 上の全体の基底に延長する。
\(L\) の基底が飽和格子の基底の右から整行列 \(A_{\mathrm{sat}}\) を
掛けたものなら、指定行の最初の \(k-l\) 成分 \(a\) は
\(aA_{\mathrm{sat}}=\pi_N w\) を満たす必要がある。
\(w\) を \(A_{\mathrm{sat}}\) の行格子を法として選ぶと、
\(a\) の法 \(\pi_N\) の値が決まるので、選択は
\([\det A_{\mathrm{sat}}]\) 通り以下である。あとの \(l\) 成分に
\(\bigl|\widehat{\mathcal O}_E/\pi_N\widehat{\mathcal O}_E\bigr|^l=N^{lf}\)
が掛かる。剰余類は行の法 \(\pi_N\) の値で区別されるため、
この格子で寄与する剰余類は高々 \([\det A_{\mathrm{sat}}]N^{lf}\) 個になる。
\(d\le D_\ast=b^{n^3 f}\) の格子の重みを合わせた剰余類上の和は
\[
  N^{lf}D_\ast^{2\varrho+6}
       (200fk\,\operatorname{rd}_E)^{f(k+1)l}2^k
  \le b^{-f}\tau^{l+1}
\]
で抑えられる。ここで \(\operatorname{rd}_E\) は \(E\) の根判別式である。
剰余類数で割ると、類を一つ当たりに平均した格子の重みは
\(N^{-kf}\tau^{l+1}\) 以下となる。全 \(n_X\) 類、\(|\bar R_0|\) 個のラベルと
一つの格子当たりの組の数を掛けると
\begin{equation}\label{eq:span-pair-count}
 \sum_{\substack{\text{近接代表元の組}\\ \operatorname{corank} L_{12}=l}} d^{-k}
 \le \tau^{l+2}n_p .
\end{equation}
\leansrc{New/Spanning/RemainingTerms/LatticesOfCorank}{剰余類上の格子の平均}
\leansrc{New/Spanning/RemainingTerms/Count}{共通格子で重み付けした計数}

一つの近接代表元の寄与には、軌道和に Schur の評価を適用した上界
\[
 \sigma_\gamma=\tau^2\left(\frac{D_j}{n_pG_\gamma^k}\right)^{1/2},
 \qquad
 D_j=\begin{cases}
       D_{\mathrm{main}} & j=\varrho+1,\\
       \tau^{-16}D_{\mathrm{id}} & j\le\varrho
     \end{cases}
\]
を用いる。すなわち、上記の残りの項のうち、左右の有理ユニポテント作用で安定な集合に対して、
切断振幅の
ノルムを \(\nu\) で割った値は、選んだ近接代表元上の \(\sum_\gamma \sigma_\gamma\)
以下である。階数 \(j=\varrho+1\) は \(\mathfrak d(\gamma)\ne0\) の場合に当たる。
\leansrc{New/Spanning/RemainingTerms/Size}{二重軌道和に対する振幅評価}
\((d\eta)^2\le G_\gamma G_{\gamma'}\) で \(D\ge D_j,D_{j'}\) なら、
\[
 \sigma_\gamma\sigma_{\gamma'}
 \le \frac{\tau^4D}{n_p\eta^k}\,d^{-k}.
\]
従って、各余階数 \(l\) で \(\tau^{l+6}D/\eta^k\le C\) なら、
\eqref{eq:span-pair-count} によりその重みの和は \(C\) 以下である。
余階数は \(n^3+1\le\tau\) 通りなので、総和は \(\tau C\) 以下となる。

同じ対応で判別式が非零の場合、\(l\le1\) であり、
\eqref{eq:span-pair-gain} により \(\eta=z\)、\(D=D_{\mathrm{main}}\) と取れる。
上の計算で \(C=\tau^7D_{\mathrm{main}}/z^k\) を使うと
\(\tau^8 D_{\mathrm{main}}/z^k=(b^{2f}/z)^kD_{\mathrm{main}}\) が得られる。
各組の還元と対応は一意に決まるので、これが \eqref{eq:span-intersect} を与える。

同じ対応で判別式が零なら、\(l=0\)、\(\eta=1\)、
\(D=\tau^{-16}D_{\mathrm{id}}\) と取れて、\(C=\tau^{-10}D_{\mathrm{id}}\) で足りる。
対応が異なる場合には \(i\ge1\)、\(l\le i+1\) で、
\(\eta=(M/2)^{if}\)、\(D=D_{\mathrm{main}}\) とする。
\(\rho_\ast=\tau(2/M)^{kf}\le1\) および
\(D_{\mathrm{main}}\le(M/b^6)^{kf}D_{\mathrm{id}}\) から
\[
 \frac{\tau^{l+6}D_{\mathrm{main}}}{\eta^k}
 \le \tau^7\rho_\ast^iD_{\mathrm{main}}
 \le \tau^8(2/b^6)^{kf}D_{\mathrm{id}}
 \le \tau^{-10}D_{\mathrm{id}} .
\]
最後の不等式には \(2\le b\) と \(\tau=b^{kf/4}\) を使った。
この二つの場合を合わせて同じ \(C=\tau^{-10}D_{\mathrm{id}}\) を使えるため、
総和は \(\tau^{-9}D_{\mathrm{id}}\le b^{-kf}D_{\mathrm{id}}\) となる。
\leansrc{New/Spanning/RemainingTerms/Sum}{対応ごとの二乗和と交差項の評価}

この計数を振幅に適用するには、ブロックを判別式が非零の各対応と
判別式が零の全項とに分け、部分ごとのノルムの和を二乗する。
非零判別式の同一対応の対角項をそのまま足すと \(\mathcal L(\Omega)\) である。
それ以外には上の \(\sigma_\gamma\sigma_{\gamma'}\) を適用する。
各組は一つの還元に属するので、ブロック全体でも同じ上界を使える。

切断に伴う誤差について、各 \((c,s)\) 行の振幅差のノルム和の上界を
\(q=\nu/(2b^{k^3 f})\) とする Bessel 関数の裾の評価を用いる。
行の二乗和はノルム和の二乗以下なので \(q^2\) 以下であり、
\(|\bar G|b^{40}\le(2b^{k^3 f})^2\) と \(n_X\le D_{\mathrm{id}}\) により、
全 \((c,s)\) 行の和を \(\nu^2\) で割った値は \(b^{-40}D_{\mathrm{id}}\) 以下になる。
この評価は切断の前後で同じ項の集合に適用する。
\leansrc{New/Spanning/RemainingTerms/Size}{切断による振幅差の行ごとの評価}
許された \(\Omega\) 上で \(E_s\) の全振幅は、残りの項の和に等しい。
実際、Bessel 部分群の外で残りの項にも入らないもののうち、有理指標の一致条件を
満たさない項は、左右作用で安定な集合上で振幅の和が消える。一致条件を満たす項は
最初の段からノルム一の係数対をもち、その還元は除外したスカラー剰余類に属する。
\leansrc{New/Spanning/Unfolding/SquareExpansion}{Bessel 部分群の外の振幅の整理}
二乗和の平方根の三角不等式と \(kf\ge40\) から
\[
 \Bigl(\sum_{s\in\Omega}\|E_s\|_{\mathrm{HS}}^2\Bigr)^{1/2}
 \le \mathcal L(\Omega)^{1/2}
       +(b^{-kf/2}+b^{-20})D_{\mathrm{id}}^{1/2}
 \le \mathcal L(\Omega)^{1/2}+b^{-15}D_{\mathrm{id}}^{1/2}.
\]
これが \eqref{eq:span-cutoff} である。

\subsection{開な還元の二次モーメント}\label{sec:span-open-proof}

\paragraph{固定条件による局所質量の減少}
有限素点の共役で安定な集合 \(\mathcal T\) を取り、法 \(M\) を支える素点を含める。
その局所積で整ユニタリー群を \(K_{\mathcal T}\)、局所レベル \(g\) の主合同部分群を
\(K_{\mathcal T}(g)\) と書く。Haar 測度を \(\mu\) とし、
\(K_{\mathcal T}(g)\subset Z\subset K_{\mathcal T}\) に対して
\[
 \mathfrak M(Z)=\int
    \left|\int_Z \chi_g(yxy^{-1})\,d\mu(y)\right|^2d\mu(x),
 \qquad
 m_Z=\frac{\mu(Z)\mu(K_{\mathcal T}(g))}{u_g^2}
\]
と置く。\(\chi_g\) はユニポテント指標を主合同部分群との積に延長し、その外で零と
した関数であり、\(u_g\) はこの延長に用いるユニポテントの測度定数である。
\(x\in K_{\mathcal T}(g)\) では \(\chi_g(yxy^{-1})=1\) であり、
この部分の積分は \(\mu(Z)^2\mu(K_{\mathcal T}(g))\) となる。

局所格子データについて、整数環に属する可逆元のうち逆元が整でないものの
モジュラスが \(Q^{-1}\) 以下であり、\(Q\ge2\) とする。合同部分群の外での
積分を分母イデアルの指数の範囲ごとに分ける局所評価により、
\(a+10^8n^4\le k\) なら零段の積分の優関数の積分は
\(m_Z2^{10^9n^4}\min(1,Q^{-a})\) 以下である。この評価では、合同部分群の外の元 \(x\) について、
どのフレーム座標にも深さ \(Q\) 以上のものがなければ、
安定なフレーム積に整単元を選べることを使う。その場合 \((x-1)/g\) は整となり、
仮定に反する。従って深いベクトルによる利得があり、その二乗を用いる
積分評価の座標数の損失と、分母イデアルの二進区分を足し合わせる損失は
それぞれ \(2^{99\cdot13n}\) と \(2^{4\cdot10^8n^4}\) で抑えられる。
その積は \(2^{10^9n^4}\) 以下である。
\(a=t+10^9n^4\) を選ぶと \(2/Q\le1\) によって係数が吸収される。
積分の優関数に \(u_g^2\) を掛け、合同部分群の寄与を足すと、
\(t+2\cdot10^9n^4\le k\) の下で
\begin{equation}\label{eq:span-local-mass}
 \mathfrak M(Z)
 \le\mu(Z)\mu(K_{\mathcal T}(g))\,\bigl(\mu(Z)+Q^{-t}\bigr).
\end{equation}
\leansrc{New/Spanning/Frames/FiniteShells}{分母イデアルによる局所積分の評価}

一つの \(w\in A\) とその共役に限って \(K_{\mathcal T}\) を \(K_w\) と書き、
固定条件の逆像を \(Z_w\) とする。還元から \(\mathrm{GL}_r(\mathbf F_{p_0})\)
への写像は全射で、共役側では逆随伴を通じて作用する。いずれも非零ベクトルへの
推移性から \([K_w:Z_w]=p_0^r-1\) となる。\(B=b^9\) と置くと
\(2\le B\le p_0\) であり、選んだ素点の局所積では非単元部分の加法体積が
\(B^{-1}\) 以下である。実際、各体積は剰余体位数の逆数、素点の個数は \(b\) 以下、
各剰余体位数は \(b^{10}\) 以上である。この上界は逆元が整でない元の
モジュラスにも適用できる。
許容条件から \(t_0=k-10^{10}n^4\) は
\[
 t_0+1+n+2\cdot10^9n^4\le k,\qquad
 6k\le9t_0,\qquad t_0+1<r
\]
を満たす。指数の公式は \(\mu(Z_w)\le B^{-t_0-1}\mu(K_w)\) を与える。
また局所積分の測度の正規化では \(\mu(K_w)\ge2^{-n}\ge B^{-n}\) である。
\eqref{eq:span-local-mass} に \(Q=B,\ t=t_0+1+n\) を代入すると
\[
 \mathfrak M(Z_w)\le
 \mu(Z_w)\mu(K_w(g))
       \bigl(\mu(Z_w)+B^{-t_0-1}\mu(K_w)\bigr).
\]
これに \(\mu(K_w)b^{6k}\) を掛け、\(b^{6k}\le B^{t_0}\) と \(2/B\le1\) を使う。
\(K_w(g)\) 上の積分だけで
\(\mathfrak M(K_w)\ge\mu(K_w)^2\mu(K_w(g))\) となるから
\[
 \mathfrak M(Z_w)\mu(K_w)b^{6k}
 \le\mathfrak M(K_w)\mu(Z_w).
\]
相異なる共役対では群と指標が直積に分かれる。従って質量積分と測度を
それぞれ掛け、残る素点では条件なしの群を使うと、同じ不等式で
\(K_w,Z_w,b^{6k}\) を \(K_{\mathcal T},Z,b^{6k|A|}\) に置き換えられる。
この比の比較は直積測度からの共通の Haar 測度の定数倍に依らない。
\leansrc{New/Spanning/Frames/FixedLines}{固定ベクトル一つ当たりの局所質量比}

\(Z=K_{\mathcal T}\) には \eqref{eq:span-local-mass} を \(Q=2,\ t=n\) で適用する。
\(\mu(K_{\mathcal T})\le1+2^{-n}\) および \(2^{-\varrho}\le b^{-6}\) を用いると
\(\mathfrak M(K_{\mathcal T})\le(1+b^{-6})\mu(K_{\mathcal T})\mu(K_{\mathcal T}(g))\)
である。零段の軌道形式は指標を平均した相関の二乗であり、その未正規化の値の絶対値に
\(u_g^2\) を掛けると \(\mathfrak M(Z)\) 以下となる。そこで \(m_Z\) で
正規化した零段の有限素点の核は
\begin{equation}\label{eq:span-local-zero}
 m_A=(1+b^{-6})b^{-6k|A|}
\end{equation}
以下に抑えられる。\leansrc{New/Spanning/MainTerm/FiniteSieveSystem}{有限素点の零段の密度}

\paragraph{Siegel の評価と最後の座標}
\(\Omega_A\) 上の非零の振幅ブロック \((c,s,c',\Lambda)\) を有限個取る。
真の首座小行列式は有理点でもすべて非零であるため、大細胞の分解によって、
各ブロックを最後の積 \(t\in E^\times\)、先行する双曲基底ベクトルをすべて
固定する代表元 \(\sigma\)、および整持ち上げ \(n_s\) で記述できる。ここで
\(|\Phi_v(t)|^2>N^2\) はすべての無限素点 \(v\) で成り立ち、
\((c,c',\operatorname{red}(n_s),\sigma)\) はブロックを定める。
\(\sigma\) を固定すると残る首座の座標は有限集合 \(T_0\subset E^\times\)
に入るので、ブロックの振幅は
\(\mathcal A=\sum_{\mathbf p\in T_0^\varrho} a(\mathbf p)\) と書ける。
\(a(\mathbf p)\) は Weyl 代表 \(w(\mathbf p,t)\sigma\) の切断されたユニポテント
二重剰余類の振幅である。\(\mathcal T\) を拡大して \(T_0\) の元とその逆元、
\(M^{-1},N^{-1}\)、格子の基底とその逆、および類の代表元で移した各 \(\sigma\) が
\(\mathcal T\) の外で整になるようにする。
\leansrc{New/Spanning/MainTerm/EnergyEstimate}{大細胞の振幅和と有限集合への帰着}

無限素点 \(v\) での第 \(j\) 座標の倍率を
\(\lambda_{v,j}=\kappa_v^{2(\varrho-j)}/N\) とする。ここで
\(\kappa_v=M\Phi_v(\varpi)/(2\pi\mathrm i)\)、\(\varpi\) は異なるイデアルの生成元である。
\(R_t(\mathbf p,\mathbf q)\) を、最後の座標を \(t\) とする無限素点の
最上段の核の積とする。各核は軌道積分を座標の倍率と頭部の密度で調整した
Gram 核である。有限素点の最後の座標の密度 \(F_t(\mathbf p,\mathbf q)\) は、
同一の最後のベクトルをもつファイバーでの Gram ペアリングを二次曲面上で
積分したものである。

\(t\) ごとに、\(\sigma\) が最後のフレームベクトルに作用してできる \(X_0\) 成分 \(y\)
を取ると \(N\langle y,y\rangle=t\bar t-N^2\) である。
\[
 G_t=\left(\prod_v(|\Phi_v(t)|^2-N^2)\right)^{k-1}>0
\]
とする。各 \(v\) の Bessel 座標の個数を \(d_{B,v}\)、その積を \(d_B\) とする。
大細胞の体積定数を \(c_G=\operatorname{vol}(U(k))^2/
(\operatorname{vol}(U(r))\operatorname{vol}(U(k-1)))\) とすると、
無限素点の軌道積分に掛ける重みは
\[
 \rho_t(\mathbf p)=
 \prod_v \frac{\sqrt{c_G}}{d_{B,v}}
 \left(\frac{|\Phi_v(t)|^2}{N^2}-1\right)^{(k-1)/2}
 \prod_{j<\varrho}|\lambda_{v,j}\Phi_v(p_j)|^{k-1+2(\varrho-j)} .
\]
対応する有限素点の頭部密度を \(h_t(\mathbf p)\) とする。各 \(p_j\) は
\(\mathcal T\) の外で単元なので、積公式により
\(\prod_v|\Phi_v(p_j)|\) は \(\mathcal T\) における \(p_j\) のモジュラスの
平方根の逆数となる。よって \(\rho_t(\mathbf p)h_t(\mathbf p)=r_0\sqrt{G_t}\) である。
ここで
\[
 r_0=\frac{1}{d_B}\prod_v
       \left(\sqrt{c_G}\prod_{j<n}\lambda_{v,j}^{k-1+2(\varrho-j)}\right).
\]
さらに \(D=|\bar R_0|M^{f\varrho(2k+2\varrho+1)}\) とし、
\(C_S=D\,n_Xc_S\) を有限素点の正規化も含む Siegel の定数とする。
同じ点 \(x=(c,y)\) のブロックを数える。
\(\Psi_t(x)\) は、類の代表元の逆で移した \(y\) から作る最後の局所ベクトルでの値とする。
\(\mathcal K=K_{\mathcal T}(g)\)、\(\mathcal C=Z/\mathcal K\) と置く。
各頭部 \(\mathbf p\) の尺度付き軌道積分は、有限商 \(\mathcal C^2\) に
座標をもつ。位置 \(\ell\) と座標 \((a,a')\in\mathcal C^2\) を固定し、
その座標を頭部ごとに並べたベクトルを \(v_{\ell,a,a'}\) とすると、
\(R_t\) を掛けた局所ペアリングの和は
\[
 Q_t(\ell)=\sum_{a,a'\in\mathcal C} v_{\ell,a,a'}^*R_t v_{\ell,a,a'}\ge0
\]
となる。類 \(c,c'\) と \(\sigma\) を固定すると、還元 \(s\) をもつ
ブロックは、\(a'\) が \(n_s^{-1}a\) の合同類になる座標を占める。
異なる \(s\) では同じ \((a,a')\) を占めることがない。
各座標の項は軌道積分を展開すると
\(\mu(\mathcal K)^2G_t r_0^2\|\mathcal A\|^2\) となる。
\(a\) について足すと \(|\mathcal C|\) 倍となり、すべてのブロックの和を
\(Q_t(\ell)\) で抑えられる。

\(H_{\mathrm{fix}}\) を最後のフレームベクトルを固定する局所群とする。
同じ点 \(x\) に属する \((c',\sigma)\) の位置 \(\ell\) は、
\(H_{\mathrm{fix}}\) の左作用で移り合い、その合同類は相異なる。
同じ合同類なら局所では合同であり、\(\mathcal T\) の外で整であることと
合わせると、始域の類も有理代表元も一致するからである。
\(Q_t\) は \(H_{\mathrm{fix}}\cap\mathcal K\) の左作用で不変である。
従ってその剰余類の測度を掛けて和を取ると、ファイバーでの積分以下になる。
\(\mu_H\) を \(H_{\mathrm{fix}}\) の Haar 測度としたとき、
\(D\mu_H(H_{\mathrm{fix}}\cap\mathcal K)\mu(\mathcal K)^2|\mathcal C|\ge m_Z\)
であり、有限素点のファイバーのペアリングを延長した非負の階段関数
\(\Psi_t\) に対して
\[
 G_t(m_Z/D)\sum_{(c,y)=x} r_0^2\|\mathcal A\|^2\le\Psi_t(x)
\]
を得る。
\leansrc{New/Spanning/Frames/TopFloor}{有限商の座標とファイバーのペアリング}

\(\beta_t=(t\bar t-N^2)/N\) とする。階数 \(k\ge3\) の正定値空間に対する
一ベクトルの Siegel の公式を、\(\mathcal T\) の外で整であり、\(\mathcal T\) では
\(\Psi_t\) を最後のベクトルに引き戻した試験関数に適用する。
これはコンパクト台をもち、局所定数で、始域のレベル群に不変である。
\(\beta_t\) は総正であり、対応する類の和は
\[
 n_X|\operatorname{disc}(E)|^{1/4}
       \prod_v\frac{\pi^k(\Phi_v(\beta_t))^{k-1}}{(k-1)!}
\]
と有限アデールのノルム密度との積になる。
\leansrc{Lib/SiegelWeil/OneVector/SiegelClassSum}{一つのノルムに対する Siegel の公式}

ノルム密度は、有限ベクトル \(w\) と自己共役な有限アデール周波数 \(u\)
にわたって、試験関数と \(\psi_f(u(w^*w-\beta_t)/2)\) の積を積分したものである。
周波数とベクトルを \(\mathcal T\) の成分で分けると、
\(\mathcal T\) での球面密度と、補集合で整であることから出る Gauss 和の
積分との積になる。Gauss 和の被積分関数を
\[
 F_{u,\Gamma}(w)=\psi_f\!\left(u(w^*w-\Gamma)/2\right),
 \qquad w\in\widehat{\mathcal O}_E^{\,k}
\]
と書く。\(\Gamma\) はノルム値から来る定数で、次の評価は \(\Gamma\) に関して一様である。
\(\int F_{u,\Gamma}(w)\,dw\) の絶対値の二乗は、
\(\overline{F_{u,\Gamma}(w)}F_{u,\Gamma}(w+h)\) の \(w,h\) に関する積分に等しい。
\(w\) の第 \(j\) 座標を \(\bar a\) だけ平行移動すると、相関には
\(\psi_f(u h_j a)\) が掛かる。従って \(w\) の積分が残るのは、
すべての \(h_j\) が導手イデアル
\[
 I_u=\{h\in\widehat{\mathcal O}_E:
            \psi_f(uha)=1\text{ for all }a\in\widehat{\mathcal O}_E\}
\]
に入る場合だけである。指標の絶対値が \(1\) なので、この部分の積分は
体積の積で抑えられ、\(q_u=|\widehat{\mathcal O}_E/I_u|\) と書くと
\[
 \left|\int F_{u,\Gamma}(w)\,dw\right|
 \le \operatorname{vol}(\widehat{\mathcal O}_E)^k q_u^{-k/2}.
\]
導手 \(I\) を固定して周波数を数えると、余不同イデアルの自己共役部分による
剰余類は、有限加法群 \(\widehat{\mathcal O}_E/I\) の指標により区別される。
一剰余類の測度は \(1\) で、その個数は高々 \([\widehat{\mathcal O}_E:I]\) 個である。
\(\mu_0=3\varrho+4\)、\(k\ge2(\mu_0+1)\) として
各導手上で積分した上界を足すと、
整アデールの有限指数イデアルのゼータ和
\(\sum_I[\widehat{\mathcal O}_E:I]^{-\mu_0}\)
に \(\operatorname{vol}(\widehat{\mathcal O}_E)^k\) を掛けたものになる。
この和は \([E:\mathbf Q]\le 2^{2\varrho}\) で Dedekind ゼータの
Euler 積と比較して \(1+2^{-\varrho}\le1+b^{-6}\) 以下である。
\leansrc{Lib/SiegelWeil/Mass}{相関の指標直交性と導手ごとの Gauss 和の評価}

\(\mathcal T\) で残ったノルム密度は、有限素点の Fourier 反転によって、
球面 \(\langle y,y\rangle=\beta_t\) 上の試験関数の積分以下である。
最後のベクトルの密度 \(F_t\) はこの積分を \(\pi_N\) のモジュラス \(N^{-f}\) で
割った値である。従って球面の積分を密度に換える際に \(N^{-f}\) を掛け、
無限素点のノルム因子の \(N^{-f(k-1)}\) と合わせると
\[
 c_S=\left(\frac{\pi^k}{(k-1)!N^k}\right)^f
       |\operatorname{disc}(E)|^{1/4}
       \operatorname{vol}(\widehat{\mathcal O}_E)^k
\]
となる。点 \(x\) ごとに一回ずつ和を取ると
\[
 \sum_x\Psi_t(x)\le(1+b^{-6})n_Xc_S G_t\,
             \operatorname{Re}\sum_{\mathbf p,\mathbf q}R_t(\mathbf p,\mathbf q)
                                                       F_t(\mathbf p,\mathbf q).
\]
両辺で \(G_t\) と \(m_Z\) を整理し、\(t\) を動かすと
\begin{equation}\label{eq:span-siegel-sum}
 r_0^2\sum_{\mathrm{blocks}}\|\mathcal A\|^2
 \le (1+b^{-6})C_S\sum_t\operatorname{Re}\sum_{\mathbf p,\mathbf q}
                         R_t(\mathbf p,\mathbf q)\frac{F_t(\mathbf p,\mathbf q)}{m_Z}.
\end{equation}
\leansrc{New/Spanning/MainTerm/SiegelTopBound}{類の和に対する Siegel の球面評価と局所密度}

\paragraph{Poisson 和公式による座標の消去}
無限素点の核と正規化した有限素点の核を各段で掛けたものを \(J_i\) とする。
\(\mathcal S_{i+1}\) を、最初の \(i\) 座標が \(\mathbf p,\mathbf q\in T_0^i\) に入り、
その次の座標を共通の \(t\in T_0\) とする \(J_{i+1}\) の和の実部とする。
頭部以降の座標は \(1\) で埋め、無限素点では各座標に \(\lambda_{v,j}\) を掛ける。
\(\mathcal S_0=\operatorname{Re}J_0\) は空の頭部に対する値である。
共通の座標ごとの和は Gram 核の正値性から非負であり、
\eqref{eq:span-siegel-sum} の核の和は \(\mathcal S_n\) 以下である。

第 \(i\) 座標の和に片側の Poisson 和公式を適用する。頭部を固定して合計した
核を非負の下半連続な関数とみなし、\(\mathcal T\) の外では整アデールの指示関数を
掛ける。零周波数は軌道形式の積分公式により \(J_i\) の全頭部の和を与え、
定数は
\[
  q_i=\frac{\operatorname{vol}(\widehat{\mathcal O}_E)}
                 {\prod_v\lambda_{v,i}^2}.
\]
残る周波数は \(\varpi\) の逆数で尺度を変えた \(\mathcal T\)-整な有理点に限られる。
局所 Fourier 評価の積の非零周波数にわたる和を \(\epsilon\)、
長さ \(i\) の頭部の非負の重みの和を \(W_i\) とすると、この寄与は
\(q_i\epsilon W_i^2\) 以下である。零周波数による全頭部の和で、
最後の座標が等しい組を残せば \(\mathcal S_i\) になる。
最後の座標 \(u,v\) が異なる組では、有限素点で
\(M^{-1}(u^{-1}-v^{-1})\) が非零周波数を与える。この点は \(\mathcal T\) の外で整であり、
同じ減衰評価を適用して、これらの組の和は \(\epsilon W_i^2\) 以下となる。
従って
\[
 \mathcal S_{i+1}\le q_i(\mathcal S_i+2\epsilon W_i^2).
\]

周波数側では \(K_0=k-n^7\) と置く。無限素点で用いる減衰関数は
\(d_{K_0}(a)=a^{K_0}\min(1,a^6)\) である。各 \(v\) の減衰因子と
その体積比は
\[
 c_{\mathrm{vol}}b^{2000n^9}d_{K_0}(|\xi_v|^{-1}),
 \qquad
 c_{\mathrm{vol}}=\pi^{(2\varrho+1)k}
            \frac{\operatorname{vol}(U(k))}{\operatorname{vol}(U(r))}.
\]
有限素点では、周波数の一般化指数の \(K_0/2\) 乗を
分母指数の三乗で割り \(b^{n^9 f}\) を掛ける。
非零の \(\mathcal T\)-整数 \(y\) に対応する周波数は
\(\xi_v=(c_v/\lambda_{v,j})\Phi_v(y)\) の形で、
二つの誤差に用いた \(c_v\) はいずれも \(|c_v|^{-1}\le|\kappa_v|/M\) を満たす。
従って
\[
 |c_v/\lambda_{v,j}|^{-1}\le
 \theta_v:=|\kappa_v|^{2\varrho+1}/(MN)
 \le B_F:=b^{23(2\varrho+1)}.
\]
\(H(y)=\operatorname{Norm}(\mathfrak a_y)\prod_v\max(1,|\Phi_v(y)|)^2\)
を \(y\) の高さとする。\(\mathfrak a_y\) は分母イデアルである。
\(\mathcal T\) の外で \(y\) は整なので、有限素点の一般化指数の平方根は
積公式により \(\prod_v|\Phi_v(y)|\) 以下である。この上界を掛けると、
無限素点で \(|\Phi_v(y)|^{-K_0}\) が消える。
\(\min(1,|\xi_v|^{-1})\le B_F/\max(1,|\Phi_v(y)|)\) と
分母指数が \(\operatorname{Norm}(\mathfrak a_y)\) 以上であることから、
\(b\) の冪と体積比の定数を除いた一つの周波数の寄与は
\(H(y)^{-3}\prod_v(\theta_v^{K_0}B_F^6)\) 以下である。
高さの逆三乗の和は \((64\cdot38\exp(4))^f\le b^{3f}\)。
これと \((bB_F^2)^3\le b^{n^9}\) から
\[
 \epsilon\le b^{2002n^9f}\mathcal D_\ast,\qquad
 \mathcal D_\ast=\prod_v c_{\mathrm{vol}}\theta_v^{K_0}.
\]
\leansrc{New/Spanning/MainTerm/GlobalSieve}{積公式と高さによる周波数の総和}

頭部の重みでは、各先頭積について有限素点の一般化指数の \(n^6/2\) 乗を掛け、
無限素点では尺度付き先頭積の絶対値の逆数と \(1\) の大きい方の \(n^6\) 乗を掛ける。
先頭積が整でないか、尺度付き絶対値が \(b^{60n^2}\) を超えるなら重みは零である。
重みが非零の頭部は、各共役で絶対値 \(b^{120n^2}\) 以下の代数的整数からなる
先頭積の列で一意に定まる。その個数は格子計数により
\((1+2b^{120n^2})^{2fi}\) 以下である。積公式によって一つの頭部の重みは
\(b^{120n^8fi}\) 以下なので、\(i\le\varrho\) で \(W_i\le b^{130n^9f}\) を得る。
一方、\(\lambda_{v,j}\le b^{46\varrho}\) と
\(\operatorname{vol}(\widehat{\mathcal O}_E)=(2/\operatorname{rd}_E)^f\) から
\(\prod_{j<i}q_j\ge b^{-100n^2f}\) である。従って
\[
 W_i^2\le b^{260n^9f+100n^2f}\prod_{j<i}q_j
       \le b^{400n^9f}\prod_{j<i}q_j .
\]
\leansrc{New/Spanning/MainTerm/GlobalSieve}{先頭積の計数と頭部の重み}
\(\deg_v\) を無限素点 \(v\) の Bessel 系の形式次数とする。
\(\mathcal S_0\le \prod_v((1+b^{-10})/\deg_v)\,m_A\) と
\(2n\le b^{n^9f}\) から、漸化式を \(n\) 回適用して
\begin{equation}\label{eq:span-sieve-sum}
 \mathcal S_n\le \left(\prod_{j<n}q_j\right)
  \left[
    \left(\prod_v\frac{1+b^{-10}}{\deg_v}\right)m_A
    + b^{2403n^9f}\mathcal D_\ast
  \right].
\end{equation}

\paragraph{定数の比較}
各座標の倍率の \(k+2(\varrho-j)\) 乗の積を \(\Lambda_\ast\) とすると、
\(r_0^2\) の式と Poisson 和公式の定数 \(q_j\) から
\[
 V_\ast=\frac{\prod_{j<n}q_j}{r_0^2}
 =\frac{\operatorname{vol}(\widehat{\mathcal O}_E)^n}{c_G^f}
    \left(\frac{d_B}{\Lambda_\ast}\right)^2
\]
となる。余体積 \(\nu\) の公式と Siegel の定数の比較は
\[
 \begin{split}
 C_S V_\ast
  &\le(1+b^{-9})D_{\mathrm{main}}\deg_B\nu^2,\\
 C_S V_\ast\mathcal D_\ast
  &\le b^{n^9f}d_BD_{\mathrm{id}}\nu^2
 \end{split}
\]
を与える。第一式は類数・体積の式を代入した主項の比較である。第二式では
\(s_d=\sqrt{\operatorname{rd}_E}\)、\(Q_d=Ms_d^2/(2\pi)\)、
\(x_d=Q_d^{2\varrho+1}/(MN)\) と置く。\(\prod_v|\kappa_v|=Q_d^f\) を使い、
\(\mathcal D_\ast\) の \(x_d^{k-n^7}\) のうち \(k\) 乗を Siegel の定数に
吸収すると、\(C_SV_\ast\mathcal D_\ast/(d_BD_{\mathrm{id}}\nu^2)\) は
次の因子の \(f\) 乗となる。
\[
 \frac{(2/s_d^2)^{n+2\varrho^2}(s_d/s_d^{2\varrho})}
      {2M^{\varrho(2\varrho+1)}}\,x_d^{-n^7}.
\]
\(2/s_d^2\le b\)、\(s_d/s_d^{2\varrho}\le1\)、
\(x_d^{-1}\le MN\le b^{23+60n}\) を用いると、この因子は \(b^{n^9}\) 以下である。
\leansrc{New/Spanning/MainTerm/EnergyEstimate}{余体積と減衰を含む定数の比較}

\(Y=b^{-6k|A|}D_{\mathrm{main}}\nu^2\) と置く。
許容条件 \(d_B\le b^{2kf}\)、\(b^{11kf}D_{\mathrm{id}}\le D_{\mathrm{main}}\)、
\(|A|\le f\)、\(512000n^9\le k\) により
\(2404n^9f+2kf+10+6k|A|\le11kf\) である。
従って \eqref{eq:span-sieve-sum} の誤差項に \(C_SV_\ast\) の対応する係数を
掛けた値、すなわち \(C_SV_\ast b^{2403n^9f}\mathcal D_\ast\) は
\(b^{-10}Y\) 以下になる。主項には
\((1+b^{-10})^f\le1+b^{-9}\) と \(\prod_v\deg_v=\deg_B\) を使う。
前者は \(f\le b/1000\) により \((1+b^{-10})^f\le
\exp(b^{-9}/1000)\le1+b^{-9}\) となることから従う。
\eqref{eq:span-siegel-sum}、\eqref{eq:span-local-zero}、
\eqref{eq:span-sieve-sum} を合わせると
\[
 \sum_{\mathrm{blocks}}\|\mathcal A\|^2
 \le (1+b^{-6})\left[(1+b^{-9})^2(1+b^{-6})+b^{-10}\right]Y
 \le (1+b^{-4})Y .
\]
最後の不等式は \(b\ge1000\) による。実際、\(s_*=b^{-4}/100\) とすると
\(b^{-6},b^{-9},b^{-10}\le s_*\le1\) なので、
係数は \((1+s_*)^4+s_*(1+s_*)\le1+17s_*\) 以下である。
有限部分和の上限を取り、\(\nu^2\) で割ると
\eqref{eq:span-open-input} が得られる。

\(\bar R_0G_i\) の元は、もう一方の素数の上で \(\bar R_0\) の元に一致する。
従って、その素数の上の \(f\) 個の素点とその共役のすべてでフレームが交わり、
\(z\) を当該素数の \(2f\) 乗に取れる。また、この集合は
非自明な \(1\) の冪根による \(\bar R_0\) の剰余類と交わらない。
各 \(p_j^2\geq M/b\) とデータの評価
\(D_{\mathrm{main}}\le (M/b^6)^{kf}D_{\mathrm{id}}\) を使うと、
\eqref{eq:span-intersect} の右辺は \(b^{-3kf}D_{\mathrm{id}}\) 以下である。
許容条件 \(kf\geq 40\) と \(b\ge 1000\) の下で \eqref{eq:span-cutoff} は
\begin{equation}\label{eq:span-bessel-near}
 \sum_{s\in\bar R_0G_i}\|E_s\|_{\mathrm{HS}}^2
 \leq b^{-12}D_{\mathrm{id}}
\end{equation}
を与える。\leanref{secondMoment_compl_besselS_le}

\subsection{射影の前後の二次モーメント}

\eqref{eq:span-equivariance} と \eqref{eq:span-bessel-near} を
\(\bar R_0\) に制限すると、\(\|E_t\|_{\mathrm{HS}}=\|P^*P-1\|_{\mathrm{HS}}\)
であることから
\begin{equation}\label{eq:span-gram-isometry}
 \|P^*P-1\|_{\mathrm{HS}}^2\leq b^{-12}D_X.
\end{equation}
部分空間 \(U\subset S_0\) への直交射影を \(Q_U\) とすれば、
\[
 \|P Q_U\|_{\mathrm{HS}}^2
 = \dim U+\operatorname{tr}\bigl(Q_U(P^*P-1)\bigr)
 \geq \dim U-b^{-6}D_X.
\]
目標への射影による損失を評価するため、\(e_0=e_{p_0}e_{p_1}\) と分解し、
\(a_i(t)\) を \(\mathbf C[G_i]\) の補射影 \(1-e_{p_i}\) の係数とする。
底本は、有限一般線形群の多項式で指定した型の数とその次数を使って
\[
 \sum_{t\in G_i}|a_i(t)|^2=a_i(1)\leq \tfrac12 b^{-8}
\]
を得る。等号は群環の自己共役な冪等元を二乗し、単位元の係数を比較して従う。
不等号は、射影で取り除く型の個数と次数の評価による。\leanref{bad_types_rare}

\((1-e_{p_i})\) は直交射影だから
\[
 \|(1-e_{p_i})P\|_{\mathrm{HS}}^2
  =\operatorname{tr}\sum_{t\in G_i} a_i(t)\bigl(r_\psi(t)+E_t\bigr).
\]
Bessel 項のトレースは \(D_X\) 倍の還元ユニポテント指標の値である。
\(a_i\) が類関数なのでこの指標を共役平均 \(\psi_U^\natural\) で置き換えられる。
\(\sum_{s\in\bar G}|\psi_U^\natural(s)|^2\le2\) という底本の評価と
Cauchy--Schwarz の不等式から、その寄与の絶対値は \(b^{-4}D_X\) 以下となる。\leanref{trace_besselTranslate_eq_mul_reducedChar}

残る項には、関数 \(s\mapsto E_s\) を \(\bar R_0G_i\) の外で零として、
\(f(s)\mapsto\sum_t a_i(t)f(st)\) を作用させる。この作用は
\(\ell^2(\bar G)\) 上の直交射影なので、二乗和を増やさない。
\(s\in\bar R_0\) での値は \eqref{eq:span-equivariance} により
\(r_\psi(s)\sum_t a_i(t)E_t\) である。従って \eqref{eq:span-bessel-near} を
\(|\bar R_0|\) で割ると
\(\|\sum_t a_i(t)E_t\|_{\mathrm{HS}}^2\le b^{-12}D_X\) となる。
そのトレースの絶対値は \(b^{-6}D_X\le b^{-4}D_X\) である。
これを二つの素数で適用し、直交射影の分解
\(1-e_{p_0}e_{p_1}=(1-e_{p_0})+e_{p_0}(1-e_{p_1})\) を用いると
\(\|(1-e_0)P\|_{\mathrm{HS}}^2\le4b^{-4}D_X\) を得る。よって
\begin{equation}\label{eq:span-near-isometry}
 \|e_0P Q_U\|_{\mathrm{HS}}^2\geq \dim U-5b^{-4}D_X
       \qquad(U\subset S_0).
\end{equation}
\noindent
この不等式が、射影後の写像の近似等長性を表す。\leanref{isNearIsometry_piTensorElem_comp}
\eqref{eq:span-near-isometry} で \(U=S_0\) とすれば \(h>0\) を得る。

次に、全群上の二次モーメントを求める。\(\beta(s)=\|P_s\|_{\mathrm{HS}}^2\) とし、
\(X_t=\rho(t)PP^*\rho(t)^*\) と置く。跡の巡回性により
\[
 \langle X_t,X_u\rangle_{\mathrm{HS}}=\beta(t^{-1}u).
\]
したがって \(\beta\) は正定値型であり、\(\beta(s)\leq\beta(1)\) である。
\eqref{eq:span-gram-isometry} から
\(\beta(1)\le(1+b^{-6})^2D_X\le2D_X\) も得る。

ベクトルの指定も動かし、有限個の選択の中で
\[
 b^{6k|A|}\sum_{s\in\Gamma_A}\beta(s)
\]
が最大となる \(A\) とベクトルを取る。\(Q=\sum_{s\in\Gamma_A}\beta(s)\) とする。
まず \(A\) が \(p_0\) の上の選んだ素点全体である場合を考える。
\(C=b^{-6kf}D_{\mathrm{main}}\) と置く。整アデールのスカラー倍による
\(\bar R_0\) の剰余類の合併には
高々 \(M^f|\bar R_0|\) 元しかない。そこでの寄与は \(2 M^fD_{\mathrm{id}}\) 以下で、
底本のパラメータ比較 \(8 M^f D_{\mathrm{id}}\le b^{-6kf}D_{\mathrm{main}}\)
により \(C/2\) 以下である。残りの元では \(z=p_0^f\) を取れ、
\eqref{eq:span-intersect} の右辺は、\(p_0\ge b^{10}\) により \(C/8\) 以下になる。
\eqref{eq:span-cutoff} の誤差項の二乗も \(C/8\) 以下なので、
平方根の和の二乗は \(C/2\) 以下である。よって \(Q\le C\) を得る。\leanref{secondMoment_fixer_univ_le}

次に \(w\notin A\) を一つ取れるとする。最大性により、\(w\) またはその共役で
指定した任意の非零フレームベクトルを固定する元に限ると、\(\beta\) の和は
\(b^{-6k}Q\) 以下になる。これを使って \(w\) での退化小行列式を数える。
剰余体の位数を \(p_0\)、その上での \(\Gamma_A\) の行列表現を \(\pi\) とする。
第 \(\ell\) 首座小行列式を
\(\det Z_s\)、\(Z_s=A_1^{\mathsf T}\pi(s)A_0\) と書く。
\(A_0\) の列はフレームの最初の \(\ell\) 本、\(A_1\) の列はそれに対応する線形形式である。
非自明な加法指標 \(\psi\) の直交性により
\[
 p_0 \sum_{\det Z_s=0}\beta(s)
 \le \sum_{s\in\Gamma_A}\beta(s)\#\{y:y^{\mathsf T}Z_s=0\}
 = \sum_{y\in\mathbf F_{p_0}^\ell}\frac{1}{p_0^\ell}
    \sum_{v\in\mathbf F_{p_0}^\ell}\ \sum_{s\in\Gamma_A}
        \psi(y^{\mathsf T}Z_s v)\beta(s).
\]
\(y=0\) と \(v=0\) に由来する寄与の和は \(2Q\) 以下である。
\(y,v\ne0\) の和には、点指示関数の正規直交基底と加法指標の正規直交基底を用いる。
具体的には、\(\pi(s)A_0v\) に対応する点指示関数と、
\(\pi(s)^{-\mathsf T}A_1y\) に対応する正規化指標の関数を、それぞれ \(X_s\) と
テンソル積にして \(s\) で足す。各和のノルムの二乗は
\(|\Gamma_A|\) と対応する固定部分群上の \(\beta\) の和との積である。
双方の内積を展開すると上式の内側の指標和が
\(|\Gamma_A|p_0^{-r/2}\) 倍されて現れる。従って Cauchy--Schwarz の不等式により
指標和の実部は \(p_0^{r/2}b^{-6k}Q\) 以下である。
以上と \(p_0^{\ell+r/2}\le b^{6k}\) から
\[
  \sum_{\det Z_s=0}\beta(s)
  \le p_0^{-1}\bigl(2+p_0^{\ell+r/2}b^{-6k}\bigr)Q
  \le 3b^{-10}Q .
\]
ここで共役素点を許したのは、\(\pi\) の反傾表現で線形形式を固定する条件も得るためである。
フレームの長さは \(b\) 以下なので、小行列式の和を取って
開な還元の補集合の寄与は \(b^{-4}Q\) 以下となる。\leanref{card_mul_sum_det_eq_zero_le}\leanref{sum_not_isOpenElement_le}

開な還元にはスカラー倍の剰余類が含まれず、\(r_\psi(s)=0\) である。
\eqref{eq:span-cutoff} と \eqref{eq:span-open-input} を適用し、
\(D_{\mathrm{id}}\le b^{-6k|A|}D_{\mathrm{main}}\) と
\(b^{-15}\le\tfrac12 b^{-4}\) を使うと、開な還元の寄与の平方根は
\((1+b^{-4})(b^{-6k|A|}D_{\mathrm{main}})^{1/2}\) 以下になる。
両方の寄与を合わせると
\[
 Q \leq (1+4b^{-4})b^{-6k|A|}D_{\mathrm{main}}+b^{-4}Q,
 \qquad
 Q\leq (1+9b^{-4})b^{-6k|A|}D_{\mathrm{main}}.
\]
最後の不等式には \((1-u)(1+9u)\ge1+4u\) を \(u=b^{-4}\) に適用した。
全素点を取った場合にもこの評価が成り立つ。
最大性を \(A=\varnothing\) に対して使えば、結論は
\begin{equation}\label{eq:span-second-moment}
 \sum_{s\in\bar G}\|P_s\|_{\mathrm{HS}}^2
    \leq (1+9b^{-4})D_{\mathrm{main}}
\end{equation}
となる。\leanref{sum_secondMoment_le}

\subsection{二次モーメントから階数へ}

\(\epsilon=(8b^2)^{-1}\) と置く。\eqref{eq:span-dimension-mass} と
\eqref{eq:span-second-moment} から \(\deg_B\,\mathrm{mass}_r\) が消え、
\(9b^{-4}\le\epsilon\) より
\begin{equation}\label{eq:span-moment-comparison}
 \frac{h}{|\bar G|}\sum_{s\in\bar G}\|P_s\|_{\mathrm{HS}}^2
    \leq (1+\epsilon)D_X^2 .
\end{equation}
任意の部分空間 \(U\subset S_0\) について、
\eqref{eq:span-near-isometry} も \(\|e_0P Q_U\|_{\mathrm{HS}}^2\geq\dim U-\epsilon D_X\)
と書ける。

\(F_1\subset V\) を部分空間とし、\(\alpha=\dim F_1/h\) とする。
\(Q_{F_1}\) を \(\mathcal H\) での \(F_1\) への直交射影とすると、二次モーメントから
\begin{align*}
 \frac1{|\bar G|}\sum_{s\in\bar G}\|Q_{F_1}\rho(s)P\|_{\mathrm{HS}}^2
 &\leq
 \left(\dim F_1\cdot
        \frac1{|\bar G|}\sum_{s\in\bar G}\|P_s\|_{\mathrm{HS}}^2\right)^{1/2}\\
 &\leq \sqrt{\alpha(1+\epsilon)}D_X
 \leq \frac{1+\alpha+\epsilon}{2}D_X.
\end{align*}
第1行は \(Q_{F_1}\) と \(\sum_s X_s\) の Hilbert--Schmidt 内積への
Cauchy--Schwarz の不等式である。実際、
\(\|\sum_s X_s\|_{\mathrm{HS}}^2=|\bar G|\sum_s\beta(s)\) となる。
よってある \(s\) で、この平均以下の値を取る。
\(B_s=e_0\rho(s)P\colon S_0\to V\) と置き、
\(U=B_s^{-1}(F_1)\) とする。\(U\) 上では \(B_s=Q_{F_1}\rho(s)P\) であり、
\(\rho(s)\) はユニタリで \(e_0\) と可換なので
\[
 \dim U-\epsilon D_X
 \ \leq\ \|e_0P Q_U\|_{\mathrm{HS}}^2
 \ \leq\ \|Q_{F_1}\rho(s)P\|_{\mathrm{HS}}^2 .
\]
従って \(\dim U\le(1+\alpha+3\epsilon)D_X/2\) となる。\leanref{isSpread_of_second_moment}

商 \(V/F_1\) への写像の階数は \(D_X-\dim U\) である。
したがって \(\operatorname{codim}F_1>16\epsilon h\) の間は
\(\operatorname{im}B_s\) を加えるたびに、欠けている次元の少なくとも
\(2D_X/(5h)\) の割合が埋まる。実際、
\[
 \frac{1-\alpha-3\epsilon}{2}D_X
 \ge \frac25(1-\alpha)D_X
 \qquad (1-\alpha>16\epsilon).
\]
この割合を \(\lambda=2D_X/(5h)>0\) とする。
\(F_1=0\) での増加は高々 \(h\) なので \(\lambda\le1\) である。
欠ける割合が \(16\epsilon\) を超える限り上の選択を続ければ、
各段階で余次元が \(1-\lambda\le e^{-\lambda}\) 倍以下となる。
よって \(\lceil\lambda^{-1}\log(1/(16\epsilon))\rceil\) 段階以内に
余次元は \(16\epsilon h=2h/b^2\) 以下になる。使った集合 \(T\) は
\[
  |T|D_X\le\tfrac52\bigl(1+\log(1/(16\epsilon))\bigr)h
               \le 5h\log(1+2b^2)
\]
を満たす。\leanref{IsSpread.exists_few_le_finrank_iSup_range}

\subsection{分岐条件を課した始域}

\(N=\pi_N\overline{\pi_N}\) とし、\(w\mid\pi_N\) ごとに
\(z_w\) を底本が取る有限アデールのスカラー作用とする。
\(S\) は、どの \(w\) についても \(\sum_{j=0}^{N-2}z_w^j\) によって零になる部分空間である。
\(z_w^{N-1}=1\) なので、この平均作用素の階数は \(z_w\) の固定空間の次元に等しい。
\(z_w\) を加えて生成したレベルの指数は \(N-1\) 以上であり、両レベルで類の
自己同型群が自明である。そのため固定空間の次元は \(D_X/(N-1)\) 以下となる。
\(w\mid\pi_N\) は \(f\) 個なので
\[
 \operatorname{codim}_{S_0} S\le \frac{f}{N-1}D_X\le b^{-8}D_X.
\]
ここで使ったのは、中心作用の固定空間の次元をレベルの指数で評価する補題である。\leanref{finrank_ker_one_sub_scalingGenerator_le}

各写像を \(S\) に制限すると、\(S\) の補空間の像一つ分ずつ次元を失う。
正規化のスカラーはいずれも非零なので
\(\operatorname{im}(B_s|_S)=\mathfrak p_sS\) であり、失う次元は高々
\[
   |T|\operatorname{codim}_{S_0}S
   \le 5b^{-8}\log(1+2b^2)h\le 2h/b^2.
\]
先に得た \(2h/b^2\) と合わせた余次元が \(4h/b^2\) 以下であること、
および \(\dim S\le D_X\) が生成定理 B の二つの不等式を与える。\leanref{enough_forms}

\section{捻りによる行列式の評価}
\label{sec:determinant-estimates}

連鎖をひねると、無限素点の評価で使ったグラム行列式を
有限素点ごとに別の行列式と比較できる。必要な捻りは、
テータ持ち上げへの射影の分母と、比較に現れる Hecke \(L\) 値を
同時に制御する。以下では、その二つの条件から各段の行列式の比を求め、
連鎖全体に移す。

\subsection{捻りの選択と射影の分母}

\(r=\dim_E Y_0\)、\(t=5r\) とし、各段の比較に用いる空間を
\(W_j,V_j\) と書く。エルミート空間としての階数はそれぞれ \(r,r+t\) である。
\(W_j\) の主レベルの形式に \(e_0\) を作用させて得られる空間を \(F_j\) とする。
レベル外の分裂素点の Hecke 固有系ごとに、\(W_j\) と \(Y_0\) で
この射影後の重複度が一致する。特に \(\dim_{\mathbb C} F_j=h\) である。
\leansrc{New/Spanning/Multiplicities/EqualMultiplicities}{重複度の一致}

捻り \(\xi\) は、無限素点で重さ \(0\) の \(U(1)\) の保型指標で、
その分岐素点の集合を \(Q_\xi\) とする。
考える \(\xi\) はすべて \(1\le j\le m\) に対し、\(\xi^j\) が
\(Q_\xi\) の各素点で分岐する。したがって、どの \(\xi^j\) でも
レベルに現れる数 \(N_\xi=\prod_{w\in Q_\xi} \mathrm Nw\) は同じである。
\(N_\xi\) が \(\ell\) と互いに素なとき、\(\xi\) を \(\ell\) で許容する。

重さ \(-c(\tau)\) の指標 \(\nu\) と、\(1\le u\le t\) かつ
\(u\le c(\tau)\le4b\) を満たす整数 \(u\) の組を考える。
\(\nu\) の剰余体指標（tame part）が指定レベル素点の外で自明なとき、この組を臨界対と呼ぶ。
指定レベル素点のすべてで剰余体指標が非自明なら、
この対を分岐対と呼ぶ。以下で用いる比は
\[
 \mathcal L_{\nu,u}(\vartheta;\varepsilon)
 =
 \frac{|L(u,(\nu\vartheta^\varepsilon)^{-1})|}
      {|L(u,\nu^{-1})|},
 \qquad \varepsilon\in\{1,-1\}.
\]
ここで \(L\) は \(E\) 上の Hecke \(L\) 関数で、指標は
\(z\mapsto z/\bar z\) を通じて \(U(1)\) から引き戻す。
分子、分母ともに指定されたレベル素点とその共役の Euler 因子を除き、
分子ではさらに \(Q_\vartheta\) とその共役の因子を除く。
原典の主対とは、\(j\) に対し \(m<u\le r\) で、剰余体指標が
第 \(j\) 段の指標に一致し、無限素点の負の重さが
第 \(j\) 段の指標の負の重さ \(w_j(\tau)\) に対して
\(w_j(\tau)+r-t\le c(\tau)\le w_j(\tau)\) を満たす臨界対である。
主対に付随する階数 \(2u\) の空間を、階数
\(12,4,4,2u-20\) の四つのブロックに分ける。
\leansrc{New/Determinant/Twist/FitTwists}{臨界対と主対}

\(u_{j,\vartheta}\) を \(F_j\) 上のテータ持ち上げ、\(P_{j,\vartheta}\) を
その像 \(I_{j,\vartheta}\) への直交射影とする。
レベル外の分裂素点 \(w\) で、第 \(j\) 段の持ち上げに対応する値に
Hecke 多項式作用素を特殊化する。形式化でいう擬似持ち上げは、
この作用素の核と \(e_0\) の像と \(I_{j,\vartheta}^{\perp}\) の交わりである。
擬似持ち上げに直交するという条件は、すべてのこのような \(w\) に課す。
\leansrc{New/OneLift/TestedForms}{擬似持ち上げの定義}

\[
 B_{\mathrm{den}}=\frac{E_3}{mk}
\]
と置く。原典の正則な捻りの存在定理は、各 \(\iota\) に対して
\(\ell\) で許容される \(\xi_\iota\) を選び、次の条件を同時に満たす。
以下 \(|v|_\iota=|\iota(v)|_\ell\) と書く。
\begin{enumerate}
\item すべての分岐対と \(1\le j\le m\) および \(\varepsilon=\pm1\) に対して、
  \(\bigl(\mathcal L_{\nu,u}(\xi_\iota^j;\varepsilon)^{-1}\bigr)_\iota\)
  の高さは \(x_1\) 以下である。
\item 各 \(j\) について、素数 \(\ell\) で添字付けられた
  \(\beta_\ell\ge 0\) があり、任意の有限な素数集合 \(T\) で
  \(\sum_{\ell\in T}\beta_\ell\le B_{\mathrm{den}}\) である。
  \(\vartheta=\xi_\iota^j\) とすると、ある \(I_{j,\vartheta}\) の自己準同型
  \(\Lambda\) が存在し、
  \[
    |\det\Lambda|_\iota\ge\exp(-hr\beta_\ell).
  \]
  このとき \(\Lambda P_{j,\vartheta}\) は、\(\iota\) で整で、
  擬似持ち上げに直交する形式を \(\iota\) で整な形式に移す。
\item 同じ \(\beta_\ell\) で、各主対の階数 \(1\) からの持ち上げが生成する
  直線への射影も制御される。具体的には、その直線への射影を \(P\) とし、
  \(\ell\) 上のすべての場所で整で、階数 \(1\) の下降に対する
  擬似持ち上げに直交する試験形式 \(\Xi,\Xi'\) をとると、
  \[
       |\langle P\Xi,P\Xi'\rangle|_\iota\le\exp(\beta_\ell).
  \]
\end{enumerate}

この存在定理の入力は、\(W_m\) の \(e_0\) 成分に現れる有限個の
Hecke 固有系に付随する階数 \(r\) のガロア表現と、各主対に
付随する階数 \(1\) の表現である。付随表現のうち、捻り付き持ち上げの
跡から指標の列を取り除いた主成分に捻り定理を適用する。
この定理は一つの反円分指標の最初の \(m\) 個の冪について、可除係数の
Selmer 群の長さと、分岐対の \(L\) 値比の付値を同時に抑える。
\(C_{\mathrm{lev}}\) を原典のレベル群の位数の積、
\(C_1\) を指定主レベルでの \(U(1)\) の類の個数とすると、
分岐対の逆比に対する \(\log|\cdot|_\iota\) の上界は
\[
  \big((4f+3)r\,v_\ell(C_{\mathrm{lev}})+2^{-\ell}\big)\log\ell
\]
である。Ribet の方法による射影の分母の評価は、Selmer 群の長さと
局所・レベルの補正を、原典の \(c_\ell=\texttt{clearingBound}\) に
まとめる。階数 \(1\) の射影内積ではさらに \(v_\ell(C_1)\log\ell\)
を見込み、共通の \(\beta_\ell=(c_\ell+v_\ell(C_1))\log\ell\) を使う。
二つの素数和の上界が、それぞれ \(x_1\) と \(B_{\mathrm{den}}\) である。
ここでは、捻り定理と Selmer 群の長さから射影分母への評価を
次の原典の形で用いる。
\leansrc{New/Congruences/Selmer/TwistingTheorem}{Selmer 群に対する捻り定理}、
\leansrc{New/Congruences/FirstOrder/DenominatorBound}{射影分母の評価}、
\leansrc{New/Congruences/FitTwistsExist}{正則な族の存在}

\subsection{主対の Hecke 値とテータ微分}

固定した主対 \(p=(\nu,u)\) に付随する階数 \(1\) から階数 \(2u\) への
正規化した持ち上げを \(\phi_\vartheta\) と書く。
その \(L\) 値比は、指標の持ち上げに対するノルム公式を二つの
捻りで比較して
\begin{equation}
 \mathcal L_p(\vartheta;1)^{-1}
   =N_\vartheta^{2u-1}\frac{\|\phi_1\|^2}{\|\phi_\vartheta\|^2}
 \label{eq:main-pair-norm}
\end{equation}
となる。これに対し有限素点の評価を直接行うため、
\(\phi_\vartheta\) を含む直線への射影を、二つの試験形式に適用する。

先に \(\|\phi_1\|^2\) を定数族とみなした高さを評価する。
同じ階数で微分次数が \(0\) の臨界対に付随する持ち上げを
\(\phi^{(0)}_1\) とする。そのノルムに整数 \(n_L\) を掛けると
すべての有限素点で整になる。
\(\|\phi_1\|^2\) の高さの定義に用いる共通の因子を
\(n_L\|\phi^{(0)}_1\|^2/\|\phi_1\|^2\) とすれば、
有限素点からの損失を \(0\) にできる。無限素点では両ノルムの
対数に共通な定数が消える。主対の微分次数は
\(c(\tau)-u\ge\kappa(\tau)\) であり、微分次数 \(0\) の項との差は
\(\sum_\tau s_\infty(\kappa(\tau))\) 以上になる。
ノルム公式の誤差と \(\log n_L\) は \(B_{\mathrm{den}}\) 以内なので
\begin{equation}
 \operatorname{ht}\big((\|\phi_1\|^2)_\iota\big)
  \le B_{\mathrm{den}}-\sum_{\tau} s_\infty(\kappa(\tau)).
 \label{eq:untwisted-rankone-height}
\end{equation}
ここでノルムの評価はまず二次体 \(K\) を固定する共役について
行われる。原典の共役の絶対値に関する補題により、
すべての \(\mathbb C\) の自己同型について同じ評価が成立する。
\leansrc{New/OneLift/HeckeValues/AlgebraicPartSize}{無捻りノルムの高さ}

有限素点側で必要な微分の付値を
\[
 F(\ell)=\frac12\sum_{\tau\in\Phi}\sum_{e\ge 1}
     (1-\chi_K(\ell^e))
     \left(\frac{\kappa(\tau)}{\ell^e}-\kappa(-\ell^e\tau)\right)_+
\]
と書く。素数冪ごとに項をまとめると、任意の有限な素数集合 \(T\) に対して
\begin{equation}
 2\sum_{\ell\in T}F(\ell)\log\ell
   \le \sum_{\tau\in\Phi}s_{\mathrm{fin}}(\kappa(\tau),\tau).
 \label{eq:local-prime-total-lower}
\end{equation}
\leansrc{New/UpperBound/Size/Estimates}{局所素数和の合計}

原典の試験形式の構成では、主対による階数 \(2u\) の持ち上げを
上の四つのブロックに展開する。階数 \(12\) のブロックで、
付値の評価に達するテータ級数の係数を一つ選ぶ。この係数により、
付値が \(F(\ell)\) と一定因子の付値の和以下となる項を得る。
階数 \(4\) の第二ブロックには、分岐対 \(q_2\) の微分次数 \(0\) の
持ち上げを使い、その \(L\) 値比を \(\rho_\vartheta=\mathcal L_{q_2}(\vartheta;1)\)
とする。残る二つのブロックには \(U(1)\) の類群の指標を二つに
分けて割り当てる。類の代表元を平行移動して自明な類の代表を単位元に取る。
指標で重みを付けた正規化類和をすべての指標について足すと、
指標の直交性により選んだ係数に戻る。非アルキメデス性により、
少なくとも一つの類和の \(\iota\)-絶対値がこの係数の絶対値以上になる。
さらに指標の分解を選び、残る二つのブロックの代数的 \(L\) 値の積の
\(\iota\)-絶対値が、その全選択にわたる幾何平均以上になるようにする。各ブロックの積で試験形式を作り、
持ち上げとの内積を類和に展開すると、以上の量が積として現れる。

この局所構成の使用形は次の通りである。捻りと有限素点に依存しない
\(c_0\ne0\) があって、すべての \(\sigma\) で
\(|\sigma(c_0)|\le\exp(fmb/k+3x^{34})\le\exp(B_{\mathrm{den}})\) である。
任意の \(\iota\) とその素数で許容される \(\vartheta\) に対し、
ある試験形式 \(\Xi\) について
\begin{equation}
 |\langle\Xi,\phi_\vartheta\rangle|_\iota
    \ge |c_0\rho_\vartheta|_\iota\,\ell^{-F(\ell)} .
 \label{eq:test-pairing-lower}
\end{equation}
ここで \(c_0\) はテータ係数に掛ける一定因子、第二ブロックの
無捻りの代数的 \(L\) 値、残る二ブロックの幾何平均の積である。
試験形式が \(\ell\) 上のすべての場所で整であることと
擬似持ち上げへの直交性も、同じ構成で得られる。
\leansrc{New/Derivatives/Sharpness/BoundAttained}{テータ係数の付値}、
\leansrc{New/OneLift/HeckeValues/TwistedValueHeights}{試験形式の構成}

\(\vartheta=\xi_\iota^j\) とする。\(\iota\) と
\(\iota\circ\mathrm{conj}\) でそれぞれ \eqref{eq:test-pairing-lower} の
試験形式 \(\Xi,\Xi'\) を選ぶ。 \(\rho_\vartheta\) は実数なので、
\(A_\iota=\langle\Xi,\phi_\vartheta\rangle\langle\phi_\vartheta,\Xi'\rangle\)
と置けば
\[
 \left| c_0\bar c_0\,\frac{\rho_\vartheta^2}{A_\iota}\right|_\iota
      \le \ell^{2 F(\ell)}.
\]
両内積は \eqref{eq:test-pairing-lower} により零でない。
\(P\) を \(\mathbb C\phi_\vartheta\) への直交射影とすると
\(\langle P\Xi,P\Xi'\rangle=A_\iota/\|\phi_\vartheta\|^2\) である。
したがって \eqref{eq:main-pair-norm} の右辺は \(N_\vartheta^{2u-1}\) と
次の四因子の積になる。
\[
 \|\phi_1\|^2,\qquad
 \langle P\Xi,P\Xi'\rangle,\qquad
 \rho_\vartheta^{-2},\qquad
 \frac{\rho_\vartheta^2}{A_\iota}.
\]
順に \eqref{eq:untwisted-rankone-height}、
射影内積の評価、分岐対の逆比の評価、
および共通因子 \(c_0\bar c_0\) と \eqref{eq:local-prime-total-lower} により、
高さの上界は
\[
 \begin{gathered}
 B_{\mathrm{den}}-\sum_\tau s_\infty(\kappa(\tau)),\qquad
 B_{\mathrm{den}},\qquad 2x_1,\\
 2B_{\mathrm{den}}+\sum_\tau s_{\mathrm{fin}}(\kappa(\tau),\tau)
 \end{gathered}
\]
となる。さらに \(N_\vartheta^{2u-1}\) は \(\iota\) で単元である。
\(x_1\le B_{\mathrm{den}}\) を使って足し合わせると
\begin{equation}
 \operatorname{ht}\bigl(
   (\mathcal L_p(\xi_\iota^j;1)^{-1})_\iota\bigr)
 \le 6 B_{\mathrm{den}}-A_0.
 \label{eq:main-pair-height-lower}
\end{equation}
\leansrc{New/OneLift/HeckeValues/TwistedValueHeights}{主対の逆比の高さ}

\subsection{一段の行列式の比較}

\(Z_j(\vartheta)=\det(u_{j,\vartheta}^*u_{j,\vartheta})\) とする。
比較作用素 \(U_{j,\vartheta}\) は、補助空間上の形式の組を
\(V_j\) 上の形式へ貼り合わせる写像 \(G_{j,\vartheta}\) とその随伴を
合成したものに、レベルの整数とフック数の二乗を掛けて得られる。
三回のテータ積分の交換をこの写像に適用すると、
\(F_j\) 上に \(\vartheta\) によらない正定値作用素 \(U_j^\circ\) があり、
\begin{equation}
 u_{j,\vartheta}^*U_{j,\vartheta}u_{j,\vartheta}
      =\rho_j(\vartheta)U_j^\circ,\qquad
 \rho_j(\vartheta)=
 \frac{\mathcal H_j(\vartheta)}{\mathcal H_j(1)}
 \label{eq:comparison-operator-lower}
\end{equation}
となる。\(\mathcal H_j\) は、交換に現れる三つの Hecke 値の積
（原典の \(\texttt{heckeValueProd}\)）である。
\(U_j^\circ\) は無捻りで得られる周期写像のグラム作用素の正の定数倍である。
その周期写像は単射なので、\(\det U_j^\circ\) は零でない。

\(\rho_j(\vartheta)\ne0\) である。この構成の局所的な性質は二つある。まず \(U_{j,\vartheta}\) は
テータ持ち上げを擬似持ち上げに直交する形式に移す。次に、
\(N_\vartheta\) と互いに素な素数では整性を保つ。
整な試験ベクトルのタイトフレームで \(G_{j,\vartheta}G_{j,\vartheta}^*\)
を階数一の作用素の和に展開し、各内積の分母をレベルの整数で払うことで
後者を得る。同じフレームに対する Cauchy--Schwarz の評価から、
無捻りの作用素は、各共役でノルムを
\(\exp(3fmb(\log b)^2)\) 倍以下にする。
\leansrc{New/OneLift/ComparisonOperator/ThreeExchanges}{三回の交換}、
\leansrc{New/OneLift/LiftVolumeHeight}{比較作用素の構成と評価}

\(B_{j,\vartheta}\) を \(P_{j,\vartheta} U_{j,\vartheta}\) の
\(I_{j,\vartheta}\) への制限とし、\(D_j(\vartheta)=\det B_{j,\vartheta}\)
と置く。\(u_{j,\vartheta}\) はその像への同型であり、
\(u_{j,\vartheta}^*\) は \(I_{j,\vartheta}^\perp\) 上で零である。
したがって
\[
 u_{j,\vartheta}^* U_{j,\vartheta} u_{j,\vartheta}
  =(u_{j,\vartheta}^*u_{j,\vartheta})
    (u_{j,\vartheta}^{-1} B_{j,\vartheta}u_{j,\vartheta}).
\]
行列式を取り、\eqref{eq:comparison-operator-lower} を使うと
\[
 Z_j(\vartheta)D_j(\vartheta)
      =\det(U_j^\circ)\rho_j(\vartheta)^h,\qquad
 \frac{Z_j(1)}{Z_j(\vartheta)}
   =\frac{D_j(\vartheta)}{D_j(1)}\,\rho_j(\vartheta)^{-h}.
\]

有限素点では、\(\vartheta=\xi_\iota^j\) に対し
\(\Lambda B_{j,\vartheta}\) が \(I_{j,\vartheta}\) の整な形式を保つ。
各固有ベクトルを、整で、さらに絶対値が \(1\) より大きい数で
掛けると整性を失うように規格化すれば、固有値は
\(\iota\)-絶対値が \(1\) 以下である。したがって
\[
 |D_j(\xi_\iota^j)|_\iota
   \le|\det\Lambda|_\iota^{-1}
   \le\exp(hr\beta_\ell).
\]
無限素点では、\(I_{j,1}\) に \(U_{j,1}\) を作用させた像は
擬似持ち上げに直交するので、そこに射影を適用しても、各共役で
ノルムは増えない。この射影の評価には、
Hecke 作用素の特殊化の共役に対する随伴性を使う。
\(C_{\mathrm{arch}}=fmb(\log b)^2\) として固有値を掛け合わせると
\(|\sigma(D_j(1))|\le\exp(3h C_{\mathrm{arch}})\) を得る。
\(D_j(1)\) を高さの共通の因子として使うと
\[
 \operatorname{ht}\bigl((D_j(\xi_\iota^j)/D_j(1))_\iota\bigr)
     \le h(3C_{\mathrm{arch}}+r B_{\mathrm{den}}).
\]
\leansrc{New/OneLift/LiftVolumeHeight}{圧縮の行列式の有限素点・無限素点評価}

三つの交換で現れる \(\rho_j(\vartheta)\) は \(N_\vartheta\) の整数冪を除くと、
\(k/2\) 個の主対の比と \(5m\) 個の分岐対の比の積の二乗である。
主対の点は \(7m/2<u\le 7m/2+k/2\) であり、
残る点の個数は \(7m/2+m+m/2=5m\) である。主対と最初の \(7m/2\) 個では
\(\varepsilon=1\)、平坦対のうち \(m\) 個で \(\varepsilon=1\)、\(m/2\) 個で \(\varepsilon=-1\) を使う。
\eqref{eq:main-pair-height-lower} と分岐対の評価を足すと
\[
 \operatorname{ht}\bigl((\rho_j(\xi_\iota^j)^{-1})_\iota\bigr)
  \le k(6B_{\mathrm{den}}-A_0)+10mx_1 .
\]
許容データの数値条件
\(3C_{\mathrm{arch}}+(r+6k)B_{\mathrm{den}}+10mx_1
        \le k B_{\mathrm{den}}\log b\)
により、各段について
\begin{equation}
 \operatorname{ht}\bigl((Z_j(1)/Z_j(\xi_\iota^j))_\iota\bigr)
   \le h\big(E_4/m-k A_0\big)
 \label{eq:single-relative-height}
\end{equation}
が従う。
\leansrc{New/OneLift/HeckeValues/TwistedValueHeights}{Hecke 値の因数分解}、
\leansrc{New/OneLift/LiftVolumeHeight}{一段の高さの評価}

\subsection{連鎖の行列式}

連鎖のノルムを計算するため、最初の空間に高さ \(t\) のブロックを
付け足していく持ち上げを三角形に配置する。節点 \((i,n)\) は、
共通の部分と第 \(i+1\) から第 \(i+n\) のブロックを持つ。
右からブロックを加える写像に対し、左からブロックを除く
転置持ち上げを比較する。両者を入れ替える二つの経路の比較が
交換恒等式を使う箇所である。

有限環上での対応は具体的には次のようになる。列と行の添字を
二組に分けたとき、可逆行列 \(g\) の先頭の列と \(g^{-1}\) の
末尾の行からなる二種類の行列に、重み \(\Psi\) と \(\Phi\) を置く。
それぞれの台をそのような行列に限り、
\(\Phi(\operatorname{invLastRows}g)=\chi(\det g)\Psi(\operatorname{firstCols}g)\) とする。
二つの経路で、底本の指定した重さ・分岐・符号・階数を満たす核を \(K,K'\) とする。
交換により
\(\sum\Psi K'=(\kappa_0/c)\sum\Phi K\) を得る。
\(|c|^2\) は剰余体の位数の冪と部分 \(L\) 値の絶対値の二乗の積で決まり、
\(\kappa_0\) は捻りと独立である。
この等式を点データと有限 Fourier 変換を施したデータに
適用することで、三角形に必要な二つの交換則を得る。
\leansrc{New/Determinant/Exchange/TheoremForPartners}{テータ積分の交換}、
\leansrc{New/Determinant/Exchange/OneSquare}{一つの正方形での交換}、
\leansrc{New/Determinant/Chain/ChainNorms}{三角形の交換則と単射性}

\(c_Q(\xi)\) を
\[
 N_\xi^{rt}\prod_{w\in Q_\xi}\prod_{t<s\le t+r}(1-(\mathrm Nw)^{-s})
\]
と書く。\(R_j(\xi^j)\) を \(V\) 上の第 \(j\) 段の Rallis の公式の
Hecke 作用素とする。単独の持ち上げに Rallis の公式を適用し、
\(W_j\) と \(Y_0\) の Hecke 固有系の重複度をそろえると、
捻りに依存しない \(A_j\ne0\) に対して
\[
 Z_j(\xi^j)=(A_jc_Q(\xi))^h\det R_j(\xi^j)
\]
を得る。連鎖については三角形の交換則でブロックを一つずつ移す。
具体的には、行 \(i\) の第 \(n\) 段までの合成を \(X_{i,n}\) と書き、
各矢印にその始域の有限レベルでの一般線型群の位数 \(w_n\) を掛ける。
\(a_{i,n}\) を右へ進む矢印、\(d_{i,n}\) を次の行へ下る矢印とする。
\(j_i\) は主レベルの空間を第 \(i\) 行の始点に等長に埋め込む写像である。
\(x'\) を \(j_{i+1}x'=d_{i,0}a_{i,0}j_i x\) により定める。
正方形ごとの交換則を合成し、随伴に関する交換則を使うと、次の漸化式を得る。
\begin{multline*}
  \langle X_{i,n+1}j_i x,X_{i,n+1}j_i y\rangle\\
   = \lambda_{i,n}\,\frac{w_0w_1w_n}{w_{n+1}}\,
     \langle X_{i+1,n}j_{i+1}x',X_{i+1,n}j_{i+1}y\rangle.
\end{multline*}
\(\lambda_{i,n}\) の絶対値は、捻りに依存しない係数と \(N_\xi\) の半整数冪を除き、
\(\prod_{i+1<j\le i+n+1}J_{i+1,j}(\xi)\) である。
右辺の \(x'\) との内積を、行 \(i\) の最初の矢印のグラム作用素による
\(x\) との内積に戻す。その際に \(w_0\) が一つ消える。
空の合成から帰納すると、行 \(i\) のグラム作用素は、
各行の最初のグラム作用素の積と \(w_0w_1^n/w_n\) に比例し、
係数に現れる対は \(i<r'<j\le m\) である。
\leansrc{New/Determinant/Chain/TriangleNorms}{重み付き連鎖の帰納公式}

行 \(0\) からの連鎖では、各行の最初のグラム作用素に Rallis の公式を適用する。
最終レベルの指数をブロック上三角行列の数で書き直すと、
\(w_0w_1^m/w_m\) が相殺される。この規格化の下で、
原典のグラム作用素の公式は
\[
 \Theta_\xi^*\Theta_\xi
       = c(\xi)c_Q(\xi)^m \prod_{j=1}^m R_j(\xi^j)
\]
であり、
\[
 |c(\xi)| = C\,N_\xi^{e/2}\prod_{1\le i<j\le m} J_{ij}(\xi),
 \qquad C>0,\quad e\in\mathbb Z.
\]
ここで、第 \(j\) 段の無捻りの指標を \(\eta_j=\mu_j^0\) と書くと
\[
 J_{ij}(\xi)=
   \prod_{u=1}^{t}\mathcal L_{\eta_j/\eta_i,u}(\xi^{j-i};1).
\]
各対 \(i<j\) はブロックを移す際に現れ、\(\eta_j/\eta_i\) の
指定レベルでの剰余体指標は固定指標の \(j-i\) 乗である。
\(1\le j-i\le m\) のため、この積に現れる対は分岐対である。

連鎖の単射性には、三角形の対応する重み付き連鎖の単射性を使う。
正方形の交換則から、長さ \(n+1\) の重み付き連鎖で
\(j_i v\) が零に送られるなら、次の行の長さ \(n\) の重み付き連鎖で
\(d_{i,0}a_{i,0}j_i v\) が零に送られる。
後者は \(j_{i+1}\) の像にあるので、長さに関する帰納と
\(d_{i,0}a_{i,0}\) の \(j_i\) の像上での単射性により \(v=0\) となる。
規格化の違いによるレベル指数と Weyl 加群の次元は正なので、
この単射性は \(\Theta_\xi\) にも移る。同様に単独の持ち上げも単射であり、
すべてのグラム行列式は正である。行列式の等式の両辺で絶対値を取り、
正のグラム行列式はそのままにすると、
\(c_Q(\xi)\) と Hecke 作用素の行列式を消去できる。
\begin{equation}
 \frac{Z(1)}{Z(\xi)}
 =
 \left(N_\xi^{e/2}\prod_{1\le i<j\le m}J_{ij}(\xi)\right)^{-h}
       \prod_{j=1}^m\frac{Z_j(1)}{Z_j(\xi^j)} .
 \label{eq:chain-relative-factorization}
\end{equation}
\(N_\xi^{e/2}\) は \(\ell\nmid N_\xi\) のとき \(\iota\) で絶対値 \(1\) である。
左の括弧の逆数の高さは、含まれる分岐対の個数から \(m^2t x_1\) 以下で、
これは \(b\) 以下である。
\eqref{eq:single-relative-height} を \(m\) 段について合計すれば
\begin{equation}
 \operatorname{ht}\bigl((Z(1)/Z(\xi_\iota))_\iota\bigr)
       \le h(b+E_4-H_0).
 \label{eq:chain-relative-height}
\end{equation}
\leansrc{New/Determinant/Chain/NormOfChain}{連鎖のグラム作用素と行列式の分解}

\subsection{基底を補う形式}

基底の補充には、射影 \(e_0\) を作用させる前の主レベルの空間を
深いレベルの形式から生成する。原典では二つの格子状の図式を作り、
その外縁に沿う最初の経路の周期を、
\(m^2\) 個の正方形で一回ずつ経路を曲げて、最後の経路の周期に移す。
各曲がりでは有限行列データの和に交換恒等式を適用する。
対角に沿う曲がりでは階数 \(1\) の核との関係を取り出してから同じ比較を行う。
この操作で生じる非零のスカラー \(c_{\mathrm{path}}(\xi)\) は、
\begin{equation}
 \left|\frac{c_{\mathrm{path}}(\xi)}{c_{\mathrm{path}}(1)}\right|
   =N_\xi^{e'/2}\prod_{0\le i,j<m} T_{ij}(\xi),\qquad
 T_{ij}(\xi)=
    \prod_{u=1}^{t}\mathcal L_{\eta_{i+1}\eta_{j+1},u}(\xi^{j+1};1)
 \label{eq:turn-ratio-lower}
\end{equation}
を満たす。剰余体指標の指数 \(i+j+2\) は \(2\) 以上 \(2m\) 以下であり、
指定レベルで非自明である。したがって各 \(T_{ij}\) に現れる対も
分岐対である。
\leansrc{New/Determinant/AuxiliaryForms/Turns}{各曲がりの \(L\) 値比}、
\leansrc{New/Determinant/AuxiliaryForms/DiagonalTurns}{全経路の比較}

最後の経路の周期は、有限行列の組 \(S\) について、
二つの半経路の持ち上げの内積を係数 \(\Psi(S)\) で足した和である。
規格化を固定すると、\(\ell\nmid N_\xi\) のとき各辺の核の分母と
類和の分母は、\(\xi\) に依存しない整数を掛けて払える。
無捻りでは、核の Weyl ノルムの評価と有限和の項数から
周期の各共役の上界が得られる。
元の持ち上げの値を復元するには、双対のデルタ形式と対にする。
双対のペアリング行列 \(D\) には、整数行列 \(C\) と正のフック数 \(H\) が
あって \(CD=HI\) となる。すべての整数ベクトルとのペアリングが
\(H\) の整倍なら、\(C\) の各行でテストし \(H\) を約して、元の値の
各座標が整であることがわかる。無限素点でも、ペアリングの上界を
双対ノルムで割ることで Weyl ノルムの上界に移せる。

以上により、次の定量的な形で主レベルの空間が生成される。
深いレベルの空間から主レベルへの全射 \(s\) と定数 \(c_{\mathrm{sc}}\ne 0\)
があり、深いレベルの空間はすべての有限素点で整で、
各共役の値の二乗 Weyl ノルムが \(\exp(x^{99}b^2)\) 以下の形式 \(\psi\)
で張られる。その像 \(c_{\mathrm{sc}}s(\psi)\) の持ち上げは
\(c_{\mathrm{path}}(\xi)/c_{\mathrm{path}}(1)\) を掛ければ許容素点で整である。
無捻りでの像の値の二乗 Weyl ノルムは、
深いレベルでの上界に \(\exp(2x^{201}b^2)\) を掛けたもの以下で、
したがって \(\exp(x^{202}b^2)\) 以下である。
これらの数値を伴う局所・無限素点での核の評価を
\leansrc{New/Determinant/AuxiliaryForms/SpanningForms}{最後の経路の周期と補助形式の構成}
から取る。

この生成系をレベル群の位数倍の射影で \(V\) に送る。
射影後も整性を保ち、無捻りの二乗ノルムに加わる対数は
\(x^{99}\) 以下である。\(x^{99}+x^{202}b^2\le 2x_1^2b^2\) より、
得られる補助形式 \(z\) は
\[
 \|\Theta_1 z\|_\sigma^2\le \exp(2x_1^2b^2)
 \qquad(\sigma\in\operatorname{Aut}(\mathbb C))
\]
を満たす。また \(\gamma(\xi)=c_{\mathrm{path}}(\xi)/c_{\mathrm{path}}(1)\) と置けば、
\(\gamma(\xi)\Theta_\xi z\) はすべての許容素点で整である。
\eqref{eq:turn-ratio-lower} と分岐対の評価により
\begin{equation}
 \operatorname{ht}\bigl((|\gamma(\xi_\iota)|^{-2})_\iota\bigr)
 \le 2m^2t x_1\le b .
 \label{eq:auxiliary-factor-height}
\end{equation}
\leansrc{New/Determinant/Twist/LongForms}{補助形式を目標空間に射影した後の評価}

\begin{proposition}\label{prop:determinant-inputs}
すべての許容データについて、\(\Theta_\xi\) は各捻り \(\xi\) に対し
単射である。また各 \(\iota\) に対しその素数で許容される捻り
\(\xi_\iota\) と、捻りごとの定数 \(\gamma(\xi)\ne0\) が存在し、
\(\gamma(1)=1\) かつ
\[
 \operatorname{ht}\bigl((Z(1)/Z(\xi_\iota))_\iota\bigr)
       \le h(b+E_4-H_0),\qquad
 \operatorname{ht}\bigl((|\gamma(\xi_\iota)|^{-2})_\iota\bigr)\le b.
\]
\(V\) は、すべての許容素点で \(\gamma(\xi)\Theta_\xi z\) が整で、
すべての \(\sigma\) に対して
\(\|\Theta_1 z\|_\sigma^2\le\exp(2x_1^2b^2)\) である補助形式 \(z\)
で生成される。
\end{proposition}
\begin{proof}
単射性は三角形の持ち上げの比較による。正則な族を選ぶと
一段ごとの評価 \eqref{eq:single-relative-height} が得られ、
連鎖の分解 \eqref{eq:chain-relative-factorization} に代入して
\eqref{eq:chain-relative-height} を得る。補助形式は深いレベルからの
全射と \(e_0\) による射影で生成され、その同じ \(\gamma(\xi)\) について
\eqref{eq:auxiliary-factor-height} が成り立つ。
\end{proof}

\section{比較不等式と主定理の証明}\label{sec:conclusion}

\begin{proposition}[対数二乗余体積と指数の下界]\label{prop:lower-bound}
\(y=(y_1,\ldots,y_n)\) は \(V\) の一次独立な族で
\(n\ge(1-4/b^2)h\) とする。
その対数二乗余体積が \(c\) 以下、対数指数が \(\delta\) 以下ならば
\begin{equation}\label{eq:lower-bound}
  h(H_0-3E_4)\le c+2\delta .
\end{equation}
\end{proposition}

\begin{proof}
命題 \ref{prop:determinant-inputs} の補助形式から
\(n'=h-n\) 個を選び、\(y\) を基底に延長する。
同じ命題の \(\xi_\iota,\gamma(\xi)\) を用いて、命題 \ref{prop:gram-product}
に \(B_{\mathrm{Ar}}=2x_1^2b^2\)、\(B_{\mathrm{fin}}=b\)、
\(B_Z=h(b+E_4-H_0)\) を代入する。すると
\[
 0\le c+2\delta+n'(2x_1^2b^2+b)+h(b+E_4-H_0).
\]
補う個数は \(n'\le4h/b^2\) であり、
許容データの \(4(2x_1^2+b^{-1})\le b\) を使うと
\(n'(2x_1^2b^2+b)\le hb\) となる。
\(b\le E_4\) を用いて \(hb\) の二つの寄与を吸収すれば
\eqref{eq:lower-bound} を得る。
\leanlinks{\leanref{integral_forms_long}}
\end{proof}

\begin{theorem}[比較不等式]\label{thm:main-inequality}
任意の許容データについて
\begin{equation}\label{eq:main-ineq}
 m\mathcal R\le fa(\log b)^5 .
\end{equation}
\end{theorem}

\begin{proof}
定理 \ref{thm:spanning} の平行移動の集合 \(T\) を取る。
定理 \ref{thm:upper-bound} (1) の短い整形式が \(S\) を張るので、
それらの像 \(\mathfrak p_g\phi\), \(g\in T\) から、
\(\sum_{g\in T}\mathfrak p_g S\) の基底 \(y_1,\ldots,y_n\) を選べる。
すると \((1-4/b^2)h\le n\le h\)、かつ \(h>0\) である。
各像に対する二乗ノルムの上界と Hadamard の不等式より、
対数二乗余体積の上界として
\[
   c=n(H_0-m\mathcal R+E_4)
\]
が使える。また定理 \ref{thm:upper-bound} (2) より対数指数の上界として
\[
  \delta=2|T|(\dim S)E_3\le 10h\log(1+2b^2)E_3
\]
が使える。ここで左辺を使った指数評価を右辺へ大きくしてもよい。
命題 \ref{prop:lower-bound} から
\[
 h(H_0-3E_4)
  \le n(H_0-m\mathcal R+E_4)+20h\log(1+2b^2)E_3.
\]
\(H_0+E_4\ge0\) なので \(n(H_0+E_4)\le h(H_0+E_4)\) である。
右辺にこの評価を入れると \(hH_0\) が消え、
\[
 n\,m\mathcal R
   \le h\bigl(4E_4+20\log(1+2b^2)E_3\bigr)
\]
が残る。許容データでは \(4/b^2\le1/2\) であり \(n\ge h/2>0\) である。
右辺の括弧の二倍は \(fa(\log b)^5\) 以下なので、\(n\) で割ればよい。
\leanlinks{\leanref{star_inequality}}
\end{proof}

\begin{corollary}[素数和の片側評価]\label{cor:prime-bound}
\(q_0\) を \(18480\) の正の倍数、\(d\equiv q_0/2-1\pmod{q_0}\) を素数とする。
任意の \(\varepsilon>0\) について、十分大きな実数 \(y\) に対し
\begin{equation}\label{eq:prime-bound}
 R_{q_0}(y)\le y^{-1/2+\varepsilon}
\end{equation}
が成り立つ。
\end{corollary}

\begin{proof}
\eqref{eq:main-ineq} を \eqref{eq:total-remainder} に代入し、
\(ma|\Phi_b|\) で割る。\(f=\varphi(q_0)|\Phi_a|\) および
\(|\Phi_a|/|\Phi_b|\le\log b\) を使うと、
\begin{equation}\label{eq:remainder-intermediate}
 R_{q_0}(x)
 \le\frac{\varphi(q_0)(\log b)^6}{m}
       +|R_{q_0}(x_1)|+\log b\,|R_1(x_1)|.
\end{equation}
\(0<\eta<1/100\)、\(7\eta\le\varepsilon\) とし、
命題 \ref{prop:admissible-data} を \(\eta\) で適用する。
十分大きな整数 \(x\) では
\(\log b\le2x^\eta\)、\(m\ge\sqrt{x}\)、\(\log x_1\ge x^{\eta/12}\) である。
補題 \ref{lem:fixed-modulus-pnt} に \(B\ge24/\eta\) を取れば、両方の
\(|R_{q_0}(x_1)|,|R_1(x_1)|\) は \(x^{-2}\) 以下となる。
したがって \eqref{eq:remainder-intermediate} の右辺は
\[
 64\varphi(q_0)x^{-1/2+6\eta}+x^{-2}+2x^{-2+\eta}
 \le (64\varphi(q_0)+3)x^{-1/2+6\eta}
\]
である。
\(x=\lfloor y\rfloor\) とすれば、補題 \ref{lem:fixed-modulus-pnt} の差の評価から
\(R_{q_0}(y)-R_{q_0}(x)\le26/y\) である。
\(\theta=-1/2+6\eta\) は \([-1,0]\) にあり、\(x\ge y/2\) なので
\(x^\theta\le2y^\theta\) である。
従って
\[
 R_{q_0}(y)\le (2(64\varphi(q_0)+3)+26)y^\theta .
\]
定数を十分大きな \(y^\eta\) に吸収し、\(7\eta\le\varepsilon\) を使えば
\eqref{eq:prime-bound} を得る。
\leanlinks{\leanref{eventually_primeRemainder_le_rpow}}
\end{proof}

\begin{proof}[定理 \ref{thm:grh} の証明]
\textup{(1)}
\(\chi\) を任意の法 \(q\) の指標とし、\(\RePart\rho>1/2\) とする。
\(\rho=1\) では主指標が除外され、非主指標については Dirichlet の
非消滅定理 \(L(1,\chi)\ne0\) による。以下では \(\rho\ne1\) とする。
\(q_0=18480q\)、\(d_0=q_0/2-1\) とおく。
\(d_0\equiv7\pmod8\) であり、\(q_0=2(d_0+1)\) かつ \(d_0\) は奇数なので
\((q_0,d_0)=1\) である。
\(1/2<\beta<\RePart\rho\) を選び、系 \ref{cor:prime-bound} を
\(\varepsilon=\beta-1/2\) で用いると、法 \(q_0\) の剰余類 \(d_0\) に属する
すべての素数 \(d\) について命題 \ref{prop:one-sided-criterion} の仮定が成立する。
\(\chi\) から誘導した法 \(q_0\) の指標を \(\chi_0\) とすれば
\(L(\rho,\chi_0)\ne0\) である。
誘導の公式
\[
 L(s,\chi_0)=L(s,\chi)
  \prod_{\substack{p\mid q_0\\p\nmid q}}\bigl(1-\chi(p)p^{-s}\bigr)
 \qquad(s\ne1)
\]
より \(L(\rho,\chi)\ne0\) を得る。

\textup{(2)}
\(0<\RePart s<1\) で \(L(s,\chi)=0\) とし、\(\chi\) を導く原始指標を \(\psi\) とする。
誘導に伴う有限個の Euler 因子はこの帯で非零なので \(L(s,\psi)=0\) である。
ガンマ因子はこの帯で正則かつ非零であるから、完備 \(L\) 関数も \(s\) で零となる。
\(\psi\) の関数等式により \(\psi^{-1}\) の完備 \(L\) 関数は \(1-s\) で零となり、
ガンマ因子を除いて \(L(1-s,\psi^{-1})=0\) を得る。
\(\psi^{-1}\) は \(\chi^{-1}\) を導くので \(L(1-s,\chi^{-1})=0\) となる。
\leanlinks{\leansrc{Lib/Analytic/ZeroSymmetry}{非原始指標を含む零点の反射}}
(1) を二つの零点にそれぞれ適用すると
\(\RePart s\le1/2\) および \(1-\RePart s\le1/2\) が従う。

\textup{(3)}
\(\zeta(s)=0\) とし、自明零点を除く。\(s=0\) は
\(\zeta(0)=-1/2\) により除かれ、\(s=1\) は極である。
\(\pi^{-s/2}\Gamma(s/2)\) はこの \(s\) で有限かつ非零なので、
完備ゼータ関数も \(s\) で零となる。
完備ゼータ関数は \(s\mapsto1-s\) で不変であり、
\(1-s\ne0,1\) でその値にガンマ因子の逆数を掛けると
\(\zeta(1-s)=0\) を得る。
法 \(1\) における (1) を二つの零点へ適用すると、
\(\RePart s=1/2\) が従う。
\leanlinks{\leanref{grh}}
\end{proof}


\begin{thebibliography}{9}
\bibitem{alpoege}
Levent Alpöge,
\emph{一般化リーマン予想の形式証明についての公表文},
2026年。
\href{https://alpo.ge/remark.pdf}{PDF}。

\bibitem{grh-source}
Levent Alpöge with Claude, \emph{GRH for Dirichlet L-functions, in Lean 4},
Luke Sernau らによる可読版。
\href{\upstream/grh/README.md}{底本の README}、
\href{\upstream/grh/GUIDE.md}{読解案内}、
\href{\upstream/grh/REFERENCES.md}{底本の文献一覧}。
コミット \texttt{72fead102f66}。

\bibitem{mv}
Hugh L. Montgomery and Robert C. Vaughan,
\emph{Multiplicative Number Theory I. Classical Theory},
Cambridge Studies in Advanced Mathematics 97,
Cambridge University Press, 2007。
\href{https://doi.org/10.1017/CBO9780511618314}{DOI}。
\end{thebibliography}
\end{document}
